\section{Eisenstein congruences for U(3) and GSp(4)}\label{sec:pf:eis}

\providecommand{\eisQZ}{\Q_\ell/\Z_\ell}
\providecommand{\eiscyc}{\chi_{\mathrm{cyc}}}
\providecommand{\eisVt}{\widetilde V}
\providecommand{\eisX}{\mathcal X}
\providecommand{\eisGG}{\mathcal G_3}
\providecommand{\eisv}{v_\ell}
\providecommand{\eisrz}{r_0}
\providecommand{\eiseis}{\mathrm{eis}}
\providecommand{\eistr}{\delta}
\providecommand{\eiscot}{\lambda^{\mathrm{cot}}}
\providecommand{\eisKill}{\mathcal N}

This section proves four results of the formal development \cite{LeanSource2026} which supply the automorphic
inputs of the plus part of the window bound (\cref{sec:window:plus,sec:window:eis}). Two of them are counting
statements in the sense of \cref{pf:vwin-he:countdef}:
\begin{itemize}
\item $\sfname{CountI}$, the bound for the lengths $V_\ell(\theta,j)$ of \cref{pf:vwin-he:V} in case I, summed over
  primes $\ell$ (\cref{thm:window:countI}, restated and proved as \cref{pf:eis:main});
\item $\sfname{CountV}$, the bound in case V (\cref{pf:eis:countV}).
\end{itemize}
The other two are automorphy statements:
\begin{itemize}
\item the symplectic automorphy statement $(\mathrm{Sp})$ of \cref{pf:vwin-he:autstmts}, which is the input of the
  counts in cases II and III (\cref{pf:eis:sp});
\item a density statement for automorphic points of unitary groups in three variables
  (\cref{pf:eis:dense}), which enters the proofs of the unitary automorphy statements $(\mathrm{U}_{\mathrm{ord}})$
  and $(\mathrm{U}_i)$ of \cref{pf:vwin-he:autstmts}.
\end{itemize}
All four are proved from the package $\sfname{HypothesesH'}$ of \cref{pf:vwin-he:H}; the counts use in addition
Lemma K (\cref{thm:window:lemmaK}), the vanishing theorem $\sfname{KatoTwoW}$ (\cref{thm:kato}) and the unitary
automorphy statements.

\emph{The counts.} Both counts are proved by the quantitative form of Ribet's method \cite{Ribet1976}, in the
setting of \cite{BellaicheChenevier2009}, described in steps (i)--(v) of \cref{sec:window:eis}. Fix a character
$\theta$, a twist $j$ and an imaginary quadratic field $K$ with quadratic character $\eta$, and put
$\Psi=\eiscyc^j\theta$. The Eisenstein representation $\eisrz=\Psi\oplus1\oplus\Psi^{-1}$ of $G_K$ extends to a
homomorphism of $G_\Q$ into the group $\eisGG=\GL_3\rtimes\{\pm1\}$ of unitary type (\cref{pf:eis:GG}), and
deformations of its reduction carry an action of the multiplicative group by conjugation with
$\mathrm{diag}(t,1,t^{-1})$, that is, a grading. The group $\Sha_{\ell N}(\Q,\OO\otimes\eisQZ(\theta^{-1}\eiscyc^{-j}))$
injects into the part of degree $-1$ of the cotangent space of a suitable deformation ring at the Eisenstein
point (lower bound, by Lemma M$'$), and the length of that part is at most twice the sum of the valuations of
$\tr r(\sigma)-\tr\eisrz(\sigma)$ over representatives $r$ of the orbits of the other characteristic-zero points
(upper bound, by Lemma M, \cref{pf:vwin-he:lemmaM}). The hypotheses of Lemma M are verified using rigidity of
these points, which in turn rests on their automorphy and, at the endoscopic points, on $\sfname{KatoTwoW}$. In
case I the field is $K=K^*_\chi=\Q(\sqrt{-p_K(\chi)})$, the automorphic points are then counted by an Eisenstein
congruence number, and two auxiliary primes allow the bound to be summed over $\ell$. In case V the field is
$\Q(i)$ and the bound is for one prime at a time.

\emph{Organisation.} \Cref{sec:pf:eis:conv} fixes conventions and introduces the group $\eisGG$, its Eisenstein
representations, and the admissible three-dimensional representations. \Cref{sec:pf:eis:countI} proves
$\sfname{CountI}$ from a \emph{number statement} (an Eisenstein congruence number for admissible
representations, proved elsewhere) and a \emph{one-character statement}; the latter is proved in
\cref{sec:pf:eis:onechar}, where the comparison of a Selmer group with a graded cotangent space is carried out.
\Cref{sec:pf:eis:countV} proves $\sfname{CountV}$ by the same method over $\Q(i)$. \Cref{sec:pf:eis:dense} proves
the density statement for unitary groups by symmetric squares of Eisenstein-congruent forms, and
\cref{sec:pf:eis:sp} reduces $(\mathrm{Sp})$ to its five inputs. The proofs of the main results are complete
down to a list of precisely stated results of the formal development which we cite; among them are the
construction of the deformation rings, the verification of the conditions of Lemma M, an abstract counting
lemma for host frames, and the potential modularity and automorphic induction statements behind $(\mathrm{Sp})$.

% ---------------------------------------------------------------------------
\subsection{Conventions, the group \texorpdfstring{$\eisGG$}{G3} and admissible representations}\label{sec:pf:eis:conv}

\subsubsection{Galois groups, places and coefficients}

We keep the conventions of \cref{sec:pf:vwin-he:counts}: $\Omega$ is an algebraic closure of $\Q$ and
$G_\Q=\Gal(\Omega/\Q)$. For a number field $K$ given with an embedding into $\Omega$, $G_K=\Gal(\Omega/K)$ is
identified with the closed subgroup of $G_\Q$ fixing $K$. Places of $\Omega$ are valuation subrings $A\subset\Omega$,
with decomposition and inertia groups $D_A\supset I_A$ in $G_\Q$ (or in $G_K$); $A$ \emph{lies over} a prime $w$ of
$K$ if an element $x\in\OO_K$ satisfies $v_A(x)<1$ if and only if $x\in w$. An element $\sigma\in D_A\cap G_K$ is a
\emph{Frobenius element at $A$} if $v_A(\sigma x-x^{\Nm w})<1$ for all $x\in A$. For $M\ge1$, a homomorphism $\chi$ on
$G_K$ \emph{has conductor dividing $M$} (in the sense of the development) if it is trivial on the elements of $G_K$
fixing all $M$-th roots of unity, i.e.\ if it factors through $\Gal(K(\mu_M)/K)$.

$\Qlbar$ is an algebraic closure of $\Ql$ with the absolute value extending $|\ell|=\ell^{-1}$, and
\begin{equation}\label{pf:eis:vl}
\eisv(x)=-\log_\ell|x|\qquad(x\in\Qlbar),
\end{equation}
so that $\eisv(\ell)=1$; with the convention $\log0=0$ of the development one has $\eisv(0)=0$, but $\eisv$ is only
evaluated at non-zero arguments below. For $0\ne n\in\Z$, $\eisv(n)$ is the $\ell$-adic valuation of $n$. There exist
field isomorphisms $\iota:\Qlbar\cong\C$ (both fields are algebraically closed of characteristic $0$ and of the
cardinality of the continuum). Coefficient rings $\OO$, their ramification indices $e_\OO$, the lengths
$V_\ell(\theta,j)$, $\eisVt_\ell(\theta,j)=\|V_\ell(\theta,j)\|_\OO$ and the convention $\|\infty\|_\OO=0$ are those of
\cref{pf:vwin-he:coef,pf:vwin-he:V}; hosts $h$, the sets $\eisX_h(\OO,j)$, the components $\psi_{(r)}$, $\psi^\flat$,
the cases and the statements $\sfname{CountI}$, \dots, $\sfname{CountV}$ are those of
\cref{pf:vwin-he:host,pf:vwin-he:cases,pf:vwin-he:countdef}.

For a Dirichlet character $\theta$ modulo $N$ with values in $\OO$ we write $\theta_G=\theta\circ\kappa_N$ as in
\cref{pf:vwin-he:V}, and for $j\ge0$
\begin{equation}\label{pf:eis:Psi}
\Psi_{\theta,j}=\eiscyc^{\,j}\,\theta_G:G_\Q\to\OO^\times .
\end{equation}
For $\sigma\in G_\Q$ the \emph{Eisenstein trace} is $\Psi_{\theta,j}(\sigma)+1+\Psi_{\theta,j}(\sigma)^{-1}$. For
$q\ge1$, $e\ge0$ and $z\in\C^\times$ put
\begin{equation}\label{pf:eis:eis3}
\eiseis_3(q,e,z)=q^{-e}z+1+q^{e}z^{-1}.
\end{equation}
If $K$ is a number field, $w$ a prime of $K$ above a prime $p\nmid N\ell$ with $\Nm w=q$, and $\sigma\in G_K$ a
Frobenius element at a place above $w$, then $\kappa_N(\sigma)=q$ and $\eiscyc(\sigma)=q$, so that
$\Psi_{\theta,j}(\sigma)=q^j\theta(q)$, and $\iota$ carries the Eisenstein trace at $\sigma^{-1}$ to
$\eiseis_3(q,j,\iota(\theta(q))^{-1})$.

A representation $r:\Gamma\to\GL_n(k)$ of a group over a field is \emph{irreducible} if $k^n$ has no
$\Gamma$-stable subspace other than $0$ and $k^n$; if $\Gamma$ is topological, $r$ is \emph{irreducible on open
subgroups} if its restriction to every open subgroup is irreducible. Two representations $r,r'$ are
\emph{isomorphic} if $r'=grg^{-1}$ for some $g\in\GL_n(k)$.

\subsubsection{Automorphy and level}

\begin{definition}[Automorphy and level]\label{pf:eis:autlevel}
Let $F$ be a number field, $\ell$ a prime, $\iota:\Qlbar\cong\C$ and $r:G_F\to\GL_n(\Qlbar)$. We use the following
notions of the development; they are only needed to state the results below.
\begin{enumerate}[label=(\alph*)]
\item $r$ is \emph{attached via $\iota$} to a cuspidal automorphic representation $\pi$ of $\GL_n(\A_F)$ if at
  almost all finite places $v$, $\pi_v$ is unramified, $r$ is unramified at $v$ and the characteristic polynomial
  of $r(\Frob_v)$ corresponds under $\iota$ to the Satake parameters of $\pi_v$.
\item For $F$ CM, $r$ is \emph{automorphic of conjugate self-dual type} if it is attached via $\iota$ to a regular
  algebraic cuspidal $\pi$ with $\pi^\vee\cong\pi^c$; and $r$ satisfies the \emph{window condition} if, in addition,
  for every such $\pi$ attached to $r$ there is $t\in\R$ such that at all finite places $v\nmid\ell$, and at the places
  $v\mid\ell$ at which $r$ has a potentially unlinked flag (no power of a ratio of two diagonal characters equals the
  same power of $\eiscyc$), $\pi_v$ has cuspidal support in the diagonal torus and the characters $\xi$ of this
  support satisfy $\Nm v^{t-1/2}<|\xi(\varpi_v)|<\Nm v^{t+1/2}$; and at places $v\mid\ell$ an unlinked flag
  (defined in (c)) is potentially unlinked. We write
  $\mathcal A_\iota(r)$ for: $r$ is automorphic of conjugate self-dual type, and if $r$ is continuous it satisfies
  the window condition.
\item Let $M\ge1$ and $m=(m_1,\dots,m_n)\in\Z^n$. $r$ has \emph{level $M$ and type $m$} if $r$ is continuous; at each
  finite place $w\nmid\ell$ some inertia group at $w$ acts on $\Qlbar^n$ diagonally through characters of conductor
  dividing $M$; and at each $w\mid\ell$, on some decomposition group $D_A$ at $w$, $r$ is upper triangular with
  diagonal characters $\kappa_1,\dots,\kappa_n$ such that $\kappa_i=\eiscyc^{m_i}\chi_i$ on $I_A$ with $\chi_i$ of
  conductor dividing $M$, and either $\kappa_i\ne\kappa_j\eiscyc$ for all $i,j$, or $r$ is potentially crystalline at
  $A$. Such a flag is \emph{unlinked} if $\kappa_i\ne\kappa_j\eiscyc$ for all $i,j$; $r$ is \emph{unlinked of level $M$ and
  type $m$} if at every $w\mid\ell$ such an unlinked flag exists.
\item A newform $f$ of weight $k$ is \emph{attached via $\iota$} to $V:G_\Q\to\GL_2(\Qlbar)$ if $V$ is semisimple and
  for almost all primes $p$, $V$ is unramified at $p$ and $\iota$ carries the characteristic polynomial of
  $V(\Frob_p^{-1})$ to the Hecke polynomial of $f$ at $p$.
\end{enumerate}
\formal{UpperTheory.V.Z.Win.IsAutCsdA}
\formal{UpperTheory.V.Af.Level}
\formal{UpperTheory.V.Th.Unl}
\end{definition}

\subsubsection{The group \texorpdfstring{$\eisGG$}{G3} and its Eisenstein representations}

\begin{definition}[The group $\eisGG$, weights and Eisenstein representations]\label{pf:eis:GG}
Let $R$ be a commutative ring and $J\in\GL_3(R)$ the antidiagonal matrix with entries $1$. Let
$\vartheta(g)=J\,{}^tg^{-1}J$ for $g\in\GL_3(R)$, an involutive automorphism, and
$\eisGG(R)=\GL_3(R)\rtimes\{\pm1\}$, where $-1$ acts by $\vartheta$. For $x=(g,s)$ write $\mathrm{mat}(x)=g$ and
$\nu(x)=s$. The \emph{weights} of the basis vectors are $w(1)=1$, $w(2)=0$, $w(3)=-1$, and the weight of the matrix
entry $(a,b)$ is $w(a,b)=w(a)-w(b)$. For $t\in R^\times$ let $\tau(t)=\mathrm{diag}(t,1,t^{-1})$; conjugation by $\tau(t)$
multiplies the $(a,b)$ entry by $t^{w(a,b)}$, and $\vartheta(\tau(t))=\tau(t)$.

Let $\Gamma$ be a group, $\eta:\Gamma\to\{\pm1\}$ and $\Psi:\Gamma\to R^\times$ homomorphisms. The \emph{Eisenstein
representation} is
\[
 \eisrz(\eta,\Psi):\Gamma\to\eisGG(R),\qquad g\mapsto\bigl(\tau(\Psi(g)),\eta(g)\bigr).
\]
For a homomorphism $r:\Gamma\to\eisGG(R)$ we say that $r$ \emph{has multiplier $\eta$} if $\nu\circ r=\eta$, and we write
$r_K$ for the restriction of $\mathrm{mat}\circ r$ to $\ker\eta$ (a homomorphism $\ker\eta\to\GL_3(R)$). Assume $R$ is a field and
$r$ has multiplier $\eta$. Then $r$ is
\begin{itemize}
\item \emph{upper} (resp.\ \emph{lower}) if the entries of $\mathrm{mat}(r(g))$ of negative (resp.\ positive) weight vanish for
  $g\in\ker\eta$;
\item \emph{of type I} if $r_K$ is irreducible;
\item \emph{of type II} if the entries of odd weight of $\mathrm{mat}(r(g))$ vanish for $g\in\ker\eta$ and the $2\times2$ matrix
  formed by the entries $(1,1),(1,3),(3,1),(3,3)$ (the \emph{corner}) is an irreducible representation of $\ker\eta$.
\end{itemize}
The \emph{adjoint representation} of $\eisGG(R)$ on $M_3(R)$ is $x\cdot X=\mathrm{mat}(x)\,X'\,\mathrm{mat}(x)^{-1}$, where $X'=X$ if
$\nu(x)=1$ and $X'=-J\,{}^tX\,J$ if $\nu(x)=-1$. Let $\mathfrak g_j\subset M_3(R)$ be the matrices supported on entries of
weight $j$.
\formal{UpperTheory.V.U3.r0}
\end{definition}

The restriction of $\eisrz(\eta,\Psi)$ to $\ker\eta$ is $\Psi\oplus1\oplus\Psi^{-1}$. If $K$ is an imaginary quadratic
field with quadratic character $\eta$ of $G_\Q$ and $r:G_\Q\to\eisGG(R)$ has multiplier $\eta$, then $r_K$ is a
representation of $G_K$; this is the passage between representations of $G_\Q$ into $\eisGG$ and three-dimensional
representations of $G_K$ of unitary type used throughout.

\begin{definition}[Admissible three-dimensional representations]\label{pf:eis:adm}
Let $K$ be an imaginary quadratic field, $\ell$ a prime, $\iota:\Qlbar\cong\C$, $e,N,M\ge1$ and
$\rho:G_K\to\GL_3(\Qlbar)$.
\begin{enumerate}[label=(\alph*)]
\item $\rho$ is \emph{admissible of unitary type} (for $\iota,e,N$) if $\rho$ is irreducible, $\mathcal A_\iota(\rho\otimes\eiscyc^{-1})$
  holds, and $\rho$ has level $N$ and type $(e,0,-e)$ and is unlinked of level $N$ and type $(e,0,-e)$.
  \formal{UpperTheory.V.Z.J.Adm3U_s}
\item $\rho$ is \emph{endoscopic} (for $\iota,e,M$) if there is an irreducible $W:G_\Q\to\GL_2(\Qlbar)$ such that a
  newform of weight $2e+1$ is attached via $\iota$ to $W\otimes\eiscyc^{-e}$; $W$ satisfies the window condition
  $\sfname{WitAQ}(\iota,W)$ of \cref{def:window:kato}; $W$ has level $M$ and type $(e,-e)$; $\det W$ is a character of
  finite order with conductor dividing some $M'\ge1$; and $\rho$ is isomorphic to
  \[
   s\longmapsto\begin{pmatrix}W(s)_{11}&0&W(s)_{12}\\0&1&0\\W(s)_{21}&0&W(s)_{22}\end{pmatrix}.
  \]
  \formal{UpperTheory.V.Z.Win.Dn.Adm2A}
\end{enumerate}
We say that $\rho$ is \emph{admissible} if it satisfies (a) for $(\iota,e,N)$ or (b) for $(\iota,e,N|d_K|)$.
\end{definition}

The endoscopic representations are the type II points of the deformation problems below; the admissible
representations of unitary type are the type I points.

% ---------------------------------------------------------------------------
\subsection{Case I: from one character to all primes}\label{sec:pf:eis:countI}

\begin{theorem}[\cref{thm:window:countI}]\label{pf:eis:main}
Assume $\sfname{HypothesesH'}$, $\sfname{LemmaK}(C_K)$ for some $C_K\in\R$, $\sfname{KatoTwoW}$ and the unitary automorphy
statement with unramified coefficients $(\mathrm{U}_{\mathrm{ur}})$ (\cref{pf:eis:Uur}). Then $\sfname{CountI}$ holds.
\formallabel{UpperTheory.V.Z.J.count_i_of_printed_bound_k_kato_automorphy}
\end{theorem}

The proof is given at the end of \cref{sec:pf:eis:onechar}. In this subsection we reduce $\sfname{CountI}$ to two
statements about a single imaginary quadratic field $K$ and a single prime $p$: a number statement, which is proved
elsewhere in the development, and a one-character statement, which is proved in \cref{sec:pf:eis:onechar}.

\subsubsection{Genericity, and the two statements}

\begin{definition}[Genericity conditions]\label{pf:eis:GD}
Let $K\subset\Omega$ be a number field, $\ell$ a prime, $e\in\N$, $N\ge1$, and $\psi$ a Dirichlet character modulo $N$ with
values in a commutative ring. We say that $\mathrm{GD}(K,\ell,e,\psi)$ holds if:
\begin{enumerate}[label=(G\arabic*)]
\item $\ell\ne2$, $e\ge1$ and $N$ is even;
\item $\ell\nmid\ord\psi$ and $\psi_{(\ell)}^2=1$;
\item $\psi(-1)=(-1)^e$;
\item $K$ is imaginary quadratic, $\ell\nmid|d_K|h_K$ and $\gcd(|d_K|,N)=1$;
\item the subgroup $\ker(\eiscyc\bmod\ell)\cap\ker\bigl((\psi^\flat)_G\bigr)$ of $G_\Q$ is not contained in $G_K$;
\item $(\ell-1)\nmid2e$.
\end{enumerate}
\formal{UpperTheory.V.HV.GD}
\end{definition}

Condition (G5) says that $K$ is not contained in the field cut out by $\eiscyc\bmod\ell$ and $\psi^\flat$. For a
subgroup $H\subset(\Z/N)^\times$ and a coefficient ring $\OO$ let $\mathrm{orb}(\OO,H)$ be the finite set of Dirichlet
characters modulo $N$ with values in $\OO$ whose kernel (on $(\Z/N)^\times$) is exactly $H$.

\begin{definition}[Number statement]\label{pf:eis:number}
$(\mathrm{Num})$ is the statement: for every imaginary quadratic field $K$ there are $c,A\in\N$ with the following
property. Let $e\ge1$, $N\ge1$, $p\ge7$ a prime with $p\nmid N$ and $p\nmid d_K$, $w$ a prime of $K$ above $p$, and
$(z_o)_{o\in\mathcal O}$ a finite family of complex roots of unity which is stable under $\Aut(\C)$ (for every
$\sigma\in\Aut(\C)$ there is a permutation $u$ of $\mathcal O$ with $z_{u(o)}=\sigma(z_o)$). Then there is an integer
$\mathcal N\ne0$ with
\[
 \log|\mathcal N|\le c\cdot\#\mathcal O\cdot(N(e+1))^A\log p
\]
such that for every prime $\ell\notin\{2,p\}$, every $\Omega$ and $\iota:\Qlbar\cong\C$, every place $B$ of $\Omega$ over $w$, every
Frobenius element $\sigma\in G_K$ at $B$, and every family $(X_o)_{o\in\mathcal O}$ of finite sets of admissible
representations $G_K\to\GL_3(\Qlbar)$ (\cref{pf:eis:adm}, for $\iota,e,N$), pairwise non-isomorphic within each
$X_o$, one has $\tr\rho(\sigma^{-1})\ne\iota^{-1}\eiseis_3(\Nm w,e,z_o)$ for $o\in\mathcal O$, $\rho\in X_o$, and
\[
 |\mathcal N|_\ell\ \le\ \prod_{o\in\mathcal O}\ \prod_{\rho\in X_o}\bigl|\tr\rho(\sigma^{-1})-\iota^{-1}\eiseis_3(\Nm w,e,z_o)\bigr|_\ell .
\]
\formal{UpperTheory.V.Z.J.NumberU3U_s}
\end{definition}

$(\mathrm{Num})$ follows from the part of $\sfname{HypothesesH'}$ consisting of classical theorems
(\lref{UpperTheory.GalN.V.Win.N.numberU3U_of_prints}: this part of the package implies $(\mathrm{Num})$). It is the
\emph{Eisenstein congruence number} of step (iv) of \cref{sec:window:eis}.

\begin{definition}[One-character statement]\label{pf:eis:onecharstmt}
$(\mathrm{One})$ is the statement: let $K\subset\Omega$ be imaginary quadratic, $\ell$ a prime, $\OO$ a coefficient ring,
$N\ge1$, $e\in\N$ and $\psi$ a Dirichlet character modulo $N$ with values in $\OO$ such that $\mathrm{GD}(K,\ell,e,\psi)$
holds; put $\Psi=\Psi_{\psi,e}$. Let $\iota:\Qlbar\cong\C$, $\sigma\in G_K$ and $t=\Psi(\sigma)+1+\Psi(\sigma)^{-1}$, and assume
that $\tr\rho(\sigma)\ne t$ for every admissible $\rho:G_K\to\GL_3(\Qlbar)$. Then $V_\ell(\psi,e)<\infty$, and there is a
finite set $X$ of pairwise non-isomorphic admissible representations such that
\[
 \eisVt_\ell(\psi,e)\ \le\ 2\sum_{\rho\in X}\eisv\bigl(\tr\rho(\sigma)-t\bigr).
\]
\formallabel{UpperTheory.V.Z.J.OneCharU_s}
\end{definition}

\begin{definition}[$\sfname{PropGK}$ and $\sfname{PropGK}^{\Sigma}$]\label{pf:eis:propGK}
\leavevmode
\begin{enumerate}[label=(\alph*)]
\item $\sfname{PropGK}$: for every number field $K$ there are $c_5,A\in\N$ with the following property. Let $e,N\ge1$,
  $H\subset(\Z/N)^\times$ a subgroup and $p$ a prime with $p\nmid30N|d_K|$. There is an integer $\mathcal N\ne0$ with
  $\log|\mathcal N|\le c_5\,\varphi([(\Z/N)^\times:H])\,(N(e+1))^A\log p$ such that for every $\Omega\supset K$, every prime
  $\ell\ne p$ and every coefficient ring $\OO$ for $\ell$: if some $\psi\in\mathrm{orb}(\OO,H)$ satisfies
  $\mathrm{GD}(K,\ell,e,\psi)$, then $V_\ell(\psi,e)<\infty$ for all $\psi\in\mathrm{orb}(\OO,H)$ and
  \[\textstyle\sum_{\psi\in\mathrm{orb}(\OO,H)}\eisVt_\ell(\psi,e)\le2\,\eisv(\mathcal N).\]
\item $\sfname{PropGK}^{\Sigma}$: for every number field $K$ there are $c_5,A\in\N$ with the following property. Let $\Omega\supset K$,
  $e,N\ge1$, $H\subset(\Z/N)^\times$, $p\ne p'$ primes not dividing $30N|d_K|$, $\OO_\ell$ a coefficient ring for each prime
  $\ell$, and $L$ a finite set of primes such that for each $\ell\in L$ some $\psi\in\mathrm{orb}(\OO_\ell,H)$ satisfies
  $\mathrm{GD}(K,\ell,e,\psi)$. Then $V_\ell(\psi,e)<\infty$ for $\ell\in L$ and $\psi\in\mathrm{orb}(\OO_\ell,H)$, and
  \[
   \sum_{\ell\in L}\log\ell\sum_{\psi\in\mathrm{orb}(\OO_\ell,H)}\eisVt_\ell(\psi,e)\ \le\
   2c_5\,\varphi([(\Z/N)^\times:H])\,(N(e+1))^A(\log p+\log p').
  \]
\end{enumerate}
\formallabel{UpperTheory.V.U3.PropGK}
\formallabel{UpperTheory.V.U3.CorGK}
\end{definition}

\subsubsection{Characters with a given kernel}

\begin{lemma}\label{pf:eis:orb}
Let $\OO$ be a domain, $N\ge1$ and $\psi$ a Dirichlet character modulo $N$ with values in $\OO$. Then
$[(\Z/N)^\times:\ker\psi]=\ord\psi$; the quotient $(\Z/N)^\times/\ker\psi$ is cyclic; a character $\psi'$ modulo $N$ with
values in $\OO$ has $\ker\psi'=\ker\psi$ if and only if $\psi'=\psi^j$ for some $j$ prime to $\ord\psi$; and there are
exactly $\varphi(\ord\psi)$ such characters.
\formallabel{UpperTheory.V.HV.orbFacts}
\end{lemma}

\begin{proof}
$\psi$ induces an injection of $C=(\Z/N)^\times/\ker\psi$ into $\OO^\times$, and finite subgroups of the unit group of a
domain are cyclic; so $C$ is cyclic, of order $n=\#\psi((\Z/N)^\times)=\ord\psi$. Let $g$ generate $C$; then $\zeta=\psi(g)$
is a primitive $n$-th root of unity, and the $n$-th roots of unity of $\OO$ are the powers of $\zeta$ (a domain has at
most $n$ of them). If $\ker\psi'=\ker\psi$, then $\psi'$ is injective on $C$, so $\psi'(g)$ has order $n$ and equals
$\zeta^j$ with $\gcd(j,n)=1$, whence $\psi'=\psi^j$; the converse is clear. Finally $j\mapsto\psi^j$ induces a bijection
from $(\Z/n)^\times$ onto the set of these characters.
\end{proof}

\begin{lemma}\label{pf:eis:orbsum}
Let $\Gamma$ be a finite abelian group and $T$ a finite set of subgroups $H\subset\Gamma$ with $\Gamma/H$ cyclic. Then
$\sum_{H\in T}\varphi([\Gamma:H])\le\#\Gamma$.
\formallabel{UpperTheory.V.HV.orbSum}
\end{lemma}

\begin{proof}
If $\Gamma/H$ is cyclic of order $n$, the characters $\Gamma\to\C^\times$ with kernel exactly $H$ correspond to the
injective characters of $\Gamma/H$, which send a generator to a primitive $n$-th root of unity; there are $\varphi(n)$ of
them. The sets of characters with kernel $H$, $H\in T$, are disjoint, and $\#\Hom(\Gamma,\C^\times)=\#\Gamma$.
\end{proof}

\begin{lemma}\label{pf:eis:gdorb}
If $\mathrm{GD}(K,\ell,e,\psi)$ holds and $\psi'$ is a character modulo $N$ with values in the same domain and
$\ker\psi'=\ker\psi$, then $\mathrm{GD}(K,\ell,e,\psi')$ holds.
\formallabel{UpperTheory.V.HV.gdFacts}
\end{lemma}

\begin{proof}
By \cref{pf:eis:orb}, $\psi'=\psi^j$ with $\gcd(j,\ord\psi)=1$, and $\ord\psi'=\ord\psi$. Conditions (G1), (G4), (G6) do not
involve $\psi$. (G2): $\ell\nmid\ord\psi'$, and $\psi'_{(\ell)}=\psi_{(\ell)}^j$ since $\chi\mapsto\chi_{(\ell)}$ is a
homomorphism, so $\psi'^2_{(\ell)}=1$. (G3): $\psi(-1)=\pm1$; if $\psi(-1)=-1$ then $\ord\psi$ is even and $j$ is odd, so
$\psi'(-1)=\psi(-1)^j=\psi(-1)$. (G5): the decomposition of a character into its parts of $\ell$-power order and of order
prime to $\ell$ commutes with taking $j$-th powers, and so does $\chi\mapsto\chi_{(\ell)}$; hence
$\psi'^\flat=(\psi^\flat)^j$. As $\ord\psi^\flat$ divides $\ord\psi$, $j$ is prime to it and $\ker(\psi'^\flat)_G=\ker(\psi^\flat)_G$.
\end{proof}

\begin{remark}\label{pf:eis:gdrem}
The formal statement \cref{pf:eis:gdorb} contains in addition several elementary consequences of
$\mathrm{GD}(K,\ell,e,\psi)$ which are not used in this section: $\ell\ge5$, $\psi$ has order prime to $\ell$,
$\psi=\psi^\flat\psi_{(\ell)}$, $(\psi^\flat)_{(\ell)}=1$, $\ell\nmid\ord\psi^\flat$, and if $\psi_{(\ell)}\ne1$ then $\ell\mid N$ and
$\psi_{(\ell)}$ has conductor $\ell$; moreover $\psi'_{(\ell)}=\psi_{(\ell)}$ for $\psi'$ as in the lemma.
\end{remark}

\begin{lemma}[Genericity in case I]\label{pf:eis:gdhost}
Let $h=(\chi,m,q_1,d)$ be a host which is big for $t$, $K^*=K^*_\chi=\Q(\sqrt{-p_K(\chi)})\subset\Omega$, $\ell\ne2$ a prime and
$0<e<t$ such that $(\ell,e)$ is in case I. Let $\theta$ be a Dirichlet character modulo $N_h$ with values in a domain,
trivial on $H_h$, with $\theta(-1)=(-1)^e$. Then $\mathrm{GD}(K^*,\ell,e,\theta)$ holds.
\formallabel{UpperTheory.V.HV.gdOfHost}
\end{lemma}

\begin{proof}
Write $p_K=p_K(\chi)$, a prime with $p_K\equiv3\bmod4$ (\lref{UpperTheory.V.pK_pOne_prime_conditions}: $p_K(\chi)$ is a
prime $\equiv3\bmod4$, $\ge7$, prime to $N_\chi$). By \lref{UpperTheory.V.HV.kstar} ($K^*$ has discriminant $-p_K$, is
totally complex and of degree $2$) $K^*$ is imaginary quadratic with $|d_{K^*}|=p_K$. We use two shared facts on
hosts: \lref{UpperTheory.V.HV.bchiFacts} (for $h$ big for $t$: $p_K\in B_\chi$, $p_K\nmid N_h$, and every odd prime
$\ell\notin B_\chi$ satisfies $\ell\ne p_K$ and $\ell\nmid h^*(\chi)$, besides other properties) and
\lref{UpperTheory.V.HV.cases} (for $h$ big for $t$, $\ell\ne2$, $0<e<t$ and $\theta$ trivial on $H_h$: if $\ell\notin B_\chi$ then
$\ell\nmid\ord\theta$ and $\theta_{(\ell)}^2=1$, besides other properties). Now (G1) holds since $2^{m+2}\mid N_h$; (G2) by
\lref{UpperTheory.V.HV.cases}, as $\ell\notin B_\chi$ in case I; (G3) by hypothesis; (G4) since $\ell\ne p_K$,
$\ell\nmid h^*(\chi)=h_{K^*}$ and $p_K\nmid N_h$; (G6) is part of case I. For (G5) let $A$ be a place of $\Omega$ above $p_K$. Since
$K^*$ is ramified at $p_K$, $I_A\not\subset G_{K^*}$ (\lref{UpperTheory.V.HV.kstarRam}). Since $\ell\ne p_K$ and $p_K\nmid N_h$, the
cyclotomic characters $\eiscyc$ and $\kappa_{N_h}$ are trivial on $I_A$, so $I_A\subset\ker(\eiscyc\bmod\ell)\cap\ker(\theta^\flat)_G$,
and this subgroup is not contained in $G_{K^*}$.
\end{proof}

\subsubsection{The reductions}

\begin{proposition}\label{pf:eis:propGKpf}
$(\mathrm{Num})$ and $(\mathrm{One})$ imply $\sfname{PropGK}$.
\formallabel{UpperTheory.V.Z.J.propGK_of_numberU3U_of_oneCharU}
\end{proposition}

\begin{proof}
Let $K$ be a number field. If $K$ is not imaginary quadratic, take $c_5=A=0$ and $\mathcal N=1$; the hypothesis of the
conclusion is never satisfied, since $\mathrm{GD}$ requires $K$ to be imaginary quadratic. Otherwise let $c,A$ be the
constants of $(\mathrm{Num})$ for $K$, and put $c_5=c$.

Let $e,N,H,p$ be given with $p\nmid30N|d_K|$; then $p\nmid N$, $p\nmid d_K$ and $p\ge7$. Choose a prime $w$ of $K$ above $p$; its
norm $q=\Nm w$ is a power of $p$, hence prime to $N$. Let $\mathcal O$ be the set of complex Dirichlet characters $\chi$
modulo $N$ with kernel $H$, and $z_\chi=\chi(q)^{-1}$, a root of unity; for $\sigma\in\Aut(\C)$, $\chi\mapsto\sigma\circ\chi$ permutes
$\mathcal O$ and $z_{\sigma\circ\chi}=\sigma(z_\chi)$. By \cref{pf:eis:orb}, $\#\mathcal O\le\varphi([(\Z/N)^\times:H])$ ($\mathcal O$ is empty
or consists of the $\varphi(\ord\chi)$ characters with the kernel of one of its elements). $(\mathrm{Num})$ gives
$\mathcal N\ne0$ with $\log|\mathcal N|\le c\,\#\mathcal O\,(N(e+1))^A\log p\le c_5\varphi([(\Z/N)^\times:H])(N(e+1))^A\log p$, and
the last assertion of \cref{pf:eis:number}, which we call (P), for all $\ell\notin\{2,p\}$.

Now let $\Omega$, $\ell\ne p$ and $\OO$ be given, and $\psi_0\in\mathrm{orb}(\OO,H)$ with $\mathrm{GD}(K,\ell,e,\psi_0)$; in particular
$\ell\ne2$. Choose $\iota:\Qlbar\cong\C$, a place $B$ of $\Omega$ over $w$ and a Frobenius element $\sigma\in G_K$ at $B$. The map
$E:\psi\mapsto\iota\circ\psi$ is a bijection from $\mathrm{orb}(\OO,H)$ onto $\mathcal O$: it is injective and maps into $\mathcal O$ since
$\iota$ is injective, and every $\chi\in\mathcal O$ is a power $(\iota\circ\psi_0)^j=E(\psi_0^j)$ with $j$ prime to the order, by
\cref{pf:eis:orb}. Let $\psi\in\mathrm{orb}(\OO,H)$. By \cref{pf:eis:gdorb}, $\mathrm{GD}(K,\ell,e,\psi)$ holds. Put
$t_\psi=\iota^{-1}\eiseis_3(q,e,z_{E(\psi)})$; by the computation following \eqref{pf:eis:eis3},
$t_\psi=\Psi_{\psi,e}(\sigma^{-1})+1+\Psi_{\psi,e}(\sigma^{-1})^{-1}$. Applying (P) to the families with $X_o=\{\rho\}$ for one
admissible $\rho$ shows $\tr\rho(\sigma^{-1})\ne t_\psi$ for every admissible $\rho$. Hence $(\mathrm{One})$, applied to $\psi$, $\iota$
and $\sigma^{-1}$, gives $V_\ell(\psi,e)<\infty$ and a finite set $X_\psi$ of pairwise non-isomorphic admissible representations
with $\eisVt_\ell(\psi,e)\le2\sum_{\rho\in X_\psi}\eisv(\tr\rho(\sigma^{-1})-t_\psi)$.

Apply (P) to the family $X_o=X_{E^{-1}(o)}$: all the differences are non-zero and
$|\mathcal N|_\ell\le\prod_o\prod_{\rho\in X_o}|\tr\rho(\sigma^{-1})-t_{E^{-1}(o)}|_\ell$. Since $|\mathcal N|_\ell=\ell^{-\eisv(\mathcal N)}$, taking
$-\log_\ell$ gives $\sum_o\sum_{\rho\in X_o}\eisv(\tr\rho(\sigma^{-1})-t_{E^{-1}(o)})\le\eisv(\mathcal N)$. Summing the bounds for
$\eisVt_\ell(\psi,e)$ over $\psi$ and reindexing by $E$,
\[
 \sum_{\psi\in\mathrm{orb}(\OO,H)}\eisVt_\ell(\psi,e)\le2\sum_{o\in\mathcal O}\sum_{\rho\in X_o}\eisv\bigl(\tr\rho(\sigma^{-1})-t_{E^{-1}(o)}\bigr)
 \le2\,\eisv(\mathcal N).\qedhere
\]
\end{proof}

\begin{proposition}\label{pf:eis:corGK}
$\sfname{PropGK}$ implies $\sfname{PropGK}^{\Sigma}$, with the same constants.
\formallabel{UpperTheory.V.U3.corGK}
\end{proposition}

\begin{proof}
Let $K$, $c_5$, $A$ be as in $\sfname{PropGK}$, and let $\Omega,e,N,H,p,p',(\OO_\ell),L$ be as in $\sfname{PropGK}^{\Sigma}$. If $L=\emptyset$ the sum
vanishes and the right-hand side is $\ge0$. Otherwise $e\ge1$ by (G1), and $\sfname{PropGK}$ gives integers $\mathcal N\ne0$ for
$p$ and $\mathcal N'\ne0$ for $p'$. For $\ell\in L$ we have $\ell\ne p$ or $\ell\ne p'$; applying $\sfname{PropGK}$ with the
corresponding integer gives finiteness and
$\sum_{\psi\in\mathrm{orb}(\OO_\ell,H)}\eisVt_\ell(\psi,e)\le2\eisv(\mathcal N)+2\eisv(\mathcal N')$. For $0\ne n\in\Z$ and a finite set $L$ of
primes, $\prod_{\ell\in L}\ell^{\eisv(n)}$ divides $n$, so
\begin{equation}\label{pf:eis:sumlog}
\textstyle\sum_{\ell\in L}\log\ell\cdot\eisv(n)\le\log|n|.
\end{equation}
Therefore the left-hand side of $\sfname{PropGK}^{\Sigma}$ is at most
$2\log|\mathcal N|+2\log|\mathcal N'|\le2c_5\varphi([(\Z/N)^\times:H])(N(e+1))^A(\log p+\log p')$.
\end{proof}

\begin{proposition}\label{pf:eis:countIcor}
$\sfname{PropGK}^{\Sigma}$ implies $\sfname{CountI}$.
\formallabel{UpperTheory.V.U3.countI}
\end{proposition}

\begin{proof}
Let $\chi$ be a primitive character, $p_K=p_K(\chi)$ and $K^*=\Q(\sqrt{-p_K})$, imaginary quadratic with $|d_{K^*}|=p_K$
(\lref{UpperTheory.V.HV.kstar}). Let $c_5,A$ be the constants of $\sfname{PropGK}^{\Sigma}$ for $K^*$. Choose primes
$r_1\ge30N_\chi p_K+1$ and $r_2>r_1$, and $k\in\N$ with $2c_5(\log r_1+\log r_2)\le16^k$. We show that $C=k+1+A$, $r_1$, $r_2$
satisfy $\sfname{CountI}$ for $\chi$.

Let $h=(\chi,m,q_1,d)$ be a host with $r_1,r_2<\min(q_1,d)$, big for $t$, let $0<e<t$, $\Omega$, coefficient rings $\OO_\ell$, and a
finite set $L$ of odd primes in case I be given; fix an embedding $K^*\to\Omega$. Write $N=N_h$ and
$G=(\Z/N)^\times$.

\emph{The auxiliary primes.} Let $r\in\{r_1,r_2\}$. Then $r\nmid30N p_K$: indeed $r>30$; the prime divisors of $N$ are $2$,
the prime divisors of $N_\chi$, $q_1$ and $d$ (\lref{UpperTheory.V.HV.hostIndex}, which states among other facts that
$16\mid N_h$ and that a prime divides $N_h$ if and only if it is $2$, divides $N_\chi$, or equals $q_1$ or $d$), and $r>N_\chi$,
$r<q_1$, $r<d$; and $r>p_K$. As $|d_{K^*}|=p_K$, the primes $r_1\ne r_2$ satisfy the hypotheses of $\sfname{PropGK}^{\Sigma}$.

\emph{Grouping by kernels.} Let $T$ be the finite set of subgroups $\ker\theta\subset G$ for $\theta\in\eisX_h(\OO_\ell,e)$,
$\ell\in L$; for $H\in T$ let $L_H$ be the set of $\ell\in L$ such that some $\theta\in\eisX_h(\OO_\ell,e)$ has kernel $H$. For
$\ell\in L_H$ such a $\theta$ lies in $\mathrm{orb}(\OO_\ell,H)$ and satisfies $\mathrm{GD}(K^*,\ell,e,\theta)$ by \cref{pf:eis:gdhost}.
Hence $\sfname{PropGK}^{\Sigma}$ applies to $(H,r_1,r_2,(\OO_\ell),L_H)$. It gives the finiteness assertion of $\sfname{CountI}$ (every
$\theta\in\eisX_h(\OO_\ell,e)$ lies in $\mathrm{orb}(\OO_\ell,\ker\theta)$ with $\ell\in L_{\ker\theta}$), and
\begin{equation}\label{pf:eis:perH}
 \sum_{\ell\in L_H}\log\ell\sum_{\psi\in\mathrm{orb}(\OO_\ell,H)}\eisVt_\ell(\psi,e)\le2c_5\varphi([G:H])(N(e+1))^A(\log r_1+\log r_2).
\end{equation}
Since $\eisVt\ge0$ and $\{\theta\in\eisX_h(\OO_\ell,e):\ker\theta=H\}\subset\mathrm{orb}(\OO_\ell,H)$, sorting the characters by their kernels gives
\[
 \sum_{\ell\in L}\log\ell\sum_{\theta\in\eisX_h(\OO_\ell,e)}\eisVt_\ell(\theta,e)\le\sum_{H\in T}\ \sum_{\ell\in L_H}\log\ell
 \sum_{\psi\in\mathrm{orb}(\OO_\ell,H)}\eisVt_\ell(\psi,e).
\]
By \cref{pf:eis:orb} each $G/H$, $H\in T$, is cyclic, so \cref{pf:eis:orbsum} gives $\sum_{H\in T}\varphi([G:H])\le\#G\le N$.
With \eqref{pf:eis:perH} the left-hand side of $\sfname{CountI}$ is therefore at most
\[
 2c_5(\log r_1+\log r_2)\cdot N\cdot(N(e+1))^A\le N^k\cdot N\cdot(Nt)^A\le(Nt)^{k+1+A},
\]
using $2c_5(\log r_1+\log r_2)\le16^k\le N^k$ (as $16\mid N$), $e+1\le t$ and $t\ge1$.
\end{proof}

% ---------------------------------------------------------------------------
\subsection{The one-character bound}\label{sec:pf:eis:onechar}

In this subsection we prove $(\mathrm{One})$ (\cref{pf:eis:onecharthm}) and then \cref{pf:eis:main}. Throughout,
$K\subset\Omega$ is an imaginary quadratic field, $\ell$ a prime, $\OO$ a coefficient ring for $\ell$ with residue field
$k$, $N\ge1$, $e\in\N$ and $\psi$ a Dirichlet character modulo $N$ with values in $\OO$ such that
$\mathrm{GD}(K,\ell,e,\psi)$ holds; $\Psi=\Psi_{\psi,e}$, and $\Psi_{\Qlbar}$ is its composite with $\OO\subset\Qlbar$. A
\emph{complete $\OO$-algebra} is a complete local $\OO$-algebra $R$ as in \cref{sec:pf:vwin-he:lemmaM} (Noetherian,
complete, with residue field $k$), and $\mathfrak m_R$ is its maximal ideal.

\subsubsection{Data and deformation rings}

\begin{definition}[Data]\label{pf:eis:dat}
A \emph{datum} $D$ for $(K,\OO,e,\psi)$ consists of: a homomorphism $\eta:G_\Q\to\{\pm1\}$ with kernel $G_K$; an element
$c\in G_\Q$ which is a complex conjugation (there is an embedding $\phi:\Omega\to\C$ with $\phi(cx)=\overline{\phi(x)}$); an
embedding $j$ of $\Omega$ into an algebraically closed complete extension of $\Ql$ and the continuous homomorphism
$\iota_\ell:G_{\Ql}\to G_\Q$ with $\sigma(j(x))=j(\iota_\ell(\sigma)x)$; a finite extension $K_v$ of $\Ql$ with
$\eta(\iota_\ell(g))=1$ if and only if $g\in G_{K_v}$; such that $\eta\circ\iota_\ell$ is trivial on the inertia group
$I_{\Ql}$, the kernel of $(\psi_{(\ell)})_G\circ\iota_\ell$ is open, all powers of $\Psi$ are continuous for the
$\mathfrak m_\OO$-adic topology, and $\Psi(\iota_\ell(g))=\eiscyc(g)^e(\psi_{(\ell)})_G(\iota_\ell(g))$ for $g\in I_{\Ql}$.
\formal{UpperTheory.V.U3.Dat}
\end{definition}

The image of $\iota_\ell$ is the decomposition group of a place of $\Omega$ above $\ell$, and $\iota_\ell$ maps $I_{\Ql}$ onto its
inertia group (\lref{UpperTheory.V.U3.dat_loc}). By
\lref{UpperTheory.V.U3.exists_dat}, $\mathrm{GD}(K,\ell,e,\psi)$ implies that a datum exists. Fix a datum $D$, and
let $\bar r_0=\eisrz(\eta,\bar\Psi)$ be the reduction of $\eisrz(\eta,\Psi)$ modulo $\mathfrak m_\OO$. A \emph{lift} of
$\bar r_0$ to a complete $\OO$-algebra $R$ is a homomorphism $r:G_\Q\to\eisGG(R)$ whose matrix entries are continuous for
the $\mathfrak m_R$-adic topology and which reduces to $\bar r_0$; lifts of $\bar r_0\circ\iota_\ell$ are defined in the same
way. For a property $P$ of lifts, a \emph{universal ring} for $P$ is a complete $\OO$-algebra $R$ with a lift $r$ such that
for every complete $\OO$-algebra $R'$ the lifts $f\circ r$, $f:R\to R'$, have property $P$, and every lift to $R'$ with
property $P$ is $f\circ r$ for a unique $f$.

Let $\eisKill_D\subset G_\Q$ be the set of elements of inertia groups $I_A$ at places $A$ not above $\ell$ with
$\bar r_0(g)=1$. An element $g_1\in G_\Q$ is \emph{marked} if $c^2=1$, $\eta(c)=-1$, $\Psi(c)=1$, $\eta(g_1)=1$ and
$\Psi(g_1)^2-1\in\OO^\times$. A homomorphism $r:G_\Q\to\eisGG(R)$ is \emph{normalized} at $(c,g_1)$ if $\mathrm{mat}(r(c))=1$
and the entries $(1,2)$ and $(2,1)$ of $\mathrm{mat}(r(g_1))-\mathrm{mat}(r(cg_1^{-1}c))$ vanish.

\begin{definition}[The deformation rings]\label{pf:eis:rings}
\leavevmode
\begin{enumerate}[label=(\alph*)]
\item $U$ is a universal ring for all lifts of $\bar r_0\circ\iota_\ell$ (no condition), with universal lift $r_U$.
\item $I_t\subset U$ is the intersection of those minimal primes of the radical of an explicit ideal $I_Y$ of $U$ which are
  contained in the kernel of every $\OO$-point $x:U\to\OO$ with $x\circ r_U=\eisrz(\eta,\Psi)\circ\iota_\ell$. Here $I_Y$ consists of
  the elements $a$ with $\delta^na\in I_{\mathrm{ord}}$ for some $n$, where $I_{\mathrm{ord}}$ is generated by the coefficients
  of the characteristic polynomials of $r_U(s)$ minus those of $\mathrm{diag}(\Psi,1,\Psi^{-1})(s)$ and by the entries of
  products of three shifted matrices $r_U(s_i)-\Psi^{w(i)}(s_i)$, for inertia elements $s,s_i$ in an open subgroup, and
  $\delta\in U$ is an explicit element (a multiple of $\ell$). It imposes an ordinarity condition at $\ell$.
\item For a marked $g_1$, $S=S(U,g_1)$ is a universal ring for the lifts $r$ of $\bar r_0$ which are normalized at
  $(c,g_1)$, trivial on $\eisKill_D$, and such that $r\circ\iota_\ell=f\circ r_U$ for some $f:U\to R$ with $f(I_t)=0$. We write
  $r_S$ for its universal lift.
\end{enumerate}
\formal{UpperTheory.V.U3.SRing}
\end{definition}

\begin{definition}[Cohomological inputs]\label{pf:eis:coh}
$(\mathrm{Coh})$ is the conjunction of the local results of the literature package (among them local Tate duality, finiteness
and Euler characteristics over finite extensions of $\Ql$), Poitou--Tate duality, global finiteness, the global Euler characteristic
formula, and $\sfname{LemmaK}(C_K)$ for some $C_K$.
\formal{UpperTheory.V.U3.CohInp}
\end{definition}

$(\mathrm{Coh})$ follows from component \ref{pf:vwin-he:H1} of $\sfname{HypothesesH'}$ and Lemma K. Under $(\mathrm{Coh})$, $U$ exists
(\lref{UpperTheory.V.U3.exists_uLoc}) and, for marked $g_1$, $S$ exists (\lref{UpperTheory.V.U3.exists_ring}).

\emph{Points.} The \emph{weight-zero subalgebra} of $S$ is the $\OO$-subalgebra generated by the products of entries of
$\mathrm{mat}(r_S(g))$ of total weight $0$. A \emph{point} of $S$ is an $\OO$-algebra homomorphism $f:S\to\Qlbar$; it is
\emph{local} if $|f(x)|<1$ for $x\in\mathfrak m_S$; it defines $r_f=f\circ r_S:G_\Q\to\eisGG(\Qlbar)$, which has multiplier
$\eta$. A \emph{base point} is an $\OO$-algebra homomorphism $x_0:S\to\OO$ with $x_0\circ r_S=\eisrz(\eta,\Psi)$. Conjugation by
$\tau(t)$, $t\in\Qlbar^\times$, acts on the representations $r_f$.

\begin{definition}[Two-signed points, trace differences, rigidity]\label{pf:eis:twopt}
Let $f$ be a local point of $S$ and $\sigma\in G_\Q$.
\begin{enumerate}[label=(\alph*)]
\item $f$ is \emph{two-signed} if $r_f$ is neither upper nor lower (\cref{pf:eis:GG}).
\item The \emph{trace difference} is $\eistr_\sigma(f)=f(\tr\mathrm{mat}(r_S(\sigma)))-\tr\mathrm{mat}(\eisrz(\eta,\Psi)(\sigma))$.
\item A \emph{first-order deformation} of $f$ is an $\OO$-algebra homomorphism $F:S\to\Qlbar[\varepsilon]$, $\varepsilon^2=0$, with
  first component $f$. Writing $\mathrm{mat}(F(r_S(g)))=\mathrm{mat}(r_f(g))+\varepsilon Y(g)$, its \emph{tangent cocycle} is
  $y_F(g)=Y(g)\,\mathrm{mat}(r_f(g))^{-1}$, a $1$-cocycle of $G_\Q$ in the adjoint representation of $r_f$. The point $f$ is
  \emph{rigid} if every tangent cocycle at $f$ is a coboundary $g\mapsto X-r_f(g)\cdot X$, and \emph{rigid in weight $j$} if
  this holds for the tangent cocycles with values in $\mathfrak g_j$.
\item For a base point $x_0$, $\mathrm{Pt}(x_0)$ is the conjunction of: every local point $f$ with $r_f$ upper and lower
  satisfies $r_f=r_{x_0}$; $x_0$ (viewed as a point) is rigid in every weight $j\le0$; every two-signed point is rigid.
\end{enumerate}
\formallabel{UpperTheory.V.U3.TwoPt}
\formallabel{UpperTheory.V.U3.trDiff}
\formallabel{UpperTheory.V.U3.PtHyp}
\end{definition}

\begin{definition}[Models]\label{pf:eis:model}
A \emph{model} $M$ of $S$ consists of a weighted ring $(A,x)$ (\cref{pf:vwin-he:wring}), integers $r\ge1$ and
$s,r',s',m\ge0$, a system of equations $Q=(F_0,\dots,F_r,G,E,E',Z)$ of type $(r,s,r',s',m)$ with ideal $I$
(\cref{pf:vwin-he:eqns}) such that $x$ is smooth with tangent dimensions $n_{r,s,r',s',m}$ (\cref{pf:vwin-he:chart}), an
$\OO$-algebra homomorphism $q:A\to S$ with kernel $I$, and a base point $x_0$ of $S$ with $x_0\circ q=x$, such that:
\begin{enumerate}[label=(\alph*)]
\item the $(a,b)$ entry of $\mathrm{mat}(r_S(g))$ is $q(f)$ for some $f\in A_{w(a,b)}$, for all $g,a,b$;
\item for $f\in A_j$ and $n\ge1$, $q(f)$ is congruent modulo $\mathfrak m_S^n$ to an $\OO$-linear combination of products of
  entries of $r_S$ of total weight $j$;
\item if condition (M2) of \cref{pf:vwin-he:eqns} holds, then for every $\Qlbar$-point $\xi$ of $A$ on $I$ there are
  $t\in\Qlbar^\times$ and a local point $f$ of $S$ with $t\cdot\xi=f\circ q$;
\item if (M2) holds, then for every local point $f$ of $S$, every point derivation of $A$ at $f\circ q$ which vanishes on $I$
  is $y\mapsto F(q(y))_\varepsilon$ for a first-order deformation $F$ of $f$;
\item a lifting property of derivations of weight $-1$ modulo non-zero elements of $\OO$, relating $S$ to $U$ and $I_t$ (used
  only through \lref{UpperTheory.V.U3.model_lower} below).
\end{enumerate}
The \emph{length} of $M$ is $\lambda(M)=\eiscot_{-1}(x;F)\in\N\cup\{\infty\}$ (\cref{pf:vwin-he:cot}).
\formal{UpperTheory.V.U3.Model}
\end{definition}

The grading of $A$ models the action of conjugation by $\tau(t)$: by (a), the torus acts on points of $A$ through the
grading and on representations through conjugation. By \lref{UpperTheory.V.U3.exists_model} ($(\mathrm{Coh})$ and Lemma P
imply that for every datum, marked $g_1$ and ring $S$ a model of $S$ over $\Qlbar$ exists), models exist; this is where
Lemma P (\lref{UpperTheory.V.Alg.lemmaP}: weight-graded presentations of framed deformation rings, with numbers of
generators and relations in each weight given by the dimensions of the graded pieces of $H^1$ and of the obstruction
space) enters.

\subsubsection{The hypotheses on automorphy and rigidity}

\begin{definition}[Unitary automorphy with unramified coefficients]\label{pf:eis:Uur}
$(\mathrm{U}_{\mathrm{ur}})$ is the statement: for all $\Omega,K,\ell,\OO,N,e,\psi$ as above with $e_\OO=1$, every datum $D$ and
every $\iota:\Qlbar\cong\C$, there is a marked $g_1\in G_\Q$ such that for all rings $U$ and $S=S(U,g_1)$:
\begin{enumerate}[label=(\alph*)]
\item the weight-zero subalgebra of $S$ is a finite $\OO$-module;
\item for every $\OO$-algebra homomorphism $\xi$ from the subalgebra of $S$ generated by the entries of $r_S$ to $\Qlbar$ such
  that $r=\xi\circ r_S$ is of type I, $\mathcal A_\iota(r_K\otimes\eiscyc^{-1})$ holds;
\item for every such $\xi$ with $r=\xi\circ r_S$ of type II and middle entry $1$ on $\ker\eta$, every
  $W:G_\Q\to\GL_2(\Qlbar)$ which equals the corner of $r$ on $\ker\eta$ and has $\tr W(c)=0$ is \emph{modular}: a newform of
  weight $2e+1$ is attached via $\iota$ to $W\otimes\eiscyc^{-e}$, and $\sfname{WitAQ}(\iota,W)$ holds.
\end{enumerate}
\formal{UpperTheory.V.Aut.AutomorphyUnramifiedCoefficientsAlt}
\end{definition}

This is the statement $(\mathrm{U}_{\mathrm{ur}})$ of \cref{pf:vwin-he:autstmts}. A continuous representation $V$ of $G_\Q$ on
a finite-dimensional $\Qlbar$-vector space has \emph{vanishing Bloch--Kato Selmer group} if every continuous $1$-cocycle
which satisfies the Bloch--Kato local conditions (unramified up to a coboundary on an inertia group at every place not above
$\ell$, and crystalline at $\ell$ in the sense of the period ring of the development) is a coboundary
\cite{BlochKato1990}.

\begin{definition}[Rigidity and admissibility statements]\label{pf:eis:rigstmts}
Let $D$ be a datum, $U$, a marked $g_1$, $S=S(U,g_1)$, $f$ a local point of $S$ and $\iota:\Qlbar\cong\C$.
\begin{enumerate}[label=(\alph*)]
\item $(\mathrm{R}_{\mathrm{I}})$: if $r_f$ is of type I and $\mathcal A_\iota((r_f)_K\otimes\eiscyc^{-1})$ holds, the Bloch--Kato Selmer group of
  the adjoint representation of $r_f$ vanishes.
\item $(\mathrm{R}_{\mathrm{II}})$: if $r_f$ is of type II and $W:G_\Q\to\GL_2(\Qlbar)$ is continuous, equals the corner of $r_f$ on
  $\ker\eta$, has $\tr W(c)=0$, $\det W=\eta$ and is modular, then every representation $R$ of $G_\Q$ on $M_2(\Qlbar)$ of
  \emph{adjoint type for $W$} has vanishing Bloch--Kato Selmer group; here $R$ is of adjoint type for $W$ if there are a
  complex conjugation $c'$ and $P\in\GL_2(\Qlbar)$ with ${}^tW(s)PW(s')=\eta(s)P$ whenever $s'=c'sc'^{-1}$ ($s\in G_K$),
  $R(s)X=W(s)XW(s)^{-1}$ for $s\in G_K$, and $R(c')X=-P^{-1}\,{}^tX\,P$.
\item $(\mathrm{A}_{\mathrm{I}})$: if $r_f$ is of type I and $\mathcal A_\iota((r_f)_K\otimes\eiscyc^{-1})$ holds, then $(r_f)_K$ is admissible of
  unitary type for $(\iota,e,N)$.
\end{enumerate}
\formallabel{UpperTheory.V.Th.RigidityFirstCase}
\formallabel{UpperTheory.V.Th.RigiditySecondCase}
\formallabel{UpperTheory.V.Z.J.AdmI_s}
\end{definition}

These three statements are theorems of the development, proved from $(\mathrm{Coh})$, $\sfname{HypothesesH'}$, an adjoint
lifting theorem and an unlinkedness statement for ordinary points (see the proof of \cref{pf:eis:main}).

\subsubsection{The upper bound}

\begin{lemma}[Upper bound]\label{pf:eis:upper}
Assume Lemma M. Let $D$ be a datum with $\ell\ne2$, $g_1$ marked, $S=S(U,g_1)$, and $M$ a model of $S$ with base point $x_0$.
Assume that the weight-zero subalgebra of $S$ is a finite $\OO$-module, that $\mathrm{Pt}(x_0)$ holds, and let $\sigma\in\ker\eta$
be such that $\eistr_\sigma(f)\ne0$ for every two-signed point $f$. Then $\lambda(M)<\infty$, and there is a finite set $T$ of
two-signed points such that every two-signed point $f$ satisfies $r_f=\tau(t)r_{f'}\tau(t)^{-1}$ for some $f'\in T$ and
$t\in\Qlbar^\times$, distinct elements of $T$ are not related in this way, and
\[
 \|\lambda(M)\|_\OO\ \le\ 2\sum_{f\in T}\eisv(\eistr_\sigma(f)).
\]
\formallabel{UpperTheory.V.U3.model_upper}
\end{lemma}

The lemma holds more generally for the universal rings of the abstract setting of the development (a group $\Gamma$ with a
subgroup mapping to it, in place of $G_\Q$ and $G_{\Ql}$); we state it in the case used.

\begin{proof}
We verify the conditions (M1)--(M5) of \cref{pf:vwin-he:eqns} for the weighted ring $(A,x)$ and the system $Q$ of $M$, with a
suitable $a\in A_0$, using the following results of the development, all for $\ell\ne2$ and marked $g_1$.
\begin{itemize}
\item \lref{UpperTheory.V.U3.modelIdeal_finiteness_of_weightZero_finite}: if the weight-zero subalgebra of $S$ is finite
  over $\OO$, then (M2) holds.
\item \lref{UpperTheory.V.U3.gradedModel_spansCot_of_rigid_weights}: (M2) and rigidity of the base point in all weights
  $j\le0$ imply (M1).
\item \lref{UpperTheory.V.U3.gradedModel_pointCondition_of_ptHyp}: (M2) and $\mathrm{Pt}(x_0)$ imply (M3).
\item \lref{UpperTheory.V.U3.gradedModel_derivationCondition_of_ptHyp}: (M2) and $\mathrm{Pt}(x_0)$ imply (M4).
\item \lref{UpperTheory.V.U3.model_signs} (signs): for local points $f$ of $S$, $f\circ q$ is negatively trivial if and only if
  $r_f$ is upper, positively trivial if and only if $r_f$ is lower, and if $r_{f'}=\tau(t)r_f\tau(t)^{-1}$ then
  $f'\circ q=t\cdot(f\circ q)$.
\end{itemize}
By the first four, (M1)--(M4) hold. By the signs, a local point $f$ is two-signed if and only if $f\circ q$ is a two-signed
$\Qlbar$-point of $A$. Conversely, if $f\circ q=t\cdot(f'\circ q)$ for local points $f,f'$, then $r_f=\tau(t)r_{f'}\tau(t)^{-1}$: by
(a) of \cref{pf:eis:model} each entry of $\mathrm{mat}(r_f(g))$ of weight $w$ is $f(q(e))$ with $e\in A_w$, and
$(t\cdot\xi)(e)=t^w\xi(e)$ for $e\in A_w$, while conjugation by $\tau(t)$ multiplies entries of weight $w$ by $t^w$.

\emph{Condition (M5).} By (a) of \cref{pf:eis:model} choose $E_{ii}\in A_0$ with $q(E_{ii})=\mathrm{mat}(r_S(\sigma))_{ii}$
($w(i,i)=0$), and put $a=E_{11}+E_{22}+E_{33}-\tr\mathrm{mat}(\eisrz(\eta,\Psi)(\sigma))\in A_0$. Then $f(q(a))=\eistr_\sigma(f)$ for
every point $f$, and $x(a)=x_0(q(a))=0$ because $x_0\circ r_S=\eisrz(\eta,\Psi)$. Let $\xi$ be a two-signed $\Qlbar$-point of $A$ on
$I$. By (c) of \cref{pf:eis:model} (which applies since (M2) holds) there are $t$ and a local point $f$ with
$t\cdot\xi=f\circ q$; as $a$ has degree $0$, $\xi(a)=(t\cdot\xi)(a)=\eistr_\sigma(f)$, and $f$ is two-signed because $t\cdot\xi$ is.
Hence $\xi(a)\ne0$, and (M5) holds.

\emph{Application of Lemma M.} The function $v=e_\OO\eisv$ on $\Qlbar$ is a normalized valuation in the sense of
\cref{sec:pf:vwin-he:lemmaM}: it is multiplicative and ultrametric on non-zero elements, and $e_\OO\eisv(c)=\ord_\OO(c)$ for
$0\ne c\in\OO$. Lemma M (\cref{pf:vwin-he:lemmaM}, with $K=\Qlbar$) gives a finite set $T_0$ of two-signed points of $A$ on
$I$ covering $I$, $\lambda(M)<\infty$, and $\lambda(M)\le2\sum_{\xi\in T}v(\xi(a))$ for every finite set $T$ covering $I$.

\emph{From points of $A$ to points of $S$.} For $\xi\in T_0$ choose by (c) an element $t_\xi$ and a local point $f_\xi$ with
$t_\xi\cdot\xi=f_\xi\circ q$, and let $T_1=\{f_\xi\}$. Each $f_\xi$ is two-signed. If $f$ is two-signed, then $f\circ q$ is a
two-signed point on $I$ ($q$ kills $I$), so $f\circ q=t\cdot\xi$ for some $\xi\in T_0$ and $t$, hence
$f\circ q=(tt_\xi^{-1})\cdot(f_\xi\circ q)$ and $r_f$ is conjugate to $r_{f_\xi}$ by an element $\tau(\cdot)$. The relation
``$r_f=\tau(t)r_{f'}\tau(t)^{-1}$ for some $t$'' is transitive, so removing successively elements of $T_1$ which are related to
another element we obtain $T\subset T_1$ with the covering property and pairwise unrelated elements.

Let $T'=\{f\circ q:f\in T\}$. It covers $I$: a two-signed point $\xi$ on $I$ satisfies $t\cdot\xi=f\circ q$ with $f$ two-signed,
$r_f=\tau(t')r_{f'}\tau(t')^{-1}$ with $f'\in T$, hence $f\circ q=t'\cdot(f'\circ q)$ by the signs, and
$\xi=(t^{-1}t')\cdot(f'\circ q)$. The map $f\mapsto f\circ q$ is injective on $T$, since $f\circ q=f'\circ q$ implies $r_f=r_{f'}$.
Lemma M for $T'$ gives $\lambda(M)\le2e_\OO\sum_{f\in T}\eisv(\eistr_\sigma(f))$; dividing by $e_\OO\ge1$ gives the assertion.
\end{proof}

\subsubsection{The lower bound}

Recall that $V_\ell(\psi,e)$ is the $\OO$-length of $\Sha_{\ell N}=\Sha_{\ell N}(\Q,\OO\otimes\eisQZ(\Psi^{-1}))$
(\cref{pf:vwin-he:V}).

\begin{lemma}[Lower bound]\label{pf:eis:lower}
Assume $(\mathrm{Coh})$ and Lemma M$'$. Let $D$ be a datum, $g_1$ marked, $S=S(U,g_1)$ and $M$ a model of $S$. Then
$V_\ell(\psi,e)\le\lambda(M)$.
\formallabel{UpperTheory.V.U3.vlen_le}
\end{lemma}

\begin{proof}
\emph{Torsion.} Every element of $\Sha_{\ell N}$ is killed by a non-zero element of $\OO$: a continuous cocycle of the compact group
$G_\Q$ with values in the discrete module $\OO\otimes\eisQZ$ takes finitely many values, each killed by a non-zero element of $\OO$,
and the product of these elements kills the cocycle. For a torsion module $X$ over the discrete valuation ring $\OO$ with
uniformizer $\varpi$, $X$ is the increasing union of the submodules $X[\varpi^k]$, so $\len_\OO X=\sup_k\len_\OO X[\varpi^k]$.
It therefore suffices to show $\len_\OO\Sha_{\ell N}[b]\le\lambda(M)$ for every $0\ne b\in\OO$.

\emph{Injection into cocycles.} For $0\ne a\in\OO$ let $\mathcal C_a$ be the set of maps $\gamma:G_\Q\to M_3(\OO/a)$ which are
locally constant, take values in $\mathfrak g_{-1}$, are $1$-cocycles for the reduction of $\eisrz(\eta,\Psi)$ modulo $a$, are
normalized at $(c,g_1)$ (that is, $\gamma(c)=0$ and the entries $(1,2),(2,1)$ of
$\gamma(g_1)\eisrz(g_1)-\gamma(cg_1^{-1}c)\eisrz(cg_1^{-1}c)$ vanish), vanish on $\eisKill_D$, and agree on $\iota_\ell(G_{\Ql})$ with
the coboundary of an element of $\mathfrak g_{-1}(\OO/a)$. By \lref{UpperTheory.V.U3.sha_inj}, for every $0\ne b\in\OO$ there are
$0\ne a\in\OO$ and an injective additive map $\gamma:\Sha_{\ell N}[b]\to\mathrm{Maps}(G_\Q,M_3(\OO/a))$ with image in $\mathcal C_a$.

\emph{Comparison with the model.} By \lref{UpperTheory.V.U3.model_lower}, if $2\in\OO^\times$, $g_1$ is marked, Lemma M$'$
holds (\lref{UpperTheory.V.Alg.M.mprime_holds}: a length bound by $\eiscot_j(x;F)$ for modules embedded into point derivations
of weight $j$ modulo $c$ vanishing on the $F_i$), the pair $(U,I_t)$ is \emph{good} in the sense of the development (a local
chart condition at the base point), all powers of $\Psi$ are continuous, and $\eisrz(\eta,\Psi)$ is trivial on $\eisKill_D$, then every
$\OO$-module $X$ with an injective additive map into $\mathrm{Maps}(G_\Q,M_3(\OO/a))$ with image in $\mathcal C_a$ has
$\len_\OO X\le\lambda(M)$. The pair $(U,I_t)$ is good by \lref{UpperTheory.V.Loc.locGood} (a consequence of the local
results contained in $(\mathrm{Coh})$), and the remaining hypotheses are elementary properties of a datum
(\lref{UpperTheory.V.U3.dat_facts}, which lists among others: $2\in\OO^\times$, continuity of $\eta$ and of the powers of
$\Psi$, and $\eisrz(\eta,\Psi)(g)=1$ for $g\in\eisKill_D$). Applied to $X=\Sha_{\ell N}[b]$ this gives the claim.
\end{proof}

\subsubsection{Rigidity at the points}

\begin{lemma}[Bloch--Kato vanishing at endoscopic points]\label{pf:eis:bktypeII}
Assume $(\mathrm{Coh})$, $\sfname{KatoTwoW}$ and $(\mathrm{R}_{\mathrm{II}})$. Let $D$ be a datum, $g_1$ marked, $S=S(U,g_1)$, $f$ a local point
with $r=r_f$ of type II, and $\iota:\Qlbar\cong\C$, and assume that every $W$ which equals the corner of $r$ on $\ker\eta$ and has
$\tr W(c)=0$ is modular. Then the Bloch--Kato Selmer group of the adjoint representation of $r$ vanishes.
\formallabel{UpperTheory.V.Z.Kato.WeilSphericalRecords.bloch_kato_selmer_vanishes_of_type_ii_modular_ext}
\end{lemma}

\begin{proof}
By \lref{UpperTheory.V.U3.isPt_push}, $r$ is a \emph{point of $D$}: it has multiplier $\eta$, continuous integral entries,
reduction $\bar r_0$, is normalized at $(c,g_1)$, trivial on $\eisKill_D$, and $r\circ\iota_\ell$ comes from a local point of $U$
killing $I_t$. By \lref{UpperTheory.V.U3.typeII_facts}, for a point of $D$ of type II the middle entry of $\mathrm{mat}(r(g))$ is
$1$ and the corner has determinant $1$ for $g\in\ker\eta$, and there is a continuous irreducible $W:G_\Q\to\GL_2(\Qlbar)$ with
$W(g)=\mathrm{corner}(\mathrm{mat}(r(g)))$ if $\eta(g)=1$, $W(g)=\mathrm{corner}(\mathrm{mat}(r(g)))\,\mathrm{diag}(-1,1)$ if $\eta(g)=-1$,
$\tr W(c)=0$ and $\det W=\eta$. In particular $W$ is modular by hypothesis.

\emph{Block form.} For $g\in\ker\eta$, $\mathrm{mat}(r(g))=\begin{psmallmatrix}a&0&b\\0&1&0\\c'&0&d\end{psmallmatrix}$ with $ad-bc'=1$,
since entries of odd weight vanish (type II). If $\eta(g)=-1$, then $r(g)=r(gc^{-1})r(c)$ with $\eta(gc^{-1})=1$ and
$\mathrm{mat}(r(c))=1$, so $\mathrm{mat}(r(g))=\mathrm{mat}(r(gc^{-1}))$ has the same form. For $x\in\eisGG$ with matrix of this form
consider the decomposition $M_3=\mathfrak a\oplus\mathfrak b_+\oplus\mathfrak b_-\oplus\mathfrak c$ given, for $X=(X_{ij})$, by
\[
 p_{\mathfrak a}(X)=\begin{psmallmatrix}X_{11}&X_{13}\\X_{31}&X_{33}\end{psmallmatrix},\quad
 p_\pm(X)=(X_{12}\pm X_{23},\,X_{32}\mp X_{21}),\quad p_{\mathfrak c}(X)=X_{22}E_{22},
\]
with sections $s_{\mathfrak a}(Y)$ (placing $Y$ in the entries $(1,1),(1,3),(3,1),(3,3)$) and
$s_\pm(v)=\begin{psmallmatrix}0&v_1&0\\\mp v_2&0&\pm v_1\\0&v_2&0\end{psmallmatrix}$, so that
$X=s_{\mathfrak a}p_{\mathfrak a}X+\tfrac12s_+p_+X+\tfrac12s_-p_-X+p_{\mathfrak c}X$. A direct computation with the formula for the
adjoint action in \cref{pf:eis:GG} shows, for $x=r(g)$:
\begin{itemize}
\item $p_{\mathfrak a}(x\cdot X)=p_{\mathfrak a}(x\cdot s_{\mathfrak a}p_{\mathfrak a}X)$ and $x\cdot s_{\mathfrak a}Y=s_{\mathfrak a}p_{\mathfrak a}(x\cdot s_{\mathfrak a}Y)$;
  so $R_{\mathfrak a}(g)Y=p_{\mathfrak a}(r(g)\cdot s_{\mathfrak a}Y)$ is a representation of $G_\Q$ on $M_2$, equal to $Y\mapsto W(g)YW(g)^{-1}$
  for $g\in G_K$ and to $Y\mapsto-J_2\,{}^tY\,J_2$ for $g=c$, where $J_2$ is the antidiagonal $2\times2$ matrix with entries $1$;
\item $p_\pm(x\cdot X)=W_\pm(g)p_\pm(X)$ and $x\cdot s_\pm(v)=s_\pm(W_\pm(g)v)$, where $W_+=W$ and $W_-=W\otimes\eta$;
\item $p_{\mathfrak c}(x\cdot X)=\eta(g)p_{\mathfrak c}(X)=\eisrz(\eta,\Psi)(g)\cdot p_{\mathfrak c}(X)$, and $x\cdot p_{\mathfrak c}(X)=\eta(g)p_{\mathfrak c}(X)$.
\end{itemize}
Moreover ${}^tW(s)J_2W(s')=\eta(s)J_2$ for $s\in G_K$ and $s'=csc^{-1}$ (as $W(c)=\mathrm{diag}(-1,1)$ and $\det W=\eta$), so $R_{\mathfrak a}$ is
of adjoint type for $W$ with $c'=c$ and $P=J_2$.

\emph{Vanishing on the pieces.} $R_{\mathfrak a}$ has vanishing Bloch--Kato Selmer group by $(\mathrm{R}_{\mathrm{II}})$. The
representations $W$ and $W\otimes\eta$ have vanishing Bloch--Kato Selmer groups by $\sfname{KatoTwoW}$: by
\lref{UpperTheory.V.Z.Kato.used_of_alt}, $W$ satisfies the configuration $\sfname{Used}$ of \cref{def:window:kato} for $K$, the
twist $e$ and the character $\psi_G$ (it is ordinary at $\ell$ with the prescribed flag, of finite level, irreducible on
$G_{K(\mu_{\ell^\infty})}$, and has the prescribed residual trace), and \lref{UpperTheory.V.Z.Kato.WeilSphericalRecords.kUse'}
deduces from $\sfname{KatoTwoW}$ and the modularity of $W$ that $W$ and $W\otimes\eta$ have vanishing Bloch--Kato Selmer groups.
Finally, a Bloch--Kato cocycle of the adjoint representation of $\eisrz(\eta,\Psi_{\Qlbar})$ with values in $\mathfrak g_0$ is a
coboundary (\lref{UpperTheory.V.U3.bloch_kato_nonpositive_weight}, which states this for values in $\mathfrak g_j$, $j\le0$).

\emph{Conclusion.} Let $\gamma$ be a Bloch--Kato cocycle of the adjoint representation of $r$. The maps $p_{\mathfrak a},p_\pm,p_{\mathfrak c}$
are continuous, linear and equivariant, so they carry $\gamma$ to cocycles of $R_{\mathfrak a}$, $W$, $W\otimes\eta$ and
$\mathrm{ad}\,\eisrz$ (values in $\mathfrak g_0$) which again satisfy the Bloch--Kato local conditions (the conditions are defined by
the existence of local coboundaries and of elements of a tensor product with the period ring, and are preserved by equivariant
linear maps). By the previous step they are coboundaries of elements $x_{\mathfrak a},x_+,x_-,x_{\mathfrak c}$. By the decomposition of
$X$ and the equivariance of the sections, $\gamma$ is the coboundary of
$s_{\mathfrak a}x_{\mathfrak a}+\tfrac12s_+x_++\tfrac12s_-x_-+p_{\mathfrak c}x_{\mathfrak c}$ ($\ell$ is odd, so $2$ is invertible).
\end{proof}

\begin{lemma}[The hypothesis $\mathrm{Pt}$]\label{pf:eis:pthyp}
Assume $(\mathrm{Coh})$, $\sfname{KatoTwoW}$, $(\mathrm{R}_{\mathrm{I}})$ and $(\mathrm{R}_{\mathrm{II}})$. Let $D$ be a datum, $\iota:\Qlbar\cong\C$, and $g_1$
marked such that (b) and (c) of \cref{pf:eis:Uur} hold for all $U$ and $S=S(U,g_1)$. Then $\mathrm{Pt}(x_0)$ holds for every
base point $x_0$ of every such $S$.
\formallabel{UpperTheory.V.Z.J.ptHyp_of_autB_of_lemW_of_isBase}
\end{lemma}

\begin{proof}
As $x_0$ is a base point, $r_{x_0}=\eisrz(\eta,\Psi_{\Qlbar})$. We use the following results of the development (for marked $g_1$):
\lref{UpperTheory.V.U3.isPt_push} (every $r_f$, $f$ local, is a point of $D$, see the proof of \cref{pf:eis:bktypeII});
\lref{UpperTheory.V.U3.pt_fixed} (a point of $D$ which is upper and lower equals $\eisrz(\eta,\Psi_{\Qlbar})$);
\lref{UpperTheory.V.U3.tanCoc_isBK} (for every first-order deformation $F$ of a local point $f$, the tangent cocycle $y_F$ is a
continuous cocycle of the adjoint representation of $r_f$ satisfying the Bloch--Kato local conditions); and
\lref{UpperTheory.V.U3.twoPt_shape} (for $\ell\ne2$, $r_f$ is of type I or of type II at every two-signed point $f$).

\emph{Fixed points.} If $r_f$ is upper and lower, then $r_f=\eisrz(\eta,\Psi_{\Qlbar})=r_{x_0}$ by the first two results.

\emph{Weights $j\le0$.} If $F$ is a first-order deformation of $x_0$ whose tangent cocycle takes values in $\mathfrak g_j$, $j\le0$,
then $y_F$ is a Bloch--Kato cocycle of the adjoint representation of $\eisrz(\eta,\Psi_{\Qlbar})$ with values in $\mathfrak g_j$, hence a
coboundary by \lref{UpperTheory.V.U3.bloch_kato_nonpositive_weight}.

\emph{Two-signed points.} Let $f$ be two-signed and $F$ a first-order deformation of $f$. Since $y_F$ is a Bloch--Kato cocycle, it
suffices that the Bloch--Kato Selmer group of the adjoint representation of $r_f$ vanishes. If $r_f$ is of type I, then
$\mathcal A_\iota((r_f)_K\otimes\eiscyc^{-1})$ holds by (b) of \cref{pf:eis:Uur} (applied to the restriction of $f$ to the subalgebra
generated by the entries), and $(\mathrm{R}_{\mathrm{I}})$ applies. If $r_f$ is of type II, its middle entry is $1$ on $\ker\eta$
(\lref{UpperTheory.V.U3.typeII_facts}), so every $W$ extending its corner with $\tr W(c)=0$ is modular by (c) of
\cref{pf:eis:Uur}, and \cref{pf:eis:bktypeII} applies.
\end{proof}

\subsubsection{Admissibility of the points}

\begin{lemma}[Type II points are endoscopic]\label{pf:eis:admII}
Assume $(\mathrm{Coh})$. Let $D$ be a datum, $g_1$ marked, $r$ a point of $D$ of type II and $\iota:\Qlbar\cong\C$, and assume
that every $W$ which equals the corner of $r$ on $\ker\eta$ and has $\tr W(c)=0$ is modular. Then $r_K$ is endoscopic for
$(\iota,e,N|d_K|)$ (\cref{pf:eis:adm}).
\formallabel{UpperTheory.GalN.V.adm2B_rhoK_of_typeII_of_modular_extensions}
\end{lemma}

\begin{proof}
Let $W$ be as in the proof of \cref{pf:eis:bktypeII} (\lref{UpperTheory.V.U3.typeII_facts}): continuous, irreducible, equal
to the corner of $r$ on $\ker\eta$, with $\tr W(c)=0$ and $\det W=\eta$. It is modular by hypothesis. On $G_K=\ker\eta$, $r_K$ is the
block embedding of $W|_{G_K}$, which is the last condition of \cref{pf:eis:adm}(b) (with conjugating matrix $1$). Since
$\gcd(|d_K|,N)=1$ and $N$ is even, $|d_K|$ is odd; an element fixing the $N|d_K|$-th roots of unity fixes $K$ (a quadratic field of odd
discriminant $d_K$ is contained in $\Q(\mu_{|d_K|})$; \lref{UpperTheory.mem_galOf_quadratic_of_fixes_discr_roots_of_unity}) and lies
in the kernel of $\psi_G$. Hence $\det W=\eta$ has conductor
dividing $N|d_K|$, so $\det W$ is a character of finite order as required. It remains to show that $W$ has level $N|d_K|$ and
type $(e,-e)$, and we write $\Psi=\eiscyc^e\psi_G$ (in $\Qlbar$).

\emph{Places $w\nmid\ell$.} Let $A$ be a place not above $\ell$. By \lref{UpperTheory.V.U3.pt_inertia_away} there is
$M\in\GL_3(\Qlbar)$ with $M\,\mathrm{mat}(r(g))=\mathrm{mat}(\eisrz(\eta,\Psi)(g))\,M$ for $g\in I_A\cap\ker\eta$. On this subgroup
$\mathrm{mat}(r(g))$ is the block embedding of $W(g)$ and $\eisrz(\eta,\Psi)(g)=\mathrm{diag}(\psi_G(g),1,\psi_G(g)^{-1})$, since $\eiscyc$
is trivial on $I_A$. By \lref{UpperTheory.diagonal_on_subgroup_of_conjugate} (if the block embedding of a $2\times2$
representation is conjugate to $\mathrm{diag}(\Psi',1,\Psi'^{-1})$ on a subgroup, the representation is diagonalizable there with
characters $\Psi'^{\pm1}$), $W$ is diagonalizable on $I_A\cap\ker\eta$ with characters $\psi_G^{\pm1}$, and by
\lref{UpperTheory.exists_isDiagOn_of_isDiagOn_inf_ker_sign} $W$ is diagonalizable on $I_A$ through characters $\chi_1,\chi_2$
which are trivial at every $g$ with $\psi_G(g)=1$ and $\eta(g)=1$, in particular on the elements fixing the $N|d_K|$-th roots of
unity.

\emph{Places $w\mid\ell$.} Let $A$ be a place above $\ell$. As $\ell\nmid d_K$, $I_A\subset\ker\eta$, and by the list of properties
of a datum (\lref{UpperTheory.V.U3.dat_facts}) there is $s_0\in I_A$ with $\Psi(s_0)^2-1\in\OO^\times$. By
\lref{UpperTheory.V.U3.pt_inertia_above} there is $M\in\GL_3(\Qlbar)$ such that $M\,\mathrm{mat}(r(s))\,M^{-1}$ has vanishing entries of
negative weight for $s\in D_A\cap\ker\eta$ and diagonal entries $\Psi(s),1,\Psi(s)^{-1}$ for $s\in I_A$. Hence $v=M^{-1}e_1$ is an
eigenvector of $\mathrm{mat}(r(s))$ with eigenvalue $\Psi(s)$ for $s\in I_A$. The middle coordinate of $v$ vanishes, since the middle
entry of $\mathrm{mat}(r(s_0))$ is $1$ and $\Psi(s_0)\ne1$; so $(v_1,v_3)\ne0$ is an eigenvector of $W|_{I_A}$ with eigenvalue $\Psi$.
As $I_A$ is normal in $D_A$ and $\det W=\eta=1$ on $I_A$, \lref{UpperTheory.exists_isFlagOn_of_eigenvector_normal_subgroup}
gives a flag of $W|_{D_A}$ with diagonal characters $\kappa_1,\kappa_2$ equal to $\Psi=\eiscyc^e\psi_G$ and $\Psi^{-1}=\eiscyc^{-e}\psi_G^{-1}$
on $I_A$. The flag is unlinked: by \lref{UpperTheory.V.U3.dat_facts} there is $s\in I_A$ such that $\eiscyc(s)^n=1$ only for
$n=0$; if $\kappa_i=\kappa_j\eiscyc$, evaluating at $s$ and raising to the power $n_0=\#(\Z/N)^\times$, which kills $\psi_G$, gives
$\eiscyc(s)^{n_0(m_i-m_j-1)}=1$ with $(m_1,m_2)=(e,-e)$, hence $m_i-m_j=1$, which is impossible since $m_i-m_j\in\{0,\pm2e\}$.
\end{proof}

\subsubsection{Proof of the one-character bound}

\begin{theorem}[One-character bound]\label{pf:eis:onecharthm}
$(\mathrm{Coh})$, Lemma M, Lemma M$'$, Lemma P, $\sfname{KatoTwoW}$, $(\mathrm{R}_{\mathrm{I}})$, $(\mathrm{R}_{\mathrm{II}})$, $(\mathrm{A}_{\mathrm{I}})$ and
$(\mathrm{U}_{\mathrm{ur}})$ imply $(\mathrm{One})$.
\formallabel{UpperTheory.V.Z.J.oneCharU_of_katoTwoW_rig_admI_automorphy}
\end{theorem}

\begin{proof}
Let $K,\ell,\OO,N,e,\psi,\iota,\sigma,t$ be as in \cref{pf:eis:onecharstmt}.

\emph{Reduction to $e_\OO=1$.} By (G2), $n=\ord\psi$ is prime to $\ell$. There is a coefficient ring $\OO'$ with $e_{\OO'}=1$
containing the $n$-th roots of unity (\lref{UpperTheory.V.HV.coefSup}, which states among other things that for $\ell\nmid n$
such an unramified coefficient ring exists); hence there is a character $\psi'$ modulo $N$ with values in $\OO'$ and the same values
as $\psi$ in $\Qlbar$. $\mathrm{GD}(K,\ell,e,\psi')$ holds: the conditions (G1)--(G6) only involve orders, kernels and the values $\pm1$,
which are preserved by injective ring homomorphisms. By \lref{UpperTheory.V.HV.coefIndep} (if two characters with values in
coefficient rings have the same values in $\Qlbar$, then $V_\ell$ is infinite for both or for neither, and $\eisVt_\ell$ agrees),
$V_\ell(\psi,e)$, $V_\ell(\psi',e)$ are finite together and $\eisVt_\ell(\psi,e)=\eisVt_\ell(\psi',e)$. Since $\Psi_{\psi,e}$ and
$\Psi_{\psi',e}$ agree in $\Qlbar$, $t$ and the hypothesis on admissible $\rho$ are the same. So we may assume $e_\OO=1$.

\emph{The rings.} Choose a datum $D$ (\lref{UpperTheory.V.U3.exists_dat}). By $(\mathrm{U}_{\mathrm{ur}})$ there is a marked
$g_1$ with (a)--(c) of \cref{pf:eis:Uur}. Choose $U$, $S=S(U,g_1)$ and a model $M$ of $S$ with base point $x_0$
(\lref{UpperTheory.V.U3.exists_uLoc}, \lref{UpperTheory.V.U3.exists_ring}, and \lref{UpperTheory.V.U3.exists_model}). For a point $f$ of
$S$ let $\rho_f=(r_f)_K:G_K\to\GL_3(\Qlbar)$.

\emph{Traces.} Since $\sigma\in G_K=\ker\eta$, $\eisrz(\eta,\Psi)(\sigma)=\mathrm{diag}(\Psi(\sigma),1,\Psi(\sigma)^{-1})$ and
$\eistr_\sigma(f)=\tr\rho_f(\sigma)-t$ for every point $f$.

\emph{Admissibility.} Let $f$ be two-signed. Then $r_f$ is a point of $D$ (\lref{UpperTheory.V.U3.isPt_push}) of type I or II
(\lref{UpperTheory.V.U3.twoPt_shape}). If it is of type I, (b) of \cref{pf:eis:Uur} and $(\mathrm{A}_{\mathrm{I}})$ show that $\rho_f$ is
admissible of unitary type. If it is of type II, its middle entry is $1$ on $\ker\eta$ (\lref{UpperTheory.V.U3.typeII_facts}),
(c) of \cref{pf:eis:Uur} gives modularity of the extensions of its corner, and \cref{pf:eis:admII} shows that $\rho_f$ is endoscopic.
In both cases $\rho_f$ is admissible, so $\eistr_\sigma(f)\ne0$ by hypothesis.

\emph{The bounds.} By \cref{pf:eis:pthyp}, $\mathrm{Pt}(x_0)$ holds; the weight-zero subalgebra of $S$ is finite by (a); $\ell\ne2$ by
(G1). \Cref{pf:eis:upper} gives $\lambda(M)<\infty$ and a finite set $T$ of two-signed points, pairwise not conjugate under the torus,
covering all two-signed points, with $\|\lambda(M)\|_\OO\le2\sum_{f\in T}\eisv(\eistr_\sigma(f))$. \Cref{pf:eis:lower} gives
$V_\ell(\psi,e)\le\lambda(M)<\infty$, and as $\|\cdot\|_\OO$ is monotone on finite values,
$\eisVt_\ell(\psi,e)\le\|\lambda(M)\|_\OO$.

\emph{The set $X$.} Let $X=\{\rho_f:f\in T\}$; its elements are admissible. If $\rho_f$ and $\rho_{f'}$ are isomorphic for $f,f'\in T$,
then by \lref{UpperTheory.V.U3.twoPt_shape} (which also states that a local point $f'$ whose $r_{f'}$ is isomorphic to $r_f$ on
$\ker\eta$, for two-signed $f$, satisfies $r_{f'}=\tau(t)r_f\tau(t)^{-1}$ for some $t$) $f$ and $f'$ are related, hence $f=f'$. So the
elements of $X$ are pairwise non-isomorphic, $f\mapsto\rho_f$ is injective on $T$, and
\[
 \eisVt_\ell(\psi,e)\le2\sum_{f\in T}\eisv(\eistr_\sigma(f))=2\sum_{\rho\in X}\eisv(\tr\rho(\sigma)-t).\qedhere
\]
\end{proof}

\subsubsection{Proof of \texorpdfstring{\cref{pf:eis:main}}{Theorem CountI}}

\begin{proof}[Proof of \cref{pf:eis:main}]
Assume $\sfname{HypothesesH'}$, $\sfname{LemmaK}(C_K)$, $\sfname{KatoTwoW}$ and $(\mathrm{U}_{\mathrm{ur}})$. The local results contained in
component \ref{pf:vwin-he:H1} of $\sfname{HypothesesH'}$ (\lref{UpperTheory.V.HV.U.locPrints_of_litN}), its global cohomological
results and $\sfname{LemmaK}(C_K)$ give $(\mathrm{Coh})$. We use the following results proved elsewhere:
\begin{itemize}
\item \lref{UpperTheory.GalN.V.Win.N.numberU3U_of_prints}: $(\mathrm{Num})$ holds;
\item \lref{UpperTheory.V.HV.DescentVanishingReductions.adjoint_lifting_of_cohomology}: $\sfname{HypothesesH'}$ implies an
  adjoint lifting theorem (vanishing of adjoint Selmer groups at locally generic automorphic points of universal ordinary
  deformation rings of rank $2\le n\le4$ over CM fields);
\item \lref{UpperTheory.V.Z.J.linkedUnlE_alt} and \lref{UpperTheory.V.Z.J.unlEnd3E}: $\sfname{HypothesesH'}$ implies a statement
  $(\mathrm{Unl})$: ordinary points of the linked unitary deformation problems, with values in a discrete valuation ring
  finite over $\OO$, which are irreducible on every open subgroup of $\ker\eta$ and have level $N'$ and type $(e,0,-e)$ at $\ell$,
  are unlinked of level $N'$ and type $(e,0,-e)$;
\item Lemma M (\cref{pf:vwin-he:lemmaM}), Lemma M$'$ (\lref{UpperTheory.V.Alg.M.mprime_holds}) and Lemma P
  (\lref{UpperTheory.V.Alg.lemmaP}).
\end{itemize}
From these, $(\mathrm{R}_{\mathrm{I}})$ follows by \lref{UpperTheory.V.Z.J.rigIUE_alt} ($(\mathrm{Coh})$, $\sfname{HypothesesH'}$, the
adjoint lifting theorem and $(\mathrm{Unl})$ imply $(\mathrm{R}_{\mathrm{I}})$), $(\mathrm{R}_{\mathrm{II}})$ by
\lref{UpperTheory.V.Z.Win.RR.rigII_a_alt} ($(\mathrm{Coh})$, $\sfname{HypothesesH'}$ and the adjoint lifting theorem imply
$(\mathrm{R}_{\mathrm{II}})$), and $(\mathrm{A}_{\mathrm{I}})$ by \lref{UpperTheory.V.Z.J.admissible_typeI_even_level_alt} ($(\mathrm{Coh})$,
$\sfname{HypothesesH'}$ and $(\mathrm{Unl})$ imply $(\mathrm{A}_{\mathrm{I}})$). \Cref{pf:eis:onecharthm} gives $(\mathrm{One})$, and
\cref{pf:eis:propGKpf,pf:eis:corGK,pf:eis:countIcor} give $\sfname{PropGK}$, $\sfname{PropGK}^{\Sigma}$ and $\sfname{CountI}$.
\end{proof}

\begin{remark}\label{pf:eis:countIrem}
(1) The hypothesis $(\mathrm{U}_{\mathrm{ur}})$ is only applied to coefficient rings with $e_\OO=1$; the reduction to this case in
\cref{pf:eis:onecharthm} uses that $\ell\nmid\ord\psi$.
(2) The rigidity statements $(\mathrm{R}_{\mathrm{I}})$, $(\mathrm{R}_{\mathrm{II}})$ and the admissibility statement $(\mathrm{A}_{\mathrm{I}})$ reduce, in the formal proof,
to the adjoint lifting theorem and to $(\mathrm{Unl})$ by placing the local point on an integral model (a discrete valuation
ring finite over $\OO$) and by the ordinarity of points of $S$ at $\ell$.
(3) $\sfname{KatoTwoW}$ enters only through \cref{pf:eis:bktypeII}, at the type II (endoscopic) points, for the two pieces $W$ and
$W\otimes\eta$ of the adjoint representation, in accordance with the description in \cref{sec:window:eis}.
\end{remark}

% ---------------------------------------------------------------------------
\subsection{Case V: host frames over \texorpdfstring{$\Q(i)$}{Q(i)}}\label{sec:pf:eis:countV}

\begin{theorem}\label{pf:eis:countV}
Assume $\sfname{HypothesesH'}$, $\sfname{LemmaK}(C_K)$ for some $C_K\in\R$, and the unitary automorphy statement over $\Q(i)$,
$(\mathrm{U}_i)$ (\cref{pf:eis:Ui}). Then $\sfname{CountV}$ holds.
\formallabel{UpperTheory.V.Z.J.count_v_of_printed_bound_k_reduced_unitary_lift}
\end{theorem}

Case V concerns $\ell\in B_\chi$ and $j=1$, under the $2$-adic condition $R_1(h,\ell,\theta)\le2^{c_2(\theta)}$ of
\cref{pf:vwin-he:countdef}. The proof follows the pattern of \cref{sec:pf:eis:onechar}, with $K=\Q(i)$, a deformation problem for
$G_{\Q,S}$ (a \emph{host frame}) in place of the ring $S$, and an abstract counting lemma (\cref{pf:eis:abscount}) which combines the
upper and lower bounds. Throughout, $\eta_4:G_\Q\to\{\pm1\}$ is the quadratic character of $\Q(i)$, whose kernel is $G_{\Q(i)}$.

\subsubsection{Frames}

\begin{definition}[Host data and host frames]\label{pf:eis:frame}
Let $\ell$ be a prime, $\OO$ a coefficient ring, $N\ge1$, $e\in\N$, $\psi$ a Dirichlet character modulo $N$ with values in $\OO$, and
$o_1\in\N$.
\begin{enumerate}[label=(\alph*)]
\item $\mathrm{H}(\ell,e,\psi,o_1)$ holds if $\ell\ne2$, $e\ge1$, $N$ is even, $\psi(-1)=(-1)^e$, $o_1\ge1$, $(\psi_{\ell'})_{(r)}^{2o_1}=1$
  for every odd prime $r$, $8o_1v^2<2^{c_2(\psi)}$ for every $v\in T_b(\ell,\psi)$, and $c_2(\psi)\ge4$ (notation of \cref{pf:vwin-he:cases}).
  Here $T_b(\ell,\psi)=\{\ell\}\cup S_b(\ell,\psi)$, where $S_b(\ell,\psi)$ is the set of primes $p\mid N$, $p\ne\ell$, with
  $(\psi^{(\ell)})_{(p)}\ne1$.
\item A \emph{frame} $F$ over $\OO$ indexed by a finite set $S$ consists of a homomorphism $\pi:G_\Q\to\Gamma$ to a profinite group,
  homomorphisms $\eta:\Gamma\to\{\pm1\}$ and $\Psi:\Gamma\to\OO^\times$, an element $c\in\Gamma$, subgroups $I_v\subset D_v$ of $G_\Q$ for
  $v\in S$, an element $v_\ell\in S$ and a subset $S_b\subset S$; put $T_b=\{v_\ell\}\cup S_b$. For $x\in G_\Q$ we write $\eta(x)$,
  $\Psi(x)$ for $\eta(\pi(x))$, $\Psi(\pi(x))$.
\item Assume $\mathrm{H}(\ell,e,\psi,o_1)$. A \emph{host frame} is a frame with $S$ the set of primes dividing $\ell N$,
  $\Gamma=G_{\Q,S}$ the Galois group of the maximal extension of $\Q$ in $\Omega$ unramified outside $S$ (and $\infty$), $\pi$ the
  projection, $\eta\circ\pi=\eta_4$, $\Psi\circ\pi=\Psi_{\psi,e}$, $c$ the image of a complex conjugation, $D_v=D_{A_v}$ and $I_v=I_{A_v}$ for
  chosen places $A_v$ of $\Omega$ over $v\in S$, $v_\ell=\ell$, and $S_b$ the set of $v\in S$ above primes of $S_b(\ell,\psi)$.
\end{enumerate}
\formal{UpperTheory.V.UnitaryThreeInertia.HostData}
\formal{UpperTheory.V.UnitaryThreeInertia.Frame}
\formal{UpperTheory.V.UnitaryThreeInertia.IsHostFrame}
\end{definition}

By \lref{UpperTheory.V.UnitaryThreeInertia.exists_hostFrame}, host frames exist whenever $\mathrm{H}$ holds. For a
frame, lifts of $\bar r_0=\eisrz(\eta,\bar\Psi)$ are defined as in \cref{sec:pf:eis:onechar} (with $\Gamma$ in place of $G_\Q$); here
complete $\OO$-algebras are complete Noetherian local $\OO$-algebras with residue field $k$, and for $s\in D_v$ we write $r(s)$, $\eta(s)$,
$\Psi(s)$ for $r(\pi(s))$, $\eta(\pi(s))$, $\Psi(\pi(s))$. Put $N_v=I_v\cap\ker\eta\cap\ker\bar\Psi$ (the inertia elements acting trivially
through $\bar r_0$). A lift $r$ is a \emph{point} of $F$ (normalized at $(c,g_1)$) if it is
normalized at $(c,g_1)$ and satisfies the local conditions
\begin{itemize}
\item[$(\mathrm P_\ell)$] there is $h\in\GL_3(R)$, $h\equiv1$ modulo $\mathfrak m_R$, such that $h^{-1}\mathrm{mat}(r(s))h$ is upper triangular for
  $s\in D_{v_\ell}\cap\ker\eta$ and has diagonal entries $\Psi(s),1,\Psi(s)^{-1}$ for $s\in I_{v_\ell}\cap\ker\eta$;
\item[$(\mathrm P_b)$] for $v\in S_b$, the same condition at $v$ with ``diagonal'' in place of ``upper triangular'';
\item[$(\mathrm P_0)$] for $v\notin T_b$, $r(s)=1$ for $s\in N_v$.
\end{itemize}
The element $g_1\in\Gamma$ is
\emph{marked} if $\eta(g_1)=1$ and $\Psi(g_1)^2-1\in\OO^\times$. A \emph{universal point} $(R,u)$ is a point $u:\Gamma\to\eisGG(R)$ such
that every point to a complete $\OO$-algebra $R'$ is $f\circ u$ for a unique $f:R\to R'$. For $j\in\Z$ the \emph{weight-$j$ part} $R_j$ of
$R$ is the set of $x\in R$ with $f'(x)=t^jf(x)$ whenever $f'\circ u=\tau(t)(f\circ u)\tau(t)^{-1}$; let $R_\bullet$ be the
$\OO$-subalgebra generated by the $R_j$. For an $\OO$-algebra homomorphism $\xi:R_\bullet\to\Qlbar$, $r_\xi:\Gamma\to\eisGG(\Qlbar)$ denotes
the specialization of $u$. A homomorphism $r:\Gamma\to\eisGG(\Qlbar)$ is a \emph{characteristic-zero point} if, up to conjugation by
some $\tau(t)$, it is a point of $F$ with values in a discrete valuation ring $A\subset\Qlbar$ finite over $\OO$; it is \emph{rigid} if
every continuous cocycle of the adjoint representation of $r$ which defines a first-order deformation satisfying the local
conditions is a coboundary.

\begin{definition}[Conditions on frames]\label{pf:eis:framecond}
Let $\bar{\mathfrak g}_j$ be the representation of $\Gamma$ on $\mathfrak g_j(k)$ through $\bar r_0$ (the adjoint action twisted as in
\cref{pf:eis:GG}). For $v\in S$ let $L_v(j)$ be the space of continuous $1$-cocycles $y:D_v\to\bar{\mathfrak g}_j$ such that $y$ vanishes on
$N_v$ if $v\ne v_\ell$, $y$ is a coboundary if $v=v_\ell$ and $j<0$, and $y$ vanishes on $I_{v_\ell}$ if $v=v_\ell$ and $j=0$ (no condition if
$v=v_\ell$ and $j>0$). The \emph{Selmer group} $H^1_{\mathrm{Sel}}(\bar{\mathfrak g}_j)$ is the space of continuous $1$-cocycles of $\Gamma$ whose
restrictions to all $D_v$ lie in $L_v(j)$, modulo coboundaries. A \emph{framed $2$-cocycle} is a pair $(\kappa,(\alpha_v)_{v\in S})$ of a
continuous $2$-cocycle $\kappa$ of $\Gamma$ with values in $\bar{\mathfrak g}_j$ and continuous maps $\alpha_v:D_v\to\bar{\mathfrak g}_j$ with
$d\alpha_v=\kappa|_{D_v}$; the \emph{framed coboundaries} are the sums of pairs $(d\beta,(\beta|_{D_v})_v)$, $\beta:\Gamma\to\bar{\mathfrak g}_j$
continuous, and $(0,(z_v)_v)$ with $z_v\in L_v(j)$; and $\mathrm{Ob}_j$ is the quotient. The frame $F$ satisfies
\begin{itemize}
\item $(\mathrm{F1})$ if a list of conditions holds: $2\in\OO^\times$; $\pi$, $\eta$ and $\Psi$ are continuous; $c^2=1$, $\eta(c)=-1$, $\Psi(c)=1$;
  $v_\ell\notin S_b$; the $D_v$ and $I_v$ are closed with $I_v$ normal in $D_v$, and there is $\varphi_v\in D_v$ which together with $I_v$
  topologically generates $D_v$ and, for every $n\ge1$, is sent to the class of $1$ by some continuous homomorphism $D_v\to\Z/n\Z$
  trivial on $I_v$; $\eta$ is trivial on $I_v$ and some $x\in D_v$ has $\eta(x)=1$, $\Psi(x)^2-1\in\OO^\times$ for $v\in T_b$; and $\Psi$
  is trivial on $N_v$ for $v\notin T_b$;
\item $(\mathrm{F2})$ if $H^2(D_v,\bar{\mathfrak g}_j)=0$ for $v\in T_b$ and $H^1(D_v,\bar{\mathfrak g}_j)=0$ for $v\in S_b$, for all $j\ne0$;
\item $(\mathrm{F3})$ if for every $j$ both $H^1_{\mathrm{Sel}}(\bar{\mathfrak g}_j)$ and $\mathrm{Ob}_j$ are finite-dimensional and
  $\dim H^1_{\mathrm{Sel}}(\bar{\mathfrak g}_j)-\dim\mathrm{Ob}_j=\epsilon_j$, where $\epsilon_{-1}=-1$, $\epsilon_1=\epsilon_2=1$ and $\epsilon_j=0$
  otherwise;
\item $(\mathrm{F4})$ if for every $j\le0$, every continuous cocycle of $\Gamma$ with values in $\mathfrak g_j(\OO)$ for $\eisrz(\eta,\Psi)$ which
  vanishes on $N_v$ ($v\notin T_b$), on $I_v$ ($v\in T_b$) if $j=0$, and, if $j<0$, agrees with a coboundary on $D_{v_\ell}\cap\ker\eta$,
  becomes a coboundary after multiplication by some $0\ne a\in\OO$;
\item $(\mathrm{F5})$ if every locally constant homomorphism from $\ker\eta$ to an abelian group which is trivial on the conjugates of the
  inertia groups $I_v$ lying in $\ker\eta$ is trivial;
\item $(\mathrm{F6})$ if $H^1_{\mathrm{Sel}}(\bar{\mathfrak g}_{-2})=0$.
\end{itemize}
For $0\ne a\in\OO$, $\Sha_F(a)\subset H^1(\Gamma,\OO/a)$ (with $\Gamma$ acting through $\Psi^{-1}$) is the $\OO$-submodule spanned by the classes
of continuous cocycles which vanish on $N_v$ for $v\notin T_b$ and are coboundaries on $D_v$ for $v\in T_b$, and
$\lambda_F(a)=\len_\OO\Sha_F(a)\in\N\cup\{\infty\}$.
\formal{UpperTheory.V.UnitaryThreeInertia.Frame.Ok}
\formal{UpperTheory.V.UnitaryThreeInertia.Frame.LocVan}
\formal{UpperTheory.V.UnitaryThreeInertia.Frame.Count}
\formal{UpperTheory.V.UnitaryThreeInertia.Frame.VanE}
\formal{UpperTheory.V.UnitaryThreeInertia.Frame.Cl1}
\formal{UpperTheory.V.UnitaryThreeInertia.Frame.Xpp}
\formal{UpperTheory.V.UnitaryThreeInertia.Frame.shaLen}
\end{definition}

\begin{definition}[The automorphy, rigidity and number statements over $\Q(i)$]\label{pf:eis:Ui}
For a homomorphism $r:\Gamma\to\eisGG(\Qlbar)$ of a host frame and $\iota:\Qlbar\cong\C$, let $\mathrm{B}'_\iota(r)$ mean: $\mathrm{mat}\circ r\circ\pi$
is irreducible on every open subgroup of $G_{\Q(i)}=\ker\eta_4$ and on $\ker\eta_4\cap\ker\eiscyc$, and its restriction $\rho$ to
$G_{\Q(i)}$ satisfies $\mathcal A_\iota(\rho\otimes\eiscyc^{-1})$.
\begin{enumerate}[label=(\alph*)]
\item $(\mathrm{U}_i)$ is the statement: for $K=\Q(i)$, all $\Omega,\ell,\iota,\OO,N,e,\psi,o_1$ with $\mathrm{H}(\ell,e,\psi,o_1)$ and every host
  frame $F$, there is a marked $g_1$ such that for every universal point $(R,u)$: $R_0$ is a finite $\OO$-module, and for every
  $\xi:R_\bullet\to\Qlbar$ with $r_\xi$ of type I, the restriction of $\mathrm{mat}\circ r_\xi\circ\pi$ to $G_{\Q(i)}$, twisted by $\eiscyc^{-1}$,
  satisfies $\mathcal A_\iota$.
\item $(\mathrm{U}_i')$ is the same statement with the conclusion $\mathrm{B}'_\iota(r_\xi)$ for $r_\xi$ of type I.
\item $(\mathrm{R}_i)$: for $K=\Q(i)$ and all data as in (a), every $\iota$ and marked $g_1$, every characteristic-zero point $r$ with
  $\mathrm{B}'_\iota(r)$ is rigid.
\item $(\mathrm{Num}_i)$: for $K=\Q(i)$ there is $C\in\N$ such that for all data as in (a), every $\iota$ and marked $g_1$, there is
  $\varphi\in\Gamma$ such that $\tr r(\varphi)\ne\tr\eisrz(\varphi)$ for every characteristic-zero point $r$ of type I with $\mathrm{B}'_\iota(r)$, and
  for every finite family $r_1,\dots,r_n$ of such points, pairwise non-isomorphic on $\ker\eta$,
  \[\textstyle\Bigl(\sum_i\eisv\bigl(\tr r_i(\varphi)-\tr\eisrz(\varphi)\bigr)\Bigr)\log\ell\le(\ell Ne)^C.\]
\item $(\mathrm{F4})_{\mathrm h}$: for all data as in (a), every host frame satisfies $(\mathrm{F4})$. $(\mathrm{F6})_{\mathrm h}$: for all data
  as in (a) and every $q\in\N$ with $\mathrm{X}(\ell,\psi,q)$, every host frame satisfies $(\mathrm{F6})$; here $\mathrm{X}(\ell,\psi,q)$ means that
  $q\ge1$, the conductor of $(\psi^\sharp)^2$ divides $q$, and $2\ell q\,4^{n_0(\ell,\psi)}\le2^{c_2(\psi)}$ (notation of \cref{pf:vwin-he:cases}).
\end{enumerate}
Here $\eisrz=\eisrz(\eta,\Psi)$.
\formal{UpperTheory.V.Aut.BPrime_s}
\formal{UpperTheory.V.Aut.ThmS6Core_s}
\formallabel{UpperTheory.V.Aut.ThmS6_s}
\formallabel{UpperTheory.V.UnitaryThreeInertia.HostRigid_s}
\formallabel{UpperTheory.V.UnitaryThreeInertia.HostNumber_s}
\end{definition}

Statement (a) is the statement $(\mathrm{U}_i)$ of \cref{pf:vwin-he:autstmts}; (d) is the Eisenstein congruence number of case V.

\subsubsection{Reductions}

\begin{lemma}\label{pf:eis:s6}
$(\mathrm{U}_i)$ implies $(\mathrm{U}_i')$.
\formallabel{UpperTheory.V.UnitaryThreeInertia.unitary_lift_of_reduced_core}
\end{lemma}

\begin{proof}
Given the data, take $g_1$ from $(\mathrm{U}_i)$; it remains to show $\mathrm{B}'_\iota(r_\xi)$ for $r_\xi$ of type I. The automorphy part is given by
$(\mathrm{U}_i)$. Host frames satisfy $(\mathrm{F1})$ (\lref{UpperTheory.V.UnitaryThreeInertia.hostFrame_ok}). As $R_0$ is finite over $\OO$,
$r_\xi$ is a characteristic-zero point (\lref{UpperTheory.V.UnitaryThreeInertia.isCharZeroPt_of_isPtHom}: under $(\mathrm{F1})$, for marked
$g_1$, a universal point with $R_0$ finite over $\OO$, every specialization $r_\xi$ is a characteristic-zero point). For characteristic-zero
points of type I of a host frame, irreducibility on every open subgroup of $\ker\eta_4$ and on $\ker\eta_4\cap\ker\eiscyc$ hold by
\lref{UpperTheory.V.UnitaryThreeInertia.matIrr_open} and \lref{UpperTheory.V.UnitaryThreeInertia.matIrr_cyc}.
\end{proof}

\begin{lemma}\label{pf:eis:hostdata}
Let $h=(\chi,m,q_1,d)$ be a host, $\OO$ a coefficient ring for a prime $\ell\ne2$, $e\ge1$ and $\theta\in\eisX_h(\OO,e)$ with
$R_1(h,\ell,\theta)\le2^{c_2(\theta)}$. Then $\mathrm{H}(\ell,e,\theta,o_\chi)$ and $\mathrm{X}(\ell,\theta,N_\chi)$ hold.
\formallabel{UpperTheory.V.UnitaryThreeInertia.hostData_of_host}
\end{lemma}

\begin{proof}
$\ell\ne2$, $e\ge1$ and $\theta(-1)=(-1)^e$ hold by hypothesis; $N_h$ is even and $o_\chi\ge1$ by
\lref{UpperTheory.V.HV.hostIndex}. The order bound on the odd components and the conductor condition of $\mathrm{X}$ are part
of \lref{UpperTheory.V.HV.hostTh1} (properties of the components of characters $\theta$ trivial on $H_h$). By
\lref{UpperTheory.V.HV.cmDatum_char_two_adic_order_bounds}, $R_1(h,\ell,\theta)\le2^{c_2(\theta)}$ implies $c_2(\theta)\ge4$ and
$8o_\chi v^2<2^{c_2(\theta)}$ for $v\in T_b(\ell,\theta)$. Finally, with $n_0=n_0(\ell,\theta)$ and $\ell\ge3$, $o_\chi\ge1$:
$2\ell N_\chi\le16o_\chi(\ell N_\chi)^2$, so $2\ell N_\chi4^{n_0}\le R_1(h,\ell,\theta)\le2^{c_2(\theta)}$, the inequality of $\mathrm{X}$.
\end{proof}

\begin{proposition}[Abstract count]\label{pf:eis:abscount}
Assume Lemma M and Lemma M$'$. Let $F$ be a frame over $\OO$ (with $\Gamma$ compact and totally disconnected and $S$ finite) satisfying
$(\mathrm{F1})$, $(\mathrm{F2})$, $(\mathrm{F3})$, $(\mathrm{F4})$, $(\mathrm{F5})$ and $(\mathrm{F6})$, $g_1$ marked, $(R,u)$ a universal point with $R_0$ a finite
$\OO$-module, $\varphi\in\Gamma$, and $\mathcal P$ a property of homomorphisms $\Gamma\to\eisGG(\Qlbar)$ such that $r_\xi$ has $\mathcal P$ whenever
$r_\xi$ is of type I, and every characteristic-zero point $r$ of type I with $\mathcal P$ is rigid and has $\tr r(\varphi)\ne\tr\eisrz(\varphi)$. Then
there are characteristic-zero points $r_1,\dots,r_n$ of type I with $\mathcal P$, pairwise non-isomorphic on $\ker\eta$, and $V\in\N$ such
that $\lambda_F(a)\le V$ for all $0\ne a\in\OO$ and
\[
 V\le2e_\OO\sum_{i=1}^n\eisv\bigl(\tr r_i(\varphi)-\tr\eisrz(\varphi)\bigr).
\]
\formallabel{UpperTheory.V.UnitaryThreeInertia.abstract_count}
\end{proposition}

\begin{proof}
The argument is that of \cref{pf:eis:upper,pf:eis:lower}, with an explicit presentation of $R$ in place of a model. We use the
following results of the development (all under $(\mathrm{F1})$, for marked $g_1$ and the universal point $u$).
\begin{enumerate}[label=(\arabic*)]
\item \lref{UpperTheory.V.UnitaryThreeInertia.exists_adapted}: there is a \emph{weighted presentation}
  $\pi_W:\OO[[x_s:s\in\sigma]]\to R$, surjective and homogeneous for weights $\mathrm{wt}(s)\in\{-2,\dots,2\}$ of the variables, adapted to $u$
  (the variables are matrix entries of $u$ minus their values at $\eisrz$, of the weight of the entry, and $\pi_W$ carries the weight
  grading to the decomposition $R_j$), and minimal.
\item \lref{UpperTheory.V.UnitaryThreeInertia.hasOrigin}: the presentation has an origin, an $\OO$-point $o$ of $R$ with all variables
  mapping to $0$, and $o\circ u=\eisrz$.
\item \lref{UpperTheory.V.UnitaryThreeInertia.exists_relations} (also using $(\mathrm{F2})$, $(\mathrm{F3})$): $\ker\pi_W$ is generated by
  elements $f_{j,i}$ of weight $j$, $1\le i\le\dim\mathrm{Ob}_j$.
\item \lref{UpperTheory.V.UnitaryThreeInertia.finrank_H1L_le}: $\dim H^1_{\mathrm{Sel}}(\bar{\mathfrak g}_j)\le\#\{s:\mathrm{wt}(s)=j\}$; with
  $(\mathrm{F3})$, $\dim\mathrm{Ob}_j+\epsilon_j\le\#\{s:\mathrm{wt}(s)=j\}$ for all $j$.
\item For $\xi:R_\bullet\to\Qlbar$, $r_\xi$ is a characteristic-zero point
  (\lref{UpperTheory.V.UnitaryThreeInertia.isCharZeroPt_of_isPtHom}); if $\xi$ is non-zero on some $R_j$ with $j<0$ and on some $R_j$
  with $j>0$ (we say $\xi$ is two-signed), then $r_\xi$ is neither upper nor lower
  (\lref{UpperTheory.V.UnitaryThreeInertia.isPtHom_sign}) and hence of type I
  (\lref{UpperTheory.V.UnitaryThreeInertia.typeI_of_twoSigned}, using $(\mathrm{F2})$ and $(\mathrm{F6})$); so $r_\xi$ has $\mathcal P$, is rigid,
  and $\tr r_\xi(\varphi)\ne\tr\eisrz(\varphi)$.
\item If $\xi$ vanishes on all $R_j$, $j\ne0$, then $r_\xi$ is upper and lower, hence equal to $\eisrz$
  (\lref{UpperTheory.V.UnitaryThreeInertia.eq_rz_of_diag}, using $(\mathrm{F5})$), and $\xi$ agrees with $o$ on all homogeneous elements
  (\lref{UpperTheory.V.UnitaryThreeInertia.isPtHom_act}: points with torus-conjugate specializations differ by the torus action).
\item At a two-signed $\xi$, every point derivation $D$ of $R_\bullet$ at $\xi$ is a multiple of the orbit derivation: $D(x)=a\,j\,\xi(x)$ for $x\in R_j$
  (\lref{UpperTheory.V.UnitaryThreeInertia.tangent_eq_orbit}, from the rigidity of $r_\xi$).
\item A tangent vector at $o$ of weight $j\le0$ which is killed by the linear parts of the relations is zero
  (\lref{UpperTheory.V.UnitaryThreeInertia.tangent_coord_eq_zero_of_nonpos_weight}, using $(\mathrm{F4})$).
\item For $0\ne a\in\OO$ there is an injective additive map from $\Sha_F(a)$ to $(\OO/a)^\sigma$, supported on the variables of weight
  $-1$ and killed by the linear parts of the relations (\lref{UpperTheory.V.UnitaryThreeInertia.mprime_coord}).
\end{enumerate}
Let $a=\tr\mathrm{mat}(u(\varphi))-\tr\eisrz(\varphi)\in R_0$; then $o(a)=0$, and $\xi(a)=\tr r_\xi(\varphi)-\tr\eisrz(\varphi)$, which is non-zero at
two-signed $\xi$ by (5). The coordinate form of Lemmas M and M$'$ (\lref{UpperTheory.V.UnitaryThreeInertia.Abs.coord_count}), applied
to the presentation, the relations, the bounds (4), the finiteness of $R_0$, the facts (6), (7), (8) and the function $a$, with the normalized
valuation $e_\OO\eisv$, gives $n$, two-signed $\xi_1,\dots,\xi_n$ no two of which are related by the torus action, and $V\in\N$, such that every
$\OO$-module with an injective map as in (9) has length $\le V$, and $V\le2\sum_ie_\OO\eisv(\xi_i(a))$. By (9), $\lambda_F(a)\le V$ for all
$a\ne0$. Put $r_i=r_{\xi_i}$; these are characteristic-zero points of type I with $\mathcal P$ by (5), and the bound for $V$ is the asserted one.
If $r_i$ and $r_{i'}$ are isomorphic on $\ker\eta$, then they are conjugate by some $\tau(t)$
(\lref{UpperTheory.V.UnitaryThreeInertia.conj_of_matIso}, for characteristic-zero points of type I), so $\xi_i$ and $\xi_{i'}$ are related by
the torus action (\lref{UpperTheory.V.UnitaryThreeInertia.isPtHom_act}) and $i=i'$.
\end{proof}

\begin{lemma}\label{pf:eis:vlenV}
Let $\mathrm{H}(\ell,e,\psi,o_1)$ hold, let $F$ be a host frame and $V\in\N$ with $\lambda_F(a)\le V$ for all $0\ne a\in\OO$. Then
$V_\ell(\psi,e)\le V$.
\formallabel{UpperTheory.V.UnitaryThreeInertia.vlen_le}
\end{lemma}

\begin{proof}
Let $S$ be the set of primes dividing $\ell N$. Since $\Psi^{-1}$ factors through $G_{\Q,S}$, the length $V_\ell(\psi,e)$ equals the length of
the group $\Sha'$ of classes in $H^1(G_{\Q,S},\OO\otimes\eisQZ(\Psi^{-1}))$ which vanish on the decomposition groups at the places $A_v$,
$v\in S$ (\lref{UpperTheory.V.Chr.litV}). The frame satisfies $(\mathrm{F1})$ (\lref{UpperTheory.V.UnitaryThreeInertia.hostFrame_ok}).
Every class of $\Sha'$ is killed by a power of $\ell$ (a representing cocycle takes finitely many values), so $\Sha'$ is the increasing
union of the submodules $\Sha'[\ell^k]$, and it suffices to show $\len_\OO\Sha'[\ell^k]\le\lambda_F(\ell^k)$ for every $k$. Let
$a=\ell^k$ and $M=\OO\otimes\eisQZ(\Psi^{-1})$; then $M$ is divisible and $M[a]$ is identified with $\OO/a$ (with the action through $\Psi^{-1}$).
If $z$ is a cocycle whose class lies in $\Sha'[a]$, then $az=dm_0$ is a coboundary, $m_0=am_1$, and $y=z-dm_1$ takes values in $M[a]=\OO/a$; it
is a continuous cocycle. For $v\in S$ we have $z=dm$ on $D_v$ for some $m\in M$. If $v\notin T_b$, then $\Psi$ is trivial on $N_v$ by $(\mathrm{F1})$,
so $y$ vanishes on $N_v$. If $v\in T_b$, choose $x\in D_v$ with $\Psi(x)^2-1\in\OO^\times$ (by $(\mathrm{F1})$); then $\Psi(x)-1$ is a unit, and from
$a(y(x))=0$ one gets $a(m-m_1)=0$, so $y=d(m-m_1)$ on $D_v$ with $m-m_1\in M[a]$. Hence the class of $y$ lies in $\Sha_F(a)$, and its image
under the natural map $J:H^1(\Gamma,\OO/a)\to H^1(\Gamma,M)$ is the class of $z$. So $\Sha'[a]\subset J(\Sha_F(a))$, and
$\len_\OO\Sha'[a]\le\len_\OO\Sha_F(a)=\lambda_F(a)\le V$.
\end{proof}

\begin{theorem}\label{pf:eis:countVof}
$(\mathrm{U}_i')$, Lemma M, Lemma M$'$, the local finiteness, Tate duality and Euler characteristic theorems over finite extensions of
$\Ql$, global finiteness, the global Euler characteristic formula, Poitou--Tate duality, $(\mathrm{F4})_{\mathrm h}$, $(\mathrm{F6})_{\mathrm h}$,
$(\mathrm{R}_i)$ and $(\mathrm{Num}_i)$ imply $\sfname{CountV}$.
\formallabel{UpperTheory.V.UnitaryThreeInertia.countV_of_finiteness_duality_euler_host}
\end{theorem}

\begin{proof}
Let $K=\Q(i)$ (\lref{UpperTheory.V.UnitaryThreeInertia.gaussian_field_range_eq_ker_quartic_char}: $\Q(i)$ exists as a CM field, and for every
embedding into $\Omega$, $G_{\Q(i)}=\ker\eta_4$ and $d_{\Q(i)}=-4$), and let $C$ be the constant of $(\mathrm{Num}_i)$. We show that
$C+1$ works in $\sfname{CountV}$, for every primitive $\chi$. Let $h$ be a host with character $\chi$, big for $t$, and let $\Omega$, an odd
prime $\ell$, $\OO$, $0<e<t$ and $\theta\in\eisX_h(\OO,e)$ in case V with $R_1(h,\ell,\theta)\le2^{c_2(\theta)}$ be given; put $N=N_h$.

By \cref{pf:eis:hostdata}, $\mathrm{H}(\ell,e,\theta,o_\chi)$ and $\mathrm{X}(\ell,\theta,N_\chi)$ hold. Choose a host frame $F$
(\lref{UpperTheory.V.UnitaryThreeInertia.exists_hostFrame}) and $\iota:\Qlbar\cong\C$. $(\mathrm{U}_i')$ gives a marked $g_1$ with its
properties. The frame satisfies $(\mathrm{F1})$ (\lref{UpperTheory.V.UnitaryThreeInertia.hostFrame_ok}) and the hypothesis needed for the
existence of universal points (\lref{UpperTheory.V.UnitaryThreeInertia.hostFrame_phi}, from global finiteness), so a universal point
$(R,u)$ exists (\lref{UpperTheory.V.UnitaryThreeInertia.exists_isUniv}); by $(\mathrm{U}_i')$, $R_0$ is finite over $\OO$ and every $r_\xi$ of
type I satisfies $\mathrm{B}'_\iota$. $(\mathrm{Num}_i)$ gives $\varphi\in\Gamma$.

We apply \cref{pf:eis:abscount} with $\mathcal P=\mathrm{B}'_\iota$. Its hypotheses hold: $(\mathrm{F2})$ by
\lref{UpperTheory.V.UnitaryThreeInertia.hostFrame_locVan} (from the local finiteness, duality and Euler characteristic theorems);
$(\mathrm{F3})$ by \lref{UpperTheory.V.UnitaryThreeInertia.hostFrame_count} (from these and the global finiteness, Euler characteristic and
Poitou--Tate theorems); $(\mathrm{F4})$ by $(\mathrm{F4})_{\mathrm h}$; $(\mathrm{F5})$ by \lref{UpperTheory.V.UnitaryThreeInertia.frame_class_one}; $(\mathrm{F6})$
by $(\mathrm{F6})_{\mathrm h}$; rigidity of characteristic-zero points with $\mathrm{B}'_\iota$ by $(\mathrm{R}_i)$, and $\tr r(\varphi)\ne\tr\eisrz(\varphi)$ by
$(\mathrm{Num}_i)$. We obtain points $r_1,\dots,r_n$ and $V$. By \cref{pf:eis:vlenV}, $V_\ell(\theta,e)\le V<\infty$. By $(\mathrm{Num}_i)$,
applied to $r_1,\dots,r_n$, $(\sum_i\eisv(\tr r_i(\varphi)-\tr\eisrz(\varphi)))\log\ell\le(\ell Ne)^C$. Hence
\[
 \eisVt_\ell(\theta,e)\log\ell\le\frac{V}{e_\OO}\log\ell\le2\Bigl(\sum_i\eisv\bigl(\tr r_i(\varphi)-\tr\eisrz(\varphi)\bigr)\Bigr)\log\ell
 \le2(\ell Ne)^C\le(\ell Ne)^{C+1},
\]
since $\ell Ne\ge2$.
\end{proof}

\begin{proof}[Proof of \cref{pf:eis:countV}]
Assume $\sfname{HypothesesH'}$, $\sfname{LemmaK}(C_K)$ and $(\mathrm{U}_i)$. By \cref{pf:eis:s6}, $(\mathrm{U}_i')$ holds. The local and global
cohomological theorems are part of component \ref{pf:vwin-he:H1} of $\sfname{HypothesesH'}$. We use:
\begin{itemize}
\item \lref{UpperTheory.V.X.So.hostVanE}: $\sfname{LemmaK}(C_K)$, Poitou--Tate duality and the local finiteness, duality and Euler
  characteristic theorems imply $(\mathrm{F4})_{\mathrm h}$;
\item \lref{UpperTheory.V.X.xpp}: $(\mathrm{F6})_{\mathrm h}$ follows from the statement that, for the residual characters of
  $G_\Q$ attached to products of a Dirichlet character of the type considered and a character modulo a power of $2$, every residual
  $1$-cocycle satisfying the relevant unramifiedness conditions is a coboundary; this statement holds by
  \lref{UpperTheory.V.X.corX}, from Lemma X (\lref{UpperTheory.V.X.lemmaX}) and the non-vanishing of residual Bernoulli
  numbers (\lref{UpperTheory.V.X.residual_bernoulli_nonvanishing});
\item \lref{UpperTheory.V.Z.J.hostRigidUE_alt}: $\sfname{HypothesesH'}$, the adjoint lifting theorem
  (\lref{UpperTheory.V.HV.DescentVanishingReductions.adjoint_lifting_of_cohomology}) and $(\mathrm{Unl})$
  (\lref{UpperTheory.V.Z.J.unlEnd3E} and \lref{UpperTheory.V.Z.J.linkedUnlE_alt}) imply $(\mathrm{R}_i)$;
\item \lref{UpperTheory.V.Z.J.hostNumberUE_alt}: $\sfname{HypothesesH'}$, $(\mathrm{Unl})$ and $(\mathrm{Num})$
  (\lref{UpperTheory.GalN.V.Win.N.numberU3U_of_prints}) imply $(\mathrm{Num}_i)$;
\item Lemma M (\cref{pf:vwin-he:lemmaM}) and Lemma M$'$ (\lref{UpperTheory.V.Alg.M.mprime_holds}).
\end{itemize}
\Cref{pf:eis:countVof} gives $\sfname{CountV}$.
\end{proof}

\begin{remark}\label{pf:eis:countVrem}
(1) In case V the bound $(\ell N_hj)^{C+1}$ depends on $\ell$, as allowed by $\mathrm{Count}(P)$ of \cref{pf:vwin-he:countdef}; the field $K=\Q(i)$ and the constant $C$ do not
depend on $\chi$. (2) The $2$-adic condition $R_1(h,\ell,\theta)\le2^{c_2(\theta)}$ is used only through \cref{pf:eis:hostdata}, to produce
$\mathrm{H}$ and $\mathrm{X}$. (3) As in case I, the Selmer group is bounded by the length of a cotangent space of weight $-1$ (via
(9) and Lemma M$'$), and this length by twice the sum of valuations of trace differences at the two-signed points (Lemma M).
\end{remark}

\subsubsection{Deformation theory of frames}\label{sec:pf:eis:framedef}

We now prove the deformation-theoretic results used above: the existence of universal points, the presentation of
universal rings by homogeneous relations (with the obstruction theory behind it), the dimension count $(\mathrm{F3})$ and the
vanishing statement $(\mathrm{F4})_{\mathrm h}$ for host frames, the coordinate form of Lemmas M and M$'$, and the rigidity statement
$(\mathrm{R}_i)$. Throughout, $F$ is a frame over $\OO$ indexed by a finite set $S$, $\Gamma$ is compact and totally
disconnected, $k$ is the (finite) residue field of $\OO$ and $\ell$ its characteristic. A matrix $X$ over a graded
$\OO$-algebra $A=\sum_iA_i$ is \emph{graded} if $X_{ab}\in A_{\mathrm{wt}(a,b)}$ for all $a,b$, where $\mathrm{wt}(a,b)$ is the weight of the
entry $(a,b)$ (\cref{pf:eis:GG}).

\begin{lemma}[Universal points]\label{pf:eis:univ}
Assume $(\mathrm{F1})$, and that $\Gamma$ satisfies Mazur's finiteness condition $\Phi_\ell$: every open subgroup of $\Gamma$ has only
finitely many continuous homomorphisms to $\Z/\ell\Z$. Then for every marked $g_1$ there is a universal point $(R,u)$.
\formallabel{UpperTheory.V.UnitaryThreeInertia.exists_isUniv}
\end{lemma}

\begin{proof}
Say that a property of $\OO$-algebra homomorphisms $f:R\to R'$ ($R'$ complete) is \emph{cut out} by an ideal $J\subset R$ if $f$ has
the property exactly when $f(J)=0$. An equality $f(x)=f(y)$ is cut out by $(x-y)$, and the conjunction of any family of properties
cut out by ideals $J_i$ is cut out by $\sum_iJ_i$.

By \lref{UpperTheory.V.U3.exists_univ} (representability of the lifts of $\bar r_0$ under $\Phi_\ell$, Mazur's theorem \cite{Mazur1989} in
the present setting; $\Phi_\ell$ passes to Hausdorff quotients of $\Gamma$, which is how the result is applied) there is a complete
$\OO$-algebra $R$ with a lift $u$ of $\bar r_0$ such that every lift to a complete $\OO$-algebra $R'$ is $f\circ u$ for a unique
$f:R\to R'$. We show that ``$f\circ u$ is a point'' is cut out by an ideal $J$. Then $R/J$, with the image of $u$, is a universal
point: a point $r$ to $R'$ is $f\circ u$ for a unique $f$, which kills $J$ and so factors uniquely through $R/J$. Moreover $R/J\ne0$,
since the base point $\eisrz(\eta,\Psi)$ is a point (\lref{UpperTheory.V.UnitaryThreeInertia.isPt_rz}) and the corresponding
homomorphism $R\to\OO$ kills $J$; and $R/J$ is again complete Noetherian local with residue field $k$.

The normalization at $(c,g_1)$ is cut out by the entries of $\mathrm{mat}(u(c))-1$ and the entries $(1,2),(2,1)$ of
$\mathrm{mat}(u(g_1))-\mathrm{mat}(u(cg_1^{-1}c))$, and $(\mathrm P_0)$ at $v\notin T_b$ by the entries of $\mathrm{mat}(u(s))-1$, $s\in N_v$
(note that $\eta(s)=1$ there). Let $v\in T_b$. By $(\mathrm{F1})$ there is $x\in D_v$ with $\eta(x)=1$ and $\Psi(x)^2-1\in\OO^\times$. Then
$X=\mathrm{mat}(u(x))$ is congruent to $\mathrm{diag}(\Psi(x),1,\Psi(x)^{-1})$, and the differences of these three values are units, being unit
multiples of $\Psi(x)-1$ or of $\Psi(x)^2-1$. By \lref{UpperTheory.V.UnitaryThreeInertia.exists_canon} there is
$h_0\in\GL_3(R)$ with $h_0\equiv1$ modulo $\mathfrak m_R$ and diagonal entries $1$ such that $h_0^{-1}Xh_0$ is diagonal. Such an $h_0$ is
unique: if $h_0'$ is another one, $h_0^{-1}h_0'$ is congruent to $1$ and conjugates one diagonal matrix with residually distinct entries
into another, hence is diagonal, hence equal to $1$. By uniqueness, $f(h_0)$ is the corresponding matrix for $f(X)$
(\lref{UpperTheory.V.UnitaryThreeInertia.canon_map}). We claim that $f\circ u$ satisfies $(\mathrm P_\ell)$ (if $v=v_\ell$),
resp.\ $(\mathrm P_b)$ (if $v\in S_b$), if and only if the condition holds with $h=f(h_0)$
(\lref{UpperTheory.V.UnitaryThreeInertia.cond_iff_canon}). Indeed, let $h$ be as in $(\mathrm P_\ell)$. Then
$T=h^{-1}f(X)h$ is upper triangular, congruent to $\mathrm{diag}(\Psi(x),1,\Psi(x)^{-1})$, so it is diagonalized by a unique upper
unitriangular matrix $n\equiv1$; for the diagonal matrix $d\equiv1$ which makes the diagonal entries of $hnd$ equal to $1$, $hnd=f(h_0)$ by
uniqueness, and conjugation by the upper triangular matrix $nd$ preserves upper triangularity and the diagonal entries. The same
argument applies to $(\mathrm P_b)$. Consequently $(\mathrm P_\ell)$ is cut out by the entries of negative weight of
$h_0^{-1}\mathrm{mat}(u(s))h_0$ for $s\in D_{v_\ell}\cap\ker\eta$ and by the differences between the diagonal entries of
$h_0^{-1}\mathrm{mat}(u(s))h_0$ and $\Psi(s),1,\Psi(s)^{-1}$ for $s\in I_{v_\ell}\cap\ker\eta$; and $(\mathrm P_b)$ by these elements for $v$ together
with the entries of positive weight. Summing all these ideals gives $J$.
\end{proof}

\paragraph{Presentations.} A \emph{weighted presentation} of $R$ is a surjective $\OO$-algebra homomorphism
$\pi_W:\OO[[x_s:s\in\sigma]]\to R$ from a power series ring in finitely many variables with weights $\mathrm{wt}(s)\in\Z$; it is
\emph{adapted} to $u$ if each $\pi_W(x_s)$ is an entry $\mathrm{mat}(u(g))_{ab}-\mathrm{mat}(\eisrz(g))_{ab}$ for some $g\in\ker\eta$ and
$\mathrm{wt}(s)=\mathrm{wt}(a,b)$, and $\pi_W$ maps the closure of the span of the monomials of weight $j$ onto $R_j$; it is \emph{minimal} if
every element of $\ker\pi_W$ has all its linear coefficients in $\mathfrak m_\OO$. Adapted minimal presentations exist under
$(\mathrm{F1})$ (\lref{UpperTheory.V.UnitaryThreeInertia.exists_adapted}). A surjection $p:A\to A'$ of complete $\OO$-algebras, with
$A$ finite and graded by $\OO$-submodules $A_i$ ($A_iA_{i'}\subset A_{i+i'}$), is a \emph{small extension of weight $j$} if
$I=\ker p$ satisfies $I\subset\mathfrak m_A$, $\mathfrak m_AI=0$ and $I\subset A_j$. Then $I$ is a $k$-vector space and $\Gamma$ acts on
$\mathfrak g_j(I)=\bar{\mathfrak g}_j\otimes_kI$ through $\bar r_0$.

\begin{lemma}[Obstructions]\label{pf:eis:obs}
Assume $(\mathrm{F1})$ and $(\mathrm{F2})$, let $g_1$ be marked, and let $p:A\to A'$ be a small extension of weight $j$ with
$\mathrm{Ob}_j$ finite and $\#\ker p>\#\mathrm{Ob}_j$. Let $r'$ be a point with values in $A'$ all of whose matrices are images of graded
matrices over $A$. Then $p$ factors as $A\xrightarrow{p_1}A''\xrightarrow{q}A'$ with $p_1$ surjective and $q$ not injective, and there is a
point $r''$ with values in $A''$ such that $q\circ r''=r'$.
\formallabel{UpperTheory.V.UnitaryThreeInertia.exists_unobstructed}
\end{lemma}

\begin{proof}
Let $I=\ker p$. \emph{Global and local lifts.} Since $A$ is finite, $r'$ is locally constant, and there is a locally constant map
$s:\Gamma\to\eisGG(A)$ with graded matrices, multiplier $\eta$ and $p\circ s=r'$. For each $v\in S$ the restriction of $r'$ to $D_v$ lifts to a
locally constant homomorphism $\rho_v:D_v\to\eisGG(A)$ with graded matrices which satisfies the local condition at $v$: at
$v=v_\ell$ by \lref{UpperTheory.V.UnitaryThreeInertia.exists_graded_lift_unitary_local_tame}, at $v\in S_b$ by
\lref{UpperTheory.V.UnitaryThreeInertia.exists_lift_d}, and at $v\notin T_b$ by
\lref{UpperTheory.V.UnitaryThreeInertia.exists_lift_u} (these lift the flag, resp.\ the diagonal form, resp.\ the triviality on
$N_v$; they use $(\mathrm{F2})$ and, for $v\notin T_b$, that $[I_v:N_v]$ divides $2\#k^\times$ and is prime to $\ell$). Write
$s(g)s(g')=(1+K(g,g'))s(gg')$ and $s(g)=(1+A_v(g))\rho_v(g)$ for $g\in D_v$. Then $K$ and $A_v$ take values in $\mathfrak g_j(I)$
(their images under $p$ vanish, and all matrices are graded with $I\subset A_j$), $K$ is a continuous $2$-cocycle, and
$dA_v=K|_{D_v}$, computed with the action through $\bar r_0$ since $\mathfrak m_AI=0$.

\emph{The obstruction.} For a $k$-linear $\phi:I\to k$, $(\phi\circ K,(\phi\circ A_v)_v)$ is a framed $2$-cocycle, and $\phi$ maps linearly
to its class in $\mathrm{Ob}_j$. As $\#\mathrm{Hom}_k(I,k)=\#I>\#\mathrm{Ob}_j$, there is $\phi\ne0$ whose class vanishes:
$\phi\circ K=d\beta$ and $\phi\circ A_v=\beta|_{D_v}+z_v$ with $z_v\in L_v(j)$. Let $A''=A/\ker\phi$ (an ideal, as $\mathfrak m_AI=0$), so that $q:A''\to A'$
has kernel $I/\ker\phi\cong k$, generated by an element $\epsilon$ with $\phi$ identified with the coordinate along $\epsilon$
(\lref{UpperTheory.exists_small_extension_factorization_of_functional}). Put $s''(g)=(1-\epsilon\beta(g))p_1(s(g))$. Since
$\epsilon\mathfrak m_{A''}=0$, $s''$ is a homomorphism (its defect is $\epsilon(\phi\circ K-d\beta)=0$), and on $D_v$ we have
$s''(g)=(1+\epsilon z_v(g))p_1(\rho_v(g))$. Since $z_v\in L_v(j)$, the twist of the local lift $p_1\circ\rho_v$ by the cocycle $z_v$ still
satisfies the local condition at $v$: vanishing on $N_v$ preserves $(\mathrm P_0)$ and $(\mathrm P_b)$, a coboundary changes the
conjugating matrix $h$ in $(\mathrm P_\ell)$ by $1+\epsilon Y$, vanishing on $I_{v_\ell}$ for $j=0$ preserves the diagonal characters, and for
$j>0$ a cocycle with values in $\mathfrak g_j$ preserves upper triangularity and the diagonal. Hence $s''$ is a lift of $\bar r_0$ with
values in $A''$ satisfying the local conditions, and $q\circ s''=r'$.

\emph{Normalization.} Finally, conjugating by a suitable matrix congruent to $1$ modulo $\ker q$ normalizes $s''$ at $(c,g_1)$
without changing $q\circ s''$ (\lref{UpperTheory.lift_along_surjection_graded}: since $\eta(c)=-1$ and $g_1$ is marked, the
normalization conditions can be solved successively, as for the base point); this gives $r''$.
\end{proof}

\begin{proposition}[Relations]\label{pf:eis:rel}
Assume $(\mathrm{F1})$, $(\mathrm{F2})$ and $(\mathrm{F3})$. Let $g_1$ be marked, $(R,u)$ a universal point and $\pi_W$ a minimal weighted
presentation adapted to $u$. Then $\ker\pi_W$ is generated by elements $f_{j,i}$, $j\in\Z$, $1\le i\le\dim_k\mathrm{Ob}_j$, with $f_{j,i}$
in the closure of the span of the monomials of weight $j$.
\formallabel{UpperTheory.V.UnitaryThreeInertia.exists_relations}
\end{proposition}

\begin{proof}
Let $P=\OO[[x_s]]$ and $J=\ker\pi_W$; write $d_j=\dim_k\mathrm{Ob}_j$, finite by $(\mathrm{F3})$.

\emph{Step 1: small quotients.} Let $\varphi:P\to A$ be a surjection onto a finite complete $\OO$-algebra, graded compatibly with the
weights of the monomials, and let $p:A\to A'$ be a small extension of weight $j$ with $\ker p=\varphi(J)$. We claim
$\#\ker p\le\#k^{d_j}$. Otherwise, let $\psi:R\to A'$ be induced by $p\circ\varphi$ (it vanishes on $J$). The point $\psi\circ u$ has matrices
which are images of graded matrices over $A$, since the entries of $u$ are images of variables of the right weight. By
\cref{pf:eis:obs} there are $A\to A''\to A'$ with $q$ not injective and a point $r''$ to $A''$ lifting $\psi\circ u$, and by universality
$r''=\sigma\circ u$ for some $\sigma:R\to A''$, with $q\circ\sigma=\psi$. The two homomorphisms $p_1\circ\varphi$ and $\sigma\circ\pi_W$ from $P$ to
$A''$ agree modulo $\ker q$, which is killed by $\mathfrak m_{A''}$. Their difference is therefore additive and vanishes on
$\mathfrak m_P^2$ (it satisfies $D(xy)=a(x)D(y)+D(x)b(y)\in\mathfrak m_{A''}\ker q=0$ for $x,y\in\mathfrak m_P$), and also on constants; since
$\mathfrak m_\OO x_s\subset\mathfrak m_P^2$, the two maps agree on every $f$ whose linear coefficients lie in $\mathfrak m_\OO$. By minimality this
applies to all $f\in J$, so $p_1(\varphi(J))=0$, i.e.\ $p_1$ kills $\ker p$. Then $q$ is injective on $A''=p_1(A)$: if $q(p_1(a))=0$ then
$a\in\ker p$ and $p_1(a)=0$. This contradiction proves the claim.

\emph{Step 2: truncation.} By the Artin--Rees lemma there is $N\ge1$ with $\mathfrak m_P^N\cap J\subset\mathfrak m_PJ$. Let $T_N\subset P$ be the
ideal of series whose coefficient at each monomial $x^n$ lies in $\mathfrak m_\OO^{\max(N-|n|,0)}$; then $T_N\subset\mathfrak m_P^N$, $T_N$ is
generated by homogeneous elements, and $B=P/T_N$ is a finite graded complete $\OO$-algebra. Let $\bar J$ be the (homogeneous)
image of $J$ in $B$. For each $j$ let $K_j\subset\bar J$ be the homogeneous ideal generated by $\mathfrak m_B\bar J$ and the homogeneous components of
$\bar J$ of weight $\ne j$; then $B/K_j\to B/\bar J$ is a small extension of weight $j$ with kernel the image of $\bar J$, so by Step 1 the
weight-$j$ part of $\bar J/\mathfrak m_B\bar J$ has at most $d_j$ dimensions. By the graded Nakayama lemma
(\lref{UpperTheory.V.UnitaryThreeInertia.Rel.exists_gens}) $\bar J$ is generated modulo $\mathfrak m_B\bar J$ by homogeneous elements $\bar f_{j,i}$
of weight $j$, $i\le d_j$.

\emph{Step 3: lifting.} Lift each $\bar f_{j,i}$ to an element of $J$ and take its homogeneous component of weight $j$, which lies in $J$
because $J$ is homogeneous; this gives $f_{j,i}\in J$ of weight $j$ with image $\bar f_{j,i}$. Let $J'\subset J$ be the ideal they generate.
For $x\in J$, the generation statement gives $y\in J'+\mathfrak m_PJ$ with $x-y\in T_N\cap J\subset\mathfrak m_P^N\cap J\subset\mathfrak m_PJ$; hence
$J=J'+\mathfrak m_PJ$, and $J=J'$ by Nakayama's lemma.
\end{proof}

\begin{lemma}[An Euler characteristic for Selmer and obstruction groups]\label{pf:eis:ecsel}
Let $j\in\Z$, assume that $\pi$ is continuous, that $H^1$ and $H^2$ of $\Gamma$ and of all $D_v$ with values in $\bar{\mathfrak g}_j$ are finite, and
that every family of continuous local $2$-cocycles $(\kappa_v)_{v\in S}$ is, up to local coboundaries, the restriction of a global
$2$-cocycle. Then $H^1_{\mathrm{Sel}}(\bar{\mathfrak g}_j)$ and $\mathrm{Ob}_j$ are finite and, with $h^i=\dim_kH^i$, $M=\bar{\mathfrak g}_j$,
\[
 \dim H^1_{\mathrm{Sel}}(M)-\dim\mathrm{Ob}_j=h^1(\Gamma,M)-h^2(\Gamma,M)+\sum_{v\in S}\bigl(h^2(D_v,M)-\dim Z^1(D_v,M)+\dim L_v(j)\bigr).
\]
\formallabel{UpperTheory.selmerComplex_finite_and_euler_characteristic}
\end{lemma}

\begin{proof}
There is an exact sequence
\[
\begin{aligned}
 0\to H^1_{\mathrm{Sel}}(M)\to H^1(\Gamma,M)&\to\bigoplus_vZ^1(D_v,M)/L_v(j)\to\mathrm{Ob}_j\\
 &\to H^2(\Gamma,M)\to\bigoplus_vH^2(D_v,M)\to0.
\end{aligned}
\]
The second map is restriction, well defined because local coboundaries lie in $L_v(j)$; its kernel is the Selmer group by definition.
The third map sends $(\alpha_v)$ to the class of $(0,(\alpha_v))$; this is zero if and only if $(0,\alpha)=(d\beta,\beta|_{D_v})+(0,z)$ with $z_v\in L_v(j)$,
i.e.\ $\beta$ is a global cocycle with $\alpha_v-\beta|_{D_v}\in L_v(j)$. The fourth map sends $(\kappa,\alpha)$ to the class of $\kappa$; if $\kappa=d\beta$,
then $(\kappa,\alpha)-(d\beta,\beta|_{D_v})=(0,\alpha-\beta|_{D_v})$ with $\alpha_v-\beta|_{D_v}$ a cocycle. Its image consists of the classes which are
locally coboundaries, which is the kernel of the last map, surjective by hypothesis. Taking alternating sums of dimensions gives the
formula.
\end{proof}

\begin{proposition}[Dimension count for host frames]\label{pf:eis:hfcount}
Assume the local finiteness, Tate duality and Euler characteristic theorems over finite extensions of $\Ql$, global finiteness, the
global Euler characteristic formula and Poitou--Tate duality. Then every host frame satisfies $(\mathrm{F3})$.
\formallabel{UpperTheory.V.UnitaryThreeInertia.hostFrame_count}
\end{proposition}

\begin{proof}
Let $F$ be a host frame, so $\Gamma=G_{\Q,S}$ with $S$ the primes dividing $\ell N$, $\ell$ odd; it satisfies $(\mathrm{F1})$
(\lref{UpperTheory.V.UnitaryThreeInertia.hostFrame_ok}) and $(\mathrm{F2})$ (\lref{UpperTheory.V.UnitaryThreeInertia.hostFrame_locVan}).
Fix $j$ and $M=\bar{\mathfrak g}_j$; then $\dim M=3-|j|$ for $|j|\le2$ and $M=0$ otherwise, and $\dim M^{c=1}$ is $1$ if $|j|\le1$ and $0$ otherwise
(\lref{UpperTheory.finrank_gradedPiece_and_complexConjugation_invariants}, an elementary computation with the action of $c$ by
$X\mapsto-J\,{}^tX\,J$).

\emph{Hypotheses of \cref{pf:eis:ecsel}.} Finiteness is part of the cited theorems. For surjectivity of $H^2(\Gamma,M)\to\bigoplus_vH^2(D_v,M)$ we
use \lref{UpperTheory.global_cohomology_finite_set}: by Poitou--Tate duality the cokernel is dual to $H^0(\Gamma,M^*)$,
$M^*=\mathrm{Hom}(M,\mu_\ell)$, which vanishes if some $\sigma\in G_\Q$ acts on $M$ by a scalar $a$ different from its cyclotomic character
$n$ modulo $\ell$ (then $\sigma$ acts on $M^*$ by $na^{-1}\ne1$). For $j\ne0$ take $\sigma$ in $D_{v_\ell}$ with $\eta(\sigma)=1$,
$\Psi(\sigma)^2-1\in\OO^\times$ and cyclotomic character $\equiv1$ modulo $\ell$
(\lref{UpperTheory.V.UnitaryThreeInertia.exists_sigma}): it acts on $M$ by $\bar\Psi(\sigma)^{\pm1}$ or $\bar\Psi(\sigma)^{\pm2}$, which is $\ne1$.
For $j=0$ take $\sigma$ fixing $\sqrt{-1}$ with cyclotomic character $\not\equiv1$ modulo $\ell$, which exists as $\Q(i)$ and $\Q(\mu_\ell)$ are
linearly disjoint (\lref{UpperTheory.exists_absGalois_fixing_sqrt_neg_one_cyclotomic_ne_one}); it acts trivially on
$\mathfrak g_0$.

\emph{Euler characteristics.} With $h^i=h^i(\Gamma,M)$ and $h^i_v=h^i(D_v,M)$: $h^1+\dim M^{c=1}=h^0+h^2+\dim M$ (global Euler characteristic
formula, \lref{UpperTheory.global_euler_characteristic_rationals_odd_prime}); $h^1_v=h^0_v+h^2_v+\dim M$ for $v=v_\ell$ and
$h^1_v=h^0_v+h^2_v$ otherwise (Tate's local formula, \lref{UpperTheory.local_euler_characteristic_rationals}); and
$\dim Z^1(D_v,M)=h^1_v+\dim M-h^0_v$. Thus the term of $v$ in \cref{pf:eis:ecsel} is $\dim L_v(j)-\dim M$ for $v\ne v_\ell$ and
$\dim L_{v_\ell}(j)-2\dim M$ for $v=v_\ell$.

\emph{Local conditions away from $v_\ell$.} Let $v\ne v_\ell$. The restriction of a cocycle $y\in L_v(j)$ to $I_v$ vanishes on $N_v$, a normal
subgroup of $I_v$ of index dividing $2\#k^\times$, hence prime to $\ell$; averaging over $I_v/N_v$ shows that $y|_{I_v}$ is a coboundary
$dm$, and conversely every $dm$ ($m\in M$) restricted to $I_v$ comes from $dm\in L_v(j)$ ($N_v$ acts trivially on $M$). The kernel of $y\mapsto y|_{I_v}$
consists of the cocycles vanishing on $I_v$, a space of dimension $\dim M^{I_v}$
(\lref{UpperTheory.finrank_unramified_cocycles_eq_finrank_invariants}: such a cocycle is determined by its value at the element
$\varphi_v$ of $(\mathrm{F1})$, which lies in $M^{I_v}$, and every value in $M^{I_v}$ occurs, using the homomorphisms $D_v\to\Z/n\Z$ of
$(\mathrm{F1})$). Hence $\dim L_v(j)=(\dim M-\dim M^{I_v})+\dim M^{I_v}=\dim M$, and the term of $v$ is $0$.

\emph{The place $v_\ell$.} If $j<0$, $L_{v_\ell}(j)$ consists of the coboundaries, of dimension $\dim M-h^0_{v_\ell}$. If $j=0$, $I_{v_\ell}$ acts
trivially on $\mathfrak g_0$ (as $\eta$ is trivial on it), so by the same lemma $\dim L_{v_\ell}(0)=\dim M$. If $j>0$, $L_{v_\ell}(j)=Z^1(D_{v_\ell},M)$
and the term is $h^2_{v_\ell}=0$ by $(\mathrm{F2})$. For $j\ne0$, $h^0=h^0_{v_\ell}=0$, since the element $x$ of $(\mathrm{F1})$ acts on $M$ by a
scalar $\ne1$; for $j=0$, $h^0=\dim M^{c=1}$, since elements of $\ker\eta$ act trivially on $\mathfrak g_0$ and the others act as $c$.

\emph{Conclusion.} By \cref{pf:eis:ecsel}, $\dim H^1_{\mathrm{Sel}}-\dim\mathrm{Ob}_j=h^0+\dim M-\dim M^{c=1}+t_{v_\ell}$, where $t_{v_\ell}$ is the term of $v_\ell$:
for $j<0$ this is $0+\dim M-\dim M^{c=1}-\dim M=-\dim M^{c=1}$, which is $-1$ for $j=-1$ and $0$ for $j\le-2$; for $j=0$ it is
$\dim M^{c=1}+\dim M-\dim M^{c=1}-\dim M=0$; for $j>0$ it is $\dim M-\dim M^{c=1}$, which is $2-1=1$ for $j=1$, $1-0=1$ for $j=2$ and $0$ for
$j\ge3$. These are the numbers $\epsilon_j$ of $(\mathrm{F3})$.
\end{proof}

\begin{proposition}[Coordinate form of Lemmas M and M$'$]\label{pf:eis:coord}
Assume Lemma M and Lemma M$'$. Let $\pi_W:\OO[[x_s:s\in\sigma]]\to R$ be a weighted presentation with weights in $\{-2,\dots,2\}$ and an origin
$o:R\to\OO$ (all $o(\pi_W(x_s))=0$); let $R_j=\pi_W(\text{weight }j)$ and $R_\bullet=\sum_jR_j$, and let $K\supset\OO$ be algebraically closed with a
normalized valuation $v$. Let $f_{j,i}$ ($j\in\Z$, $i\le b_j$) be homogeneous generators of $\ker\pi_W$ of weight $j$ with
$b_j+\epsilon_j\le\#\{s:\mathrm{wt}(s)=j\}$, and $a\in R_0$ with $o(a)=0$. Assume:
\begin{enumerate}[label=(m\arabic*)]
\item for $j\le0$, a vector $(d_s)\in\OO^\sigma$ supported on the variables of weight $j$ and orthogonal to the linear parts of all elements of
  $\ker\pi_W$ is zero;
\item $R_0$ is a finite $\OO$-module;
\item every $K$-point of $R_\bullet$ which vanishes on all $R_j$, $j\ne0$, agrees with $o$ on all $R_j$;
\item at every two-signed $K$-point $\xi$ of $R_\bullet$, every point derivation $D$ satisfies $D(x)=c\,j\,\xi(x)$ for $x\in R_j$, for some $c\in K$;
\item $\xi(a)\ne0$ at every two-signed $K$-point $\xi$.
\end{enumerate}
Then there are two-signed $K$-points $\xi_1,\dots,\xi_n$, no two in the same $K^\times$-orbit, and $V\in\N$ such that every $\OO$-module $M$
with an injective additive map $M\to(\OO/c)^\sigma$ ($0\ne c\in\OO$) supported on the variables of weight $-1$ and orthogonal to the linear
parts of the elements of $\ker\pi_W$ has length $\le V$, and $V\le2\sum_iv(\xi_i(a))$.
\formallabel{UpperTheory.V.UnitaryThreeInertia.Abs.coord_count}
\end{proposition}

\begin{proof}
If there is no variable of weight $-1$, take $n=0$ and $V=0$: the modules in question are zero. Otherwise let $A$ be the weighted ring
formed by the finite sums of homogeneous elements of $\OO[[x_s]]$ (homogeneous components being closures of spans of monomials of
a given weight), with base point $x$ the constant term (\cref{pf:vwin-he:wring}). Put $r=\#\mathrm{wt}^{-1}(-1)$, $s'=\#\mathrm{wt}^{-1}(-2)$,
$r'=\#\mathrm{wt}^{-1}(1)-1$, $s''=\#\mathrm{wt}^{-1}(2)-1$, $m=\#\mathrm{wt}^{-1}(0)$. The bounds $b_j+\epsilon_j\le\#\mathrm{wt}^{-1}(j)$ say that there are at most
$r+1$, $s'$, $r'$, $s''$, $m$ generators of weights $-1,-2,1,2,0$, and none of other weights; padding with zeros, the generators form a
system of equations of type $(r,s',r',s'',m)$ (\cref{pf:vwin-he:eqns}) generating the ideal $I=\ker\pi_W\cap A$, whose quotient $A/I$ is
$R_\bullet$ (\lref{UpperTheory.exists_weighted_equations_presenting_ker}). We verify (M1)--(M5):
(M1): $x$ is smooth with tangent dimensions $n_{r,s',r',s'',m}$ (the variables give a chart,
\lref{UpperTheory.smoothPt_weightedGradedAlgebra_basePoint}), and by (m1) the equations of weights $-1,-2,0$ span the cotangent
spaces in these degrees, since a functional on the cotangent lattice of degree $j$ killing them would be a vector as in (m1)
(\lref{UpperTheory.spansCot_of_linear_terms_nondegenerate}); (M2) is (m2), as $A_0/I_0\cong R_0$; (M3) holds since by (m3) a
$K$-point on $I$ which is negatively and positively trivial equals $x$; (M4) follows from (m4), as two derivations which are multiples of the
same derivation are linearly dependent; (M5) is (m5) for a homogeneous lift of $a$. Lemma M gives a finite set covering $I$; replace it by
a subset $T$ of representatives of distinct orbits which still covers $I$
(\lref{UpperTheory.exists_cover_with_distinct_orbit_representatives}), and put $V=\eiscot_{-1}(x;F)$ (the
saturated cotangent length of \cref{pf:vwin-he:cot}). Lemma M gives $V\le2\sum_{\xi\in T}v(\xi(a))$, the points of $T$ are two-signed points of
$R_\bullet$, and Lemma M$'$ (\lref{UpperTheory.length_le_cotLen_of_injective_tangent_map}) bounds the length of
every module $M$ as in the statement by $V$.
\end{proof}

\begin{proposition}[Vanishing for host frames]\label{pf:eis:vanE}
$\sfname{LemmaK}(C_K)$, Poitou--Tate duality, and the local finiteness, Tate duality and Euler characteristic theorems imply
$(\mathrm{F4})_{\mathrm h}$.
\formallabel{UpperTheory.V.X.So.hostVanE}
\end{proposition}

\begin{proof}
Let $F$ be a host frame, $j\le0$, and $\gamma:\Gamma\to\mathfrak g_j(\OO)$ a continuous cocycle for $\eisrz=\eisrz(\eta,\Psi)$ satisfying the hypotheses of
$(\mathrm{F4})$. Write $\bar a=4-a$ for $a\in\{1,2,3\}$.

\emph{Lines.} For $s\in\{0,1\}$ let $\theta_{j,s}=\Psi^j\eta^s$. For every entry $(a,b)$ of weight $j$, $z_{ab,s}=\gamma_{ab}-(-1)^s\gamma_{\bar b\bar a}$ is a
cocycle for the character $\theta_{j,s}$ (\lref{UpperTheory.V.X.So.vanEOf}: for $X\in M_3(\OO)$,
$(\mathrm{ad}\,\eisrz(g)X)_{ab}-(-1)^s(\mathrm{ad}\,\eisrz(g)X)_{\bar b\bar a}=\theta_{\mathrm{wt}(a,b),s}(g)(X_{ab}-(-1)^sX_{\bar b\bar a})$, since $\eisrz(g)$ is diagonal
for $\eta(g)=1$ and the twisted action of $c$ exchanges the entries $(a,b)$ and $(\bar b,\bar a)$ with a sign), and it inherits the local
hypotheses of $\gamma$. It suffices to show that each $z_{ab,s}$ becomes a coboundary after multiplication by some $0\ne a_{ab,s}\in\OO$: then
$2a\gamma$ is a coboundary for $a=\prod a_{ab,s}$, as $2\gamma_{ab}=z_{ab,0}+z_{ab,1}$.

\emph{Weight $0$.} Let $j=0$ and $z=z_{ab,s}$; $\theta=\eta^s$ is trivial on $\ker\eta$, so $z$ is additive on $\ker\eta$. It vanishes on $I_v\cap\ker\eta$ for
every $v$: for $v\in T_b$ by hypothesis, and for $v\notin T_b$ because $t^{\#k^\times}\in N_v$ for $t\in I_v\cap\ker\eta$, so that $\#k^\times z(t)=0$. By the cocycle
relation $z$ vanishes on all conjugates of these elements lying in $\ker\eta$. For each $N$, $z$ modulo $\mathfrak m_\OO^N$ is a locally constant
homomorphism on $\ker\eta$ trivial on these conjugates, hence trivial by $(\mathrm{F5})$, which holds for host frames
(\lref{UpperTheory.V.UnitaryThreeInertia.frame_class_one}). So $z=0$ on $\ker\eta$, a subgroup of index $2$, and then $2z$ is a
coboundary (\lref{UpperTheory.V.X.So.lineZero}).

\emph{Negative weight.} Let $j=-n<0$. Pulled back to $G_\Q$, $\theta_{j,s}^{-1}=\eiscyc^{en}\psi'$ with $\psi'=\psi_G^n\eta_4^{-s}$ of finite order, so $z$ is a
cocycle for $(\eiscyc^{en}\psi')^{-1}$, a character of negative weight. Locally, after multiplication by a non-zero constant, $z$ is a
coboundary on every decomposition group: at $\ell$ because it is one on the subgroup $D_{v_\ell}\cap\ker\eta$ of index $\le2$ (hypothesis of $(\mathrm{F4})$),
and at $p\ne\ell$ because local $H^1$ of a character of non-zero weight is torsion (\lref{UpperTheory.V.X.So.locTors}: $H^0$ and, by
local duality, $H^2$ vanish rationally, so $H^1$ does by Tate's local Euler characteristic formula). By
\lref{UpperTheory.V.X.soBar} (a cocycle of $G_\Q$ for $\psi^{-1}\eiscyc^{-e}$, $e\ge1$, with values in $\Qlbar$, which is locally a
coboundary on all inertia groups and at $\ell$, is a coboundary; a consequence of $\sfname{LemmaK}$ and Poitou--Tate duality), $z$ is a
coboundary over $\Qlbar$, hence $az$ is a coboundary over $\OO$ for some $0\ne a\in\OO$ (\lref{UpperTheory.V.X.So.lineNeg}).
\end{proof}

\begin{proposition}[Rigidity over $\Q(i)$]\label{pf:eis:rigi}
$\sfname{HypothesesH'}$, the adjoint lifting theorem
(\lref{UpperTheory.V.HV.DescentVanishingReductions.adjoint_lifting_of_cohomology}) and $(\mathrm{Unl})$ imply $(\mathrm{R}_i)$.
\formallabel{UpperTheory.V.Z.J.hostRigidUE_alt}
\end{proposition}

\begin{proof}
Let $F$ be a host frame, $g_1$ marked, $r$ a characteristic-zero point with $\mathrm B'_\iota(r)$, and $y$ a continuous cocycle of the adjoint
representation of $r$ defining a first-order deformation $r_\epsilon=(1+\epsilon y)r$ (with values in $\Qlbar[\epsilon]$) which satisfies the local
conditions. We show that $y$ is a coboundary in two steps.

\emph{Step 1: the Bloch--Kato Selmer group of $\mathrm{ad}\,r$ vanishes.} Write $r=\tau(t)^{-1}r'\tau(t)$ with $r'$ a point with values in a discrete
valuation ring $B$ finite over $\OO$. Then $r'$ is an ordinary point of the linked unitary deformation problem of
\cref{sec:pf:eis:onechar} for $\eta_4$ (\lref{UpperTheory.V.Pr.ordPt_of_isPt}): the local conditions of the host
frame at $v_\ell$ are the ordinarity conditions. The restriction $\rho'$ of $r'$ to $G_{\Q(i)}$ is isomorphic to that of $r$, so by
$\mathrm B'_\iota(r)$ it is irreducible on open subgroups and $\rho'\otimes\eiscyc^{-1}$ is automorphic of conjugate self-dual type with the
window condition. By $(\mathrm{Unl})$ $\rho'$ is unlinked of level $\ell N$ and type $(e,0,-e)$
(\lref{UpperTheory.V.Z.J.hostUnlE}), hence ``free'' at $\ell$ in the sense required by the adjoint lifting theorem
(\lref{UpperTheory.V.Z.J.freeCU}). The rigidity theorem for ordinary points of the linked problems,
\lref{UpperTheory.V.Th.rigEnd3} (a consequence of $\sfname{HypothesesH'}$ and the adjoint lifting theorem), then shows that the
Bloch--Kato Selmer group of $\mathrm{ad}\,r'$ vanishes, and this property is invariant under conjugation.

\emph{Step 2: $y$ is a Bloch--Kato cocycle.} At places outside $S$, $r$ and $y$ are trivial on inertia ($\Gamma=G_{\Q,S}$). At $v\in S_b$,
$(\mathrm P_b)$ for $r_\epsilon$ gives $h=h_0(1+\epsilon H_1)$ with $h^{-1}r_\epsilon(s)h=\mathrm{diag}(\Psi(s),1,\Psi(s)^{-1})$ for $s\in I_v$ ($\eta$ is trivial on
$I_v$); comparing with $h_0^{-1}r(s)h_0=\mathrm{diag}(\Psi(s),1,\Psi(s)^{-1})$ gives $y(s)=X-\mathrm{ad}(r(s))X$ on $I_v$ with $X=h_0H_1h_0^{-1}$
(\lref{UpperTheory.V.UnitaryThreeInertia.rig_diag}). At $v\notin T_b$, $y$ vanishes on $N_v$, a subgroup of finite index of $I_v$, and
in characteristic $0$ averaging over $I_v/N_v$ shows that $y|_{I_v}$ is a coboundary. At $v_\ell$, $(\mathrm P_\ell)$ for $r_\epsilon$ says that after
conjugation $r_\epsilon$ is upper triangular on $D_{v_\ell}\cap\ker\eta$ with diagonal entries $\Psi,1,\Psi^{-1}$ on inertia. There is an open normal
subgroup $H\subset D_{v_\ell}\cap\ker\eta$ on whose inertia $\Psi=\eiscyc^e$, containing an element whose cyclotomic character minus $1$ is a unit, at
which the diagonal entries of $r$ are pairwise distinct (\lref{UpperTheory.V.UnitaryThreeInertia.rig_sub} and \lref{UpperTheory.V.UnitaryThreeInertia.rig_eig}).
By \lref{UpperTheory.V.UnitaryThreeInertia.rig_loc} (a consequence of the local results of the literature package: a first-order
deformation of an ordinary representation which remains upper triangular, with fixed diagonal characters $\eiscyc^{e},1,\eiscyc^{-e}$ of
pairwise distinct Hodge--Tate weights on inertia, has its class in the Bloch--Kato subspace), $y|_{D_{v_\ell}}$ is a Bloch--Kato cocycle.
Hence $y$ is a Bloch--Kato cocycle, and by Step 1 it is a coboundary.
\end{proof}

% ---------------------------------------------------------------------------
\subsection{Density of automorphic points for \texorpdfstring{$\Uni(3)$}{U(3)}}\label{sec:pf:eis:dense}

The unitary automorphy statements $(\mathrm{U}_{\mathrm{ord}})$ and $(\mathrm{U}_i)$ are proved elsewhere by propagating automorphy
along chains of ordinary deformations, starting from \emph{seeds}: three-dimensional automorphic representations with prescribed
residual and local behaviour. The following statement provides seeds of arbitrarily large weight; they are symmetric squares of
modular forms congruent to Eisenstein series.

\begin{definition}[Seeds]\label{pf:eis:seed}
Let $L\subset\Omega$ be a CM field, $\ell$ a prime, $\iota:\Qlbar\cong\C$, $c\in G_\Q$ a complex conjugation, $n\ge1$,
$\psi_i,\tau_i:G_L\to\Qlbar^\times$ ($1\le i\le n$), $m\in\Z^n$ and $\kappa:G_L\to\Qlbar^\times$. A homomorphism $r:G_L\to\GL_n(\Qlbar)$ is a
\emph{seed} for $(\iota,c,\psi,\tau,m,\kappa)$ if:
\begin{enumerate}[label=(\alph*)]
\item $r$ is continuous and irreducible;
\item there is a symmetric $P\in\GL_n(\Qlbar)$ with ${}^tr(s)Pr(s')=P$ whenever $s'=csc^{-1}$;
\item the coefficients of the characteristic polynomial of $r(s)$ are congruent modulo the maximal ideal to those of
  $\prod_i(X-\psi_i(s))$, for all $s$;
\item at every place not above $\ell$, some inertia group acts diagonally through $\tau_1,\dots,\tau_n$;
\item at every place $A$ above $\ell$, $r|_{D_A}$ is upper triangular with diagonal characters $\kappa_i$ such that $\kappa_i=\eiscyc^{m_i}$ on
  $I_A$ and $\kappa_i\equiv\psi_i$ modulo the maximal ideal;
\item $\mathcal A_\iota(r\otimes\kappa)$ holds.
\end{enumerate}
\formal{UpperTheory.V.Sd.IsSeedAlt}
\end{definition}

\begin{definition}[Admissible auxiliary data]\label{pf:eis:fieldpair}
Let $K,L\subset\Omega$ be number fields, $\ell$ a prime, $c\in G_\Q$, $N,w,j\in\N$ and $\chi',\eta,\theta:G_\Q\to\Qlbar^\times$. A character
$\phi$ is \emph{split at $(N,w)$} if it has conductor dividing $N$ and $\phi(s)=1$ whenever $s$ acts on the $N$-th roots of unity by
$\zeta\mapsto\zeta^w$. The tuple satisfies the \emph{auxiliary conditions} if: $c$ is a complex conjugation; $\ell\ne2$, $N\ge1$,
$\ell\nmid N$, $w$ is a prime not dividing $N$, $\theta\ne1$, $\theta^\ell=1$ and $\theta$ has conductor dividing $w$; $K$ is imaginary
quadratic, $L$ is abelian over $\Q$ and $G_L\subset G_K$; $\ker\eta=G_K$; $\eta$ and $\chi'$ are split at $(N,w)$; $\theta$ is trivial on $G_L$;
$\chi'(c)=(-1)^j$; and there is $s\in G_L$ with $(\eiscyc^j\chi')(s)^2\not\equiv1$ modulo the maximal ideal. For $M\ge1$, the CM field $L$
is \emph{unramified and $M$-split} if $L$ is unramified over its maximal totally real subfield at all finite places and every finite
place of $L$ dividing $M$ is split over that subfield.
\formal{UpperTheory.V.Sd.FieldPairCharacterTripleConditions}
\formal{UpperTheory.V.Z.Sym.UnrSplit}
\end{definition}

\begin{definition}[Density statement]\label{pf:eis:densestmt}
$(\mathrm{D}_3)$ is the statement: let $K,L\subset\Omega$ with $L$ CM, $\ell$, $\iota$, $c$, $N,w,j$, $\chi',\eta,\theta$ satisfy the auxiliary
conditions, with $\chi'\ne\eta$; let $M\ne0$ with $2\ell\mid M$, $L$ unramified and $M$-split, and $\chi'|_{G_L}$ unramified at all places of $L$
not dividing $M$. Then for every $B\in\N$ there are $m\ge B$ with $m\equiv j\pmod{\ell-1}$ and $r:G_L\to\GL_3(\Qlbar)$, irreducible on open
subgroups, which is a seed for $(\iota,c,\psi,\tau,(m,0,-m),\eiscyc^{-1})$, where $\psi=(\eiscyc^j\chi',1,(\eiscyc^j\chi')^{-1})|_{G_L}$ and
$\tau=(\chi',1,\chi'^{-1})|_{G_L}$.
\formal{UpperTheory.V.Z.Sym.UnitaryThreeAutomorphyDense}
\end{definition}

For a two-dimensional $\rho$ let $\mathrm{ad}^0\rho=\mathrm{Sym}^2\rho\otimes(\det\rho)^{-1}$. The proof uses three statements of the
development.

\begin{definition}[Inputs of the density statement]\label{pf:eis:denseinputs}
\leavevmode
\begin{enumerate}[label=(\alph*)]
\item $(\mathrm{E})$: for $\ell,\Omega,\iota$, a complex conjugation $c$, and $N,w,\phi,\theta$ such that $(N,w,\theta)$ satisfy the conditions
  on $(N,w,\theta)$ listed in \cref{pf:eis:fieldpair}, $\phi$ is split at $(N,w)$ and $\phi\ne1$, there is $k_0$ such that for every $k\ge k_0$
  with $\phi(c)=(-1)^k$ there is $\rho:G_\Q\to\GL_2(\Qlbar)$ which is an \emph{Eisenstein representation} for $(\phi,\theta,k)$ --
  continuous, irreducible, attached via $\iota$ to a cuspidal automorphic representation of $\GL_2(\A_\Q)$, with
  $\det\rho=\eiscyc^{k-1}\phi\theta$, residually of characteristic polynomial that of $1\oplus\eiscyc^{k-1}\phi$, with inertia away from $\ell$
  acting diagonally through $(1,\phi\theta)$, and, at every place above $\ell$, upper triangular on the decomposition group with diagonal
  characters equal to $\eiscyc^{k-1}$ and $1$ on inertia and congruent to $\eiscyc^{k-1}\phi$ and $1$ -- and which satisfies two further
  conditions of the development (a condition $\mathrm{PR}$, and unlinkedness of $\mathrm{ad}^0\rho$).
\item $(\mathrm{T})$: for data as in $(\mathrm{D}_3)$ (with $\ell\mid M$), $m\ge2$ and an Eisenstein representation $\rho$ for $(\chi'\eta,\theta,m+1)$
  which is irreducible on open subgroups and satisfies the two further conditions, with $(\chi'\eta)(c)=(-1)^{m+1}$, there is
  $\kappa:G_L\to\Qlbar^\times$ such that $\rho|_{G_L}\otimes\kappa$ satisfies the condition $\mathcal W(\iota,M,m)$ of the development (a
  two-dimensional representation of $G_L$ which is continuous, automorphic of conjugate self-dual type, irreducible on open subgroups,
  unramified at the places not dividing $M$, and ordinary at $\ell$ with weight difference $m$).
\item $(\mathrm{Ad})$: for a CM field $L$, $\ell\ne2$, $M\ne0$ with $2\ell\mid M$, $L$ unramified and $M$-split, and $r:G_L\to\GL_2(\Qlbar)$ with
  $\mathcal W(\iota,M,m)$, the representation $\mathrm{ad}^0(r)\otimes\eiscyc^{-1}$ is automorphic of conjugate self-dual type and satisfies the
  window condition at all places.
\end{enumerate}
\formallabel{UpperTheory.V.Z.Sym.W.PropEGW}
\formallabel{UpperTheory.V.Z.Sym.W.SeedInW}
\formallabel{UpperTheory.V.Z.Sym.W.AdAutW}
\end{definition}

$(\mathrm{E})$ is the existence of cusp forms of large weight congruent to Eisenstein series, with the information on their
Galois representations needed below; $(\mathrm{Ad})$ is an automorphy statement for adjoint (symmetric square) lifts, which over
$\Q$ is the theorem of Gelbart and Jacquet \cite{GelbartJacquet1978}. They are proved from $\sfname{HypothesesH'}$ in the development
(\lref{UpperTheory.V.Z.Sym.W.prop_eg_w_of_printed_hida}, \lref{UpperTheory.V.Z.Sym.W.seedInW_alt}, and
\lref{UpperTheory.V.Z.Sym.W.adjoint_automorphy_w_of_printed_hida_window}, the last also from the statement $\sfname{Row0b}$ with the
local parameter predicate of \cref{sec:galois}).

\begin{proposition}\label{pf:eis:denseprop}
$(\mathrm{E})$, $(\mathrm{T})$ and $(\mathrm{Ad})$ imply $(\mathrm{D}_3)$.
\formallabel{UpperTheory.V.Z.Sym.W.adjoint_lift_from_eisenstein_of_propagation}
\end{proposition}

\begin{proof}
Let the data and $B$ be given. Since $\ker\eta=G_K$ and $K$ is quadratic, $\eta^2=1$; $\eta(c)=-1$, since $c\notin G_K$ ($K$ is imaginary). Put
$\phi=\chi'\eta$. Then $\phi$ is split at $(N,w)$ (both factors are), and $\phi\ne1$ (otherwise $\chi'=\eta^{-1}=\eta$). Moreover $\phi$ has finite
order, $\eta$ and $\theta$ are trivial on $G_L$ (as $G_L\subset G_K$), and $\chi'$ takes values of absolute value $1$. Let $k_0$ be given by
$(\mathrm{E})$ for $(N,w,\phi,\theta)$. Put $t=B+k_0+2$ and $m=j+(\ell-1)t$; then $m\ge t\ge B$, $m\equiv j\pmod{\ell-1}$ and $m+1\ge k_0$. As
$\ell-1$ is even, $\phi(c)=\chi'(c)\eta(c)=(-1)^{j+1}=(-1)^{m+1}$. Let $\rho$ be an Eisenstein representation for $(\phi,\theta,m+1)$ with the two
further properties.

Put $\Phi=\eiscyc^m\phi$ and $\Psi=\eiscyc^j\chi'$. For $s\in G_L$, $\Phi(s)\equiv\Psi(s)$ modulo the maximal ideal, since $\eta(s)=1$ and
$\eiscyc(s)^{(\ell-1)t}\equiv1$. By the last auxiliary condition there is $s\in G_L$ with $\Psi(s)^2\not\equiv1$, hence $\Phi(s)^2\not\equiv1$; with
this, \lref{UpperTheory.V.Sd.sI} (for $\ell\ne2$, $k\ge2$, $\phi$ of finite order and an Eisenstein representation $\rho$ for $(\phi,\theta,k)$:
if $(\eiscyc^{k-1}\phi)(s)^2\not\equiv1$ for some $s$, then $\rho$ is irreducible on open subgroups) shows that $\rho$ is irreducible on open
subgroups. Then $\mathrm{ad}^0\rho$ is irreducible on open subgroups (\lref{UpperTheory.V.Sd.ad0Irr}), so $r=\mathrm{ad}^0\rho|_{G_L}$ is
irreducible on open subgroups. We check that $r$ is a seed for $(\iota,c,\psi,\tau,(m,0,-m),\eiscyc^{-1})$.

(a) $r$ is continuous and irreducible. (b) A symmetric pairing $P$ with ${}^t\mathrm{ad}^0\rho(s)P\,\mathrm{ad}^0\rho(s')=P$ for $s'=csc^{-1}$ exists
for the adjoint representation of any two-dimensional $\rho$ (\lref{UpperTheory.V.Sd.exists_symmetric_pairing_adjoint_zero}; for $s'=csc^{-1}$
the trace form on trace-zero matrices composed with conjugation by $\rho(c)$). (c) Residually $\rho$ has the characteristic polynomial of
$1\oplus\Phi$, so $\mathrm{ad}^0\rho$ has that of $\Phi\oplus1\oplus\Phi^{-1}$ (\lref{UpperTheory.V.Sd.ad0Facts}), which on $G_L$ is congruent to
that of $\Psi\oplus1\oplus\Psi^{-1}$ (\lref{UpperTheory.V.Sd.shapeFacts}: congruences of characters are preserved by these operations).
(d) Away from $\ell$, inertia acts on $\rho$ diagonally through $(1,\phi\theta)$, so on $\mathrm{ad}^0\rho$ through $(\phi\theta,1,(\phi\theta)^{-1})$, which
on $G_L$ equals $(\chi',1,\chi'^{-1})=\tau$. (e) At $A\mid\ell$, $\rho$ is upper triangular on $D_A$ with diagonal characters equal to $\eiscyc^m$ and $1$ on inertia and congruent to
$(\Phi,1)$, hence $\mathrm{ad}^0\rho$ is upper triangular with diagonal characters $\eiscyc^{m},1,\eiscyc^{-m}$ on inertia, congruent to
$(\Phi,1,\Phi^{-1})\equiv(\Psi,1,\Psi^{-1})$ on $G_L$ (\lref{UpperTheory.V.Sd.order_at_adjoint_zero}). (f) By $(\mathrm{T})$ (with
$\ell\mid M$) there is $\kappa$ with $\mathcal W(\iota,M,m)$ for $\rho|_{G_L}\otimes\kappa$, and by $(\mathrm{Ad})$,
$\mathrm{ad}^0(\rho|_{G_L}\otimes\kappa)\otimes\eiscyc^{-1}$ is automorphic of conjugate self-dual type and satisfies the window condition at all
places. Since $\mathrm{ad}^0$ is insensitive to twists, $\mathrm{ad}^0(\rho|_{G_L}\otimes\kappa)=r$, and $\mathcal A_\iota(r\otimes\eiscyc^{-1})$ follows
(\lref{UpperTheory.V.Z.Win.Birth.mkCsdW'}, using that $r$ is potentially unlinked at $\ell$ by (e), as $m\ge2$:
\lref{UpperTheory.V.Z.Win.Birth.potOrd}).
\end{proof}

\begin{theorem}\label{pf:eis:dense}
$\sfname{HypothesesH'}$ implies $(\mathrm{D}_3)$. More precisely, $\sfname{HypothesesH'}$ and $\sfname{Row0b}$ (with the local parameter
predicate $\sfname{IsRecT}$) imply $(\mathrm{D}_3)$, and $\sfname{HypothesesH'}$ contains $\sfname{Row0b}$.
\formallabel{UpperTheory.V.Z.Sym.W.unitary_three_automorphy_dense_of_prints}
\formallabel{UpperTheory.V.Z.Sym.W.unitary_three_automorphy_dense_of_prints_and_row}
\end{theorem}

\begin{proof}
Given $\sfname{HypothesesH'}$ and $\sfname{Row0b}$, \lref{UpperTheory.V.Z.Sym.W.prop_eg_w_of_printed_hida} gives $(\mathrm{E})$,
\lref{UpperTheory.V.Z.Sym.W.seedInW_alt} gives $(\mathrm{T})$, and
\lref{UpperTheory.V.Z.Sym.W.adjoint_automorphy_w_of_printed_hida_window} gives $(\mathrm{Ad})$; \cref{pf:eis:denseprop} gives
$(\mathrm{D}_3)$. The statement $\sfname{Row0b}$ is part of the literature package \ref{pf:vwin-he:H1} contained in
$\sfname{HypothesesH'}$ (the existence of Galois representations attached to polarized regular algebraic cuspidal representations,
\cref{thm:row0b}).
\end{proof}

\subsubsection{Proofs of the inputs \texorpdfstring{$(\mathrm{E})$, $(\mathrm{T})$ and $(\mathrm{Ad})$}{(E), (T) and (Ad)}}

The first input rests on Ribet's construction of cusp forms congruent to Eisenstein series \cite{Ribet1976}, made explicit with an
auxiliary prime.

\begin{lemma}[Eisenstein-congruent eigenforms]\label{pf:eis:propEC}
Let $\ell$ be an odd prime and $\iota:\Qlbar\cong\C$. Let $M>1$ with $\ell\nmid M$, $\phi_0$ a primitive Dirichlet character modulo $M$, $w$ a
prime with $w\nmid M$, $\ell\mid w-1$ and $\phi_0(w)=1$, $\theta_0\ne1$ a character modulo $w$ with $\theta_0^\ell=1$, and $\chi=\phi_0\theta_0$ (modulo
$Mw$). There is $k_0$ such that for every $k\ge k_0$ with $\phi_0(-1)=(-1)^k$ there is a normalized cuspidal eigenform $g$ of weight $k$,
level $\Gamma_1(Mw)$ and character $\chi$, whose conjugates under $\mathrm{Aut}(\C)$ are again eigenforms of this kind, with
\[
 \iota^{-1}\bigl(a_p(g)\bigr)\equiv\iota^{-1}\bigl(1+\chi(p)p^{k-1}\bigr)\quad\text{modulo the maximal ideal, for all primes }p.
\]
\formallabel{UpperTheory.V.Z.Sym.Unl.propEC}
\end{lemma}

\begin{proof}
The character $\chi$ is primitive modulo $Mw$, and $\chi(-1)=\phi_0(-1)=(-1)^k$ since $\theta_0$ has odd order. Let $E$ be the Eisenstein series
of weight $k$ and character $\chi$ with $a_n(E)=\sigma_{k-1,\chi}(n)=\sum_{d\mid n}\chi(d)d^{k-1}$ for $n\ge1$; its constant term $a_0(E)$ is a
multiple of $L(1-k,\chi)$ and is $\ell$-adically small: modulo the maximal ideal $\theta_0\equiv1$, and the Euler factor at $w$ gives
$L(1-k,\phi_0\theta_0)\equiv L(1-k,\phi_0)(1-\phi_0(w)w^{k-1})\equiv0$, as $\phi_0(w)=1$ and $w\equiv1$ modulo $\ell$. By
\lref{UpperTheory.V.Z.Sym.Unl.unitFormC} there is $k_0$ such that for every $k\ge k_0$ of the right parity there is a modular form $G$
of weight $k$, level $Mw$ and character $\chi$ with $\ell$-integral coefficients, constant term $1$ at $\infty$ and $0$ at the other cusps at
which $E$ has no constant term, whose conjugates are forms of the same kind (a combination of Eisenstein series with coefficients
$\chi(a)^{-1}$, normalized by a unit). Then $f=E-a_0(E)G$ is a cusp form (\lref{UpperTheory.V.Z.Sym.Unl.cusp_comb}) with
$\ell$-integral coefficients, $a_1(f)=1-a_0(E)a_1(G)$ is a unit, and $a_n(f)\equiv\sigma_{k-1,\chi}(n)$. The Hecke relations of $E$ show that
$T_pf-(1+\chi(p)p^{k-1})f=-a_0(E)\bigl(T_pG-(1+\chi(p)p^{k-1})G\bigr)$ has coefficients in the maximal ideal. By the Deligne--Serre lifting lemma
\cite[Lemme~6.11]{DeligneSerre1974}, in the form \lref{UpperTheory.V.Z.Sym.Unl.eigLiftC}, there is a normalized eigenform $g$ of the
same weight, level and character whose eigenvalues are congruent to $1+\chi(p)p^{k-1}$; the conjugates of $g$ are eigenforms because those
of $f$, $E$ and $G$ are of the same kind.
\end{proof}

\begin{proposition}\label{pf:eis:propE}
$\sfname{HypothesesH'}$ implies $(\mathrm{E})$.
\formallabel{UpperTheory.V.Z.Sym.W.prop_eg_w_of_printed_hida}
\end{proposition}

\begin{proof}
Let $\ell,\Omega,\iota,c,N,w,\phi,\theta$ be as in $(\mathrm{E})$. Write $\phi$ and $\theta$ as the Galois characters of Dirichlet characters $\phi_0$, primitive of
conductor $M\mid N$, and $\theta_0$ modulo $w$; then $\theta_0\ne1$ and $\theta_0^\ell=1$, so $\ell\mid w-1$; $M>1$ since $\phi\ne1$; $\ell\nmid Mw$; and $\phi_0(w)=1$,
since $\phi$ is split at $(N,w)$. \Cref{pf:eis:propEC} gives $k_0$; let $k\ge\max(k_0,3)$ with $\phi(c)=(-1)^k$, and let $g$ be the eigenform. By the
results of the literature package contained in $\sfname{HypothesesH'}$ (the Galois representations attached to cuspidal eigenforms,
with their local properties, in particular ordinarity at $\ell$ when the $U_\ell$-eigenvalue is a unit), there is a continuous
$r:G_\Q\to\GL_2(\Qlbar)$ attached to $g$, ordinary at $\ell$, with $\det r=D^{-1}$, $D=\eiscyc^{k-1}\phi\theta$. Put $\rho=r\otimes D$. Then:
$\det\rho=D$; $\tr\rho(\Frob_p)\equiv1+\Phi(\Frob_p)$ for $p\nmid\ell Mw$, with $\Phi=\eiscyc^{k-1}\phi$, since $\theta\equiv1$ modulo the maximal ideal,
so the residual characteristic polynomial is that of $1\oplus\Phi$ (Chebotarev); $\rho$ is irreducible because $g$ is cuspidal
(\lref{UpperTheory.V.Sd.ordinary_irreducible_of_printed_hida}) and is attached to a cuspidal automorphic representation of
$\GL_2(\A_\Q)$ (a twist of that of $g$); at $p\nmid\ell$, inertia acts diagonally through $(1,\phi\theta)$
(\lref{UpperTheory.V.Z.Win.RR.ordAway_a_alt}, using that $\chi$ is primitive); and at $\ell$, $\rho$ is upper triangular on the
decomposition group with diagonal characters $\eiscyc^{k-1}$ and $1$ on inertia, congruent to $\Phi$ and $1$
(\lref{UpperTheory.V.Z.Win.RR.ordAbove_of_printsH}: the unramified quotient has Frobenius eigenvalue the unit root, which is
$\equiv1$). The condition $\mathrm{PR}$ (the attached cuspidal representation has ``plain'' local components) follows from ordinarity
(\lref{UpperTheory.V.Z.Win.Birth.prOrd}). Finally, by Deligne's estimates the Frobenius eigenvalues $\alpha,\beta$ of $\rho$ at unramified
$p\ne\ell$ satisfy $|\iota\alpha|<p|\iota\beta|$ and $|\iota\alpha|^2\ne p|\iota\beta|^2$, and these imply that $\mathrm{ad}^0\rho$ is unlinked
(\lref{UpperTheory.V.Z.Sym.Unl.adUnl}).
\end{proof}

\begin{proposition}\label{pf:eis:propT}
$\sfname{HypothesesH'}$ implies $(\mathrm{T})$.
\formallabel{UpperTheory.V.Z.Sym.W.seedInW_alt}
\end{proposition}

\begin{proof}
Let the data, $m\ge2$ and $\rho$ be as in $(\mathrm{T})$, and put $\Theta=\chi'\eta\theta$, a character of finite order with $\Theta(c)=(-1)^{m+1}$ (as
$\theta(c)=1$, $\theta$ having odd order), $\det\rho=\eiscyc^m\Theta$ and $\Theta|_{G_L}=\chi'|_{G_L}$ (as $\eta,\theta$ are trivial on $G_L$), unramified
outside $M$.

\emph{The twist.} By \lref{UpperTheory.V.Z.Sym.Pl.twEx} (an existence theorem for algebraic Hecke characters of the CM field $L$
with prescribed restriction to the ideles of $L^+$ and prescribed infinity type, proved by global class field theory from the literature
package, \lref{UpperTheory.V.Z.Sym.Pl.lemTWM}) there is a continuous algebraic character $\kappa:G_L\to\Qlbar^\times$, unramified at
the places not dividing $M$, with $\kappa^k$ a power of the cyclotomic character on inertia above $\ell$ for some $k\ge1$, and with
$\kappa(s)\kappa(csc^{-1})=\Theta(s)^{-1}\eiscyc(s)^{-(m+1)}$ for $s\in G_L$. Then $\det\rho(s)\kappa(s)\kappa(csc^{-1})=\eiscyc(s)^{-1}$, so
$\rho|_{G_L}\otimes\kappa$ is conjugate self-dual with multiplier $\eiscyc^{-1}$ for every complex conjugation (any two complex conjugations
differ by an element of $G_L$).

\emph{Automorphy.} $\rho|_{G_K}$ is automorphic, by base change from $\Q$ to the imaginary quadratic field $K$
(\lref{UpperTheory.V.Sd.up_quad_of_printed_results}), and so is $\rho|_{G_L}$, by solvable base change along the abelian extension
$L/K$ (\lref{UpperTheory.V.Z.Bc.upSol_of_alt}), using that $\rho$ is irreducible on open subgroups. Twisting by the algebraic character
$\kappa$ preserves automorphy (\lref{UpperTheory.V.Sd.twist_conjugate_selfDual}), and with the pairing above,
$\rho|_{G_L}\otimes\kappa$ is automorphic of conjugate self-dual type
(\lref{UpperTheory.GalN.V.aut_csd_zero_of_complex_conj_pair_automorphic}). The window condition holds for $\rho$ over $K$ by the
condition $\mathrm{PR}$, is preserved by base change to $L$ and by the twist, and at the places above $\ell$ uses that $\rho$ is ordinary
(\lref{UpperTheory.V.Z.Win.Birth.potOrd}).

\emph{The remaining properties.} $\mathrm{ad}^0(\rho|_{G_L}\otimes\kappa)=\mathrm{ad}^0\rho|_{G_L}$ is unlinked, since $\mathrm{ad}^0\rho$ is
(\lref{UpperTheory.V.Z.Sym.W.adjoint_zero_twist} and \lref{UpperTheory.GalN.V.L.wup}); irreducibility on open subgroups is inherited
from $\rho$; at a place not dividing $M$, inertia acts on $\rho$ through $(1,\Theta)$ and $\Theta=\chi'$ and $\kappa$ are unramified there, so the twist
is unramified; and at a place above $\ell$, the ordinary flag of $\rho$ with diagonal characters $\eiscyc^m$ and $1$ on inertia, twisted by $\kappa$,
gives an ordinary flag with weight difference $m$. Hence $\rho|_{G_L}\otimes\kappa$ satisfies $\mathcal W(\iota,M,m)$.
\end{proof}

The third input is an automorphic adjoint lift for unitary groups in two and three variables, with control of the local components.

\begin{proposition}\label{pf:eis:propAd}
$\sfname{HypothesesH'}$ and $\sfname{Row0b}$ imply $(\mathrm{Ad})$.
\formallabel{UpperTheory.V.Z.Sym.W.adjoint_automorphy_w_of_printed_hida_window}
\end{proposition}

\begin{proof}
We use two statements proved in the development.
\begin{enumerate}[label=(\alph*)]
\item A descent package compatible with weights (\lref{UpperTheory.V.Z.Sym.W.weight_compatible_descent}, from
  $\sfname{HypothesesH'}$ and $\sfname{Row0b}$): for CM fields $L$ unramified over $L^+$, automorphic representations of unitary groups
  over $L^+$ split by $L$ have Galois representations satisfying local--global compatibility at the split places, compatible with the
  Galois representations of their base change to $\GL_n(\A_L)$, and conjugate self-dual cuspidal representations of $\GL_n(\A_L)$,
  $2\le n\le4$, regular algebraic, descend to the unitary groups.
\item An adjoint lift for $U(2)$ with prescribed local components
  (\lref{UpperTheory.V.Z.Sym.W.adjoint_lift_w_of_literal_automorphy_spec_out}, from the literature package and a local
  statement on Weil--Deligne representations and depths, \lref{UpperTheory.V.Z.Sym.W.specOutS}): let $\sigma$ be an automorphic
  representation of a unitary group in two variables over $L^+$, of weight $j\ge1$ at all infinite places, unramified outside a finite
  set $S_0$ of split places, and let $P$ be a set of split places disjoint from $S_0$ at which the local components of $\sigma$ are
  unramified regular principal series; prescribe the local components of $\sigma$ at $S_0\cup P$ as principal series $(\psi_{w,0},\psi_{w,1})$, with
  $\psi_{w,0}/\psi_{w,1}$ not linked at $S_0$. Then either there is an automorphic representation $\sigma_3$ of a unitary group in three variables
  whose local components at $S_0\cup P$ are the principal series $(\psi_{w,0}/\psi_{w,1},1,\psi_{w,1}/\psi_{w,0})$ and which is unramified wherever
  $\sigma$ is, or $\sigma$ is endoscopic (induced from a character of a quadratic extension).
\end{enumerate}
Let $r$ satisfy $\mathcal W(\iota,M,m)$, and fix a place $w_0\mid M$ of $L$.

\emph{The adjoint lift.} By (a), $r$ is attached to an automorphic representation $\sigma$ of the unitary group in two variables, of weight
$m-1$ (\lref{UpperTheory.V.Z.Sym.W.aboveWeight}). Let $S$ consist of $w_0$ and one place above each place of $L^+$ dividing $M$ (all split, as
$L$ is $M$-split), and let $Q$ be an infinite set of split places not dividing $M$ at which $r$ has a Frobenius element whose eigenvalue ratio avoids the
finitely many degenerate values excluded in \lref{UpperTheory.V.Z.Sym.Ga.Bad} (\lref{UpperTheory.V.Z.Sym.Ga.select}). The local components of $\sigma$ at $S\cup Q$ are principal series
whose characters are read off from $r$: at $Q$ they are unramified and regular
(\lref{UpperTheory.V.Z.Sym.W.good_w_statement}); at the places of $S$ above $\ell$ they come from the ordinary flag of $r$ with weight gap
$m\ge2$ (\lref{UpperTheory.V.Z.Sym.W.above_w_statement}); at the other places of $S$ they are not linked because
$\mathrm{ad}^0r$ is unlinked (\lref{UpperTheory.V.Z.Sym.W.away_w_statement}). Apply (b) with $S_0=S$ and $P=Q$. The endoscopic
alternative is impossible: an endoscopic $\sigma$ would make $r$ induced from a character, contradicting irreducibility on open subgroups
at the places of $Q$ (\lref{UpperTheory.V.Z.Sym.W.branch_two_w_statement}). In the other alternative, the base change of
$\sigma_3$ to $\GL_3(\A_L)$ is a regular algebraic conjugate self-dual cuspidal representation $\pi$ attached via $\iota$ to
$\mathrm{ad}^0r\otimes\eiscyc^{-1}$, unramified at the places not dividing $M$
(\lref{UpperTheory.V.Z.Sym.W.branch1W_alt}), whose local component at $w_0$ is a principal series of the plain type required by the
window condition (\lref{UpperTheory.V.Z.Sym.W.plain_at_w_statement}).

\emph{Gluing.} So $\mathrm{ad}^0r\otimes\eiscyc^{-1}$ is automorphic of conjugate self-dual type. The window condition is a condition on each place
separately: at the places not dividing $M$ it holds for the lift obtained with any $w_0$, the local components being unramified with Satake
parameters of the right size; at a place $w\mid M$ it holds for the lift obtained with $w_0=w$; since all these lifts are attached to the same
Galois representation, they satisfy the condition simultaneously (\lref{UpperTheory.V.Z.Sym.W.glueW}).
\end{proof}
% ---------------------------------------------------------------------------
\subsection{The symplectic automorphy statement}\label{sec:pf:eis:sp}

Cases II and III are treated with four-dimensional representations of $G_\Q$ of $\GSp_4$-type whose residual representations are
sums of four characters with Hodge--Tate weights $(j+1,j,1,0)$. The counting argument for these cases is carried out elsewhere
(\lref{UpperTheory.V.HV.DescentVanishingReductions.countII_and_countIII_of_printsH_of_sp4In}); its automorphic input is the
statement $(\mathrm{Sp})$, which we now state precisely and reduce to its inputs.

\begin{definition}[Raw data]\label{pf:eis:raw}
Let $\ell$ be a prime and $\OO$ a coefficient ring. A \emph{raw datum} $D$ over $\OO$ consists of a finite set $S$ of primes, $e\in\N$,
characters $\psi,\xi,\omega:G_\Q\to\OO^\times$, elements $c,\phi,g_0\in G_\Q$, for each $p\in\N$ a place $A_p$ of $\Omega$ and an element
$\sigma_p\in G_\Q$, and an element $s_0\in D_{A_\ell}$. Put
\[
 \Psi=\eiscyc^e\psi,\quad\alpha=\eiscyc\,\omega^{-1}\xi,\quad\beta=\Psi\alpha^{-1},\quad\vartheta=\alpha\Psi,\quad
 r_0=\mathrm{diag}(\vartheta,\Psi,\alpha,1),
\]
a homomorphism $G_\Q\to\GL_4(\OO)$ with Hodge--Tate weights $(e+1,e,1,0)$, and $\bar r$ the reduction of $r_0$. Let $S_b$ be the set of $p\in S$, $p\ne\ell$, with
$\ell\mid\#\psi(I_{A_p})$, and $T=\{\ell\}\cup S_b$. A matrix $g\in M_4(R)$ is in $\GSp_4$ with multiplier $\nu$ if
$\mathrm{adj}(g)g=\nu\cdot1$, where $\mathrm{adj}(X)_{ij}=\epsilon_i\epsilon_jX_{5-j,5-i}$ with $\epsilon=(1,1,-1,-1)$; $r_0$ takes values in $\GSp_4$ with
multiplier $\vartheta$. The entry $(i,j)$ has weight $(\lambda_i-\lambda_j,\mu_i-\mu_j)\in\Z^2$, with $\lambda=(1,0,1,0)$ and $\mu=(1,1,0,0)$.

The datum is \emph{good} if: $\ell\ne2$; $S$ consists of primes and contains $2$ and $\ell$; $e\ge2$; $A_p$ lies over $p$ for every prime $p$; $c$ is a
complex conjugation; $\psi$ has open kernel and is unramified outside $S$; $\Psi(c)=1$; $\omega^{\ell-1}=1$ and $\omega\equiv\eiscyc$ modulo the
maximal ideal; $\xi$ has open kernel and order prime to $\ell$, is unramified outside $S$ and at the primes of $T$, and $\xi(c)=-1$; $\sigma_v$ is
a Frobenius element at $A_v$ for $v\in T$ and $\eiscyc(\sigma_\ell)=1$; $\bar\xi\bar\Psi^i(\sigma_v)\ne1$ for $v\in T$ and $i\in\Z$; for
$\gamma\in\{\alpha,\beta,\vartheta\}$ and its inverse, every continuous residual cocycle for $\bar\gamma$ which is unramified outside $T$ (in the
sense above) and locally a coboundary at the primes of $T$ is a coboundary; $\phi=\lim_n\sigma_\ell^{\ell^{n!}}$; $\bar\alpha(g_0)\ne1$ and
$\bar\Psi(g_0)\ne1$; and $s_0\in I_{A_\ell}\cap\ker\psi\cap\ker\omega$ with $\eiscyc(s_0)$ of infinite order.
\formal{UpperTheory.V.GradeZeroInterpolatingElement.Raw}
\formal{UpperTheory.V.GradeZeroInterpolatingElement.Raw.Good}
\end{definition}

\begin{definition}[Points and the conditions of $(\mathrm{Sp})$]\label{pf:eis:rawpts}
Let $D$ be a raw datum. A \emph{point} of $D$ over a complete $\OO$-algebra $R$ is a homomorphism $r:G_\Q\to M_4(R)$ with continuous entries
reducing to $\bar r$ such that: $r(g)$ is in $\GSp_4$ with multiplier $\vartheta(g)$; $r(g)=1$ for $g$ in inertia groups at places not above
$T$ with $\bar r(g)=1$; $r(\phi)$ is block diagonal for $\lambda$ and the entries $(1,3)$, $(3,1)$ of $r(g_0)+\vartheta(g_0)r(g_0^{-1})$ vanish; and
the restriction of $r$ to $D_{A_\ell}$ (resp.\ $D_{A_p}$, $p\in S_b$) factors through a domain on which the base representation $r_0$ satisfies
the local condition at $\ell$ (multiplier $\vartheta$, and ordinarity with respect to the eigenvalues of $r(s_0)$ and the characters
$(\vartheta,\Psi,\alpha,1)$ on inertia), resp.\ at $p$ (multiplier $\vartheta$ and triviality on $I_{A_p}\cap\ker\psi$). A $\Qlbar$-point is the image of
a point under a homomorphism $R\to\Qlbar$. Then:
\begin{itemize}
\item $(\mathrm{S1})$: for every point $r$ over $R$, the $\OO$-subalgebra of $R$ generated by the products of entries of $r$ of total
  weight $0$ is a finite $\OO$-module;
\item $(\mathrm{S2})$: for every $\Qlbar$-point $r$ which is block diagonal for $\lambda$, the matrices formed by the entries
  $(1,1),(1,3),(3,1),(3,3)$ of the $r(g)$ span a proper subspace of $M_2(\Qlbar)$;
\item for $\iota:\Qlbar\cong\C$, $\mathrm{Aut}_4(\iota,r)$ means: if $r$ is continuous, then for every imaginary quadratic field $K$ such that
  $r|_{G_K}$ is irreducible, $r|_{G_K}$ is attached via $\iota$ to a regular algebraic cuspidal automorphic representation of $\GL_4(\A_K)$
  and satisfies the window condition;
\item $(\mathrm{S3})_\iota$: every $\Qlbar$-point $r$ whose image spans $M_4(\Qlbar)$ satisfies $\mathrm{Aut}_4(\iota,r)$.
\end{itemize}
For a host $h$, $e$ and $\theta\in\eisX_h(\OO,e)$, $D$ is \emph{of the case of $(h,e,\theta)$} if either $(\ell,e,\theta)$ is in case II, $e_D=e$,
$\psi=\theta_G$, $S$ is the set of prime divisors of $2\ell N_h$, and $\xi=\xi'_G$ for a primitive odd Dirichlet character $\xi'$ of conductor
$2^{n_{II}}$; or $(\ell,e,\theta)$ is in case III, $e_D=e$, $\psi=\theta_G$, $S$ is the set of prime divisors of $2\ell N_hp_1(\chi)$, and
$\xi=\xi_{0,G}\xi'_G$ with $\xi'$ primitive and odd of conductor $2^{n_{III}}$ and $\xi_0$ a character modulo $p_1(\chi)$ of order $\ord\theta$.
Here $n_{II}=\lceil\log_2\max(\ell N_h4^{n_0},2^{c_2(\theta)+6}o_\chi\ell^3N_\chi^2)\rceil$ and
$n_{III}=\lceil\log_2\max(\ell N_\chi p_1(\chi)4^{n_0},8\ell^2)\rceil$, $n_0=n_0(\ell,\theta)$.
\formal{UpperTheory.V.Aut.OfCase}
\end{definition}

\begin{definition}[$(\mathrm{Sp})$]\label{pf:eis:spstmt}
$(\mathrm{Sp})$ is the statement: for every $\Omega$, every host $h$ big for $t$, every prime $\ell\ne2$, every coefficient ring $\OO$, every
$0<e<t$, every $\theta\in\eisX_h(\OO,e)$ and every good raw datum $D$ over $\OO$ of the case of $(h,e,\theta)$, $(\mathrm{S1})$ and
$(\mathrm{S2})$ hold, and $(\mathrm{S3})_\iota$ holds for some $\iota:\Qlbar\cong\C$.
\formal{UpperTheory.V.Aut.SymplecticFourImageDense}
\end{definition}

This is the statement $(\mathrm{Sp})$ of \cref{pf:vwin-he:autstmts}. Its proof rests on five inputs which concern \emph{towers} attached to
a raw datum: finite families of totally real and CM fields over which the residual representation becomes induced or a tensor product
of two-dimensional representations, together with framed deformation problems over these fields; the conditions $\mathrm{Ok}$,
$\mathrm{Two}$, $\mathrm{Deep}$, $\mathrm{Fit}$ and $\mathrm{Four}$ on towers and the contexts II and III of a raw datum are those of the development.

\begin{definition}[Inputs of $(\mathrm{Sp})$]\label{pf:eis:spinputs}
\leavevmode
\begin{enumerate}[label=(I\arabic*)]
\item\label{pf:eis:I1} The local results of the literature package, Poitou--Tate duality, global finiteness, the
  global Euler characteristic formula, orthogonality of unramified classes and the Greenberg--Wiles formula.
\item\label{pf:eis:I2} The theorem of Ferrero--Washington (\cref{lit:window:fw}) \cite{FerreroWashington1979}, in the form: for every abelian
  number field $K$ containing $\mu_\ell$ ($\mu_4$ if $\ell=2$) there is $r$ such that the class groups of the layers $K(\mu_{\ell^n})$ have
  $\ell$-rank at most $r$.
\item\label{pf:eis:I3} Seeds in context II: for every raw datum in context II and every tower $T$ attached to it with $\mathrm{Ok}$, $\mathrm{Two}$
  and $\mathrm{Deep}$, and every $\iota$, after enlarging $\OO$ there is a \emph{seed pair} over the top field of $T$ whose reduction lies in the
  space prescribed by the frame of $T$ and has a non-zero entry of root weight $-(1,0)$, or non-zero entries of root weights $(-1,1)$ and
  $-(1,1)$.
\item\label{pf:eis:I4} Seeds in context III: the analogous statement for context III, with an auxiliary prime $w$ split in the tower, for all
  towers above a given one in which $w$ dies.
\item\label{pf:eis:I5} Potential modularity towers: for every raw datum, every tower $T$ with $\mathrm{Fit}$ and $\mathrm{Deep}$, and every $\iota$, there is a
  coefficient ring $\OO_1$ such that for every larger coefficient ring and every space of points of the frame of $T$ there is a potential
  modularity tower.
\end{enumerate}
\formal{UpperTheory.V.GradeZeroInterpolatingElement.CohomologyHyp}
\formal{UpperTheory.V.Ws.SeedIIAlt}
\formal{UpperTheory.V.Cu.SeedIVAlt}
\formal{UpperTheory.V.Ws.PMExAlt}
\end{definition}

\begin{lemma}\label{pf:eis:hfit}
Every tower over a raw datum in context III which satisfies $\mathrm{Ok}$ and $\mathrm{Four}$ satisfies $\mathrm{Fit}$.
\formallabel{UpperTheory.V.Cu.hFit}
\end{lemma}

\begin{proof}
By definition $\mathrm{Fit}$ is the conjunction of $\mathrm{Ok}$ and of ``context II and $\mathrm{Two}$, or context III and $\mathrm{Three}_0$'', where
$\mathrm{Three}_0$ means $\ell\ge5$, $S_b=\emptyset$ and $\psi$ has order prime to $\ell$; and $\mathrm{Four}$ consists of these three conditions and
a fourth one.
\end{proof}

\begin{theorem}\label{pf:eis:sp}
$\sfname{HypothesesH'}$ implies $(\mathrm{Sp})$. More precisely, \ref{pf:eis:I1}--\ref{pf:eis:I5} imply $(\mathrm{Sp})$.
\formallabel{UpperTheory.V.HV.DescentVanishingReductions.sp4In0_alt}
\formallabel{UpperTheory.V.Cu.sp4In_alt}
\end{theorem}

\begin{proof}
The rank-four reduction theorem of the development (\lref{UpperTheory.V.Cu.main4_alt}) states that \ref{pf:eis:I1},
\ref{pf:eis:I2}, the statement of \cref{pf:eis:hfit}, \ref{pf:eis:I3}, \ref{pf:eis:I4} and \ref{pf:eis:I5} imply $(\mathrm{Sp})$. With
\cref{pf:eis:hfit} this proves the second assertion.

Assume $\sfname{HypothesesH'}$. \ref{pf:eis:I1} is contained in it (\lref{UpperTheory.V.HV.DescentVanishingReductions.adjoint_lift_of_prints},
which extracts from $\sfname{HypothesesH'}$ a package containing \ref{pf:eis:I1}); \ref{pf:eis:I2} is
\lref{UpperTheory.V.FW.ferreroWashington}; \ref{pf:eis:I3} is \lref{UpperTheory.V.Z.Bc.seedIIs_alt}; \ref{pf:eis:I4} is
\lref{UpperTheory.V.Z.Win.Birth.fourth_initial_of_prints_h}; and \ref{pf:eis:I5} is
\lref{UpperTheory.V.HV.DescentVanishingReductions.pmEx_of_printsH}, each of the last three a consequence of
$\sfname{HypothesesH'}$.
\end{proof}

\subsubsection{Proofs of the inputs of \texorpdfstring{$(\mathrm{Sp})$}{(Sp)}}

We now prove the rank-four reduction theorem used in \cref{pf:eis:sp} and the inputs \ref{pf:eis:I3}--\ref{pf:eis:I5}. The
arguments are long in the formal development (roughly $66{,}000$ lines); we give their main steps and cite the intermediate results.
The objects of the development used below are: \emph{towers} $T$ attached to a raw datum (a finite family of totally real and CM
fields, with a top field $F_T$, over which the residual representation becomes induced or a tensor product of two-dimensional
representations, together with a framed deformation problem over $F_T$ and the conditions $\mathrm{Ok}$, $\mathrm{Two}$,
$\mathrm{Four}$, $\mathrm{Fit}$, $\mathrm{Deep}$); \emph{spaces of points} $X$ of the frame of $T$ (an algebra $A_X$ of finite type over a
complete local ring $B_X$, carrying a family of four-dimensional points); \emph{potential modularity towers} $(X,J)$, $J\subset B_X$
an ideal (the \emph{modular locus}) such that $B_X/J$ is Artinian modulo $\mathfrak m_\OO$ and the inertial characters, every point of $X$ on
$V(J)$ with the right Hodge--Tate weights descends to a pure representation of $G_\Q$ which is automorphic when irreducible on
$G_{L_T}$, and seeds and points with distinct inertial characters on $V(J)$ satisfy a local vanishing condition; and \emph{promodular}
points: points of a raw datum which are conjugate to the specialization of a point of a potential modularity tower lying on $V(J)$.

\begin{definition}[Hilbert modular inputs]\label{pf:eis:hilb}
Let $F$ be a totally real field, $\ell$ an odd prime and $\iota:\Qlbar\cong\C$.
\begin{enumerate}[label=(\alph*)]
\item $(\mathrm{Hol})$: for every totally odd Hecke character $\omega$ of $F$ of finite order, of conductor $\mathfrak c\ne1$ prime to $\ell$, and every
  set $P$ of finite places containing a place $w\mid\ell$ with $\omega(\varpi_w)\not\equiv1$ and omitting another such place, and every $B$, there are
  $k\ge B$ with $(\ell-1)\mid k$ and a cuspidal Hilbert eigenform of parallel weight $k+1$, level $\mathfrak c$ with Iwahori level at $\ell$ and
  character $\omega$ whose eigenvalues are congruent, via $\iota$, to those of the Eisenstein series of this weight, level and character which is
  ordinary in the direction prescribed by $P$.
\item $(\mathrm{WC})$: for $[F:\Q]\ge2$, $\ell$ split in $F$ (each infinite place $v$ determining via $\iota$ a place $w_v\mid\ell$), $k_0\ge4$, a level
  $\mathfrak c$ with Iwahori level at all $w_v$ and divisible by a place of norm $2$, a Hecke character $\omega$ of finite order, and a cuspidal
  eigenform $\varphi$ of parallel weight $k_0$, level $\mathfrak c$ and character $\omega$ with eigenvalues $\lambda_w$ such that the $U_{w_v}$-eigenvalues are
  $\ell$-adic units: for every set $S$ of infinite places and every $B$ there are $t\ge B$ with $(\ell-1)\mid t$ and a cuspidal eigenform $\varphi'$ of
  weight $k_0+2t$ at $S$ and $k_0$ elsewhere, of level $\mathfrak c$ and character $\omega$, with eigenvalues $\lambda'_w$ such that
  $\lambda'_w\prod_{v\in S,\,w_v=w}\ell^t\equiv\lambda_w$ modulo the maximal ideal for all $w$.
\item $(\mathrm{Mix})$: for $[F:\Q]$ even, $\ell$ split in $F$, a totally odd character $\xi$ of $G_F$ of finite order, unramified above $\ell$, ramified at a
  place of norm $2$, with $\xi^2\not\equiv1$ and $\xi\not\equiv1$ on each decomposition group above $\ell$, and a set $P$ of places above $\ell$, neither empty nor
  all of them: there are $k,t\ge2$ with $(\ell-1)\mid k,t$, a cuspidal $\pi$ of $\GL_2(\A_F)$ and $r:G_F\to\GL_2(\Qlbar)$ attached to $\pi$, continuous and
  irreducible on open subgroups, with $\det r=\eiscyc^k\xi$, residually $1\oplus\xi$, inertia away from $\ell$ acting through $(1,\xi)$, the condition
  $\mathrm{PR}$, and ordinary above $\ell$ with weights $(k+t,-t)$ at the places of $P$ and $(k,0)$ at the others.
\item $(\mathrm{Ind})$: for a quadratic extension $M/K$ of totally imaginary number fields and a cuspidal $\pi$ of $\GL_2(\A_M)$ of mixed weight
  (the infinitesimal characters at the two embeddings of $M$ above an embedding of $K$ differ) whose ramified places are split over $K$,
  there is a cuspidal $P$ of $\GL_4(\A_K)$ whose Satake parameters at the places where both are unramified are those of the representation
  induced from $\pi$, and whose infinitesimal character at each embedding of $K$ is the sum of those of $\pi$ at the two embeddings above it.
\item $(\mathrm{Ind}')$: the four-dimensional representations $R$ of $G_\Q$ obtained from the data of a tower by inducing (a twist of) a
  two-dimensional representation $r$ of a totally real field of the tower which satisfies $\mathrm{PR}$ are, restricted to the CM field of the
  tower, automorphic with the window condition at all places.
\end{enumerate}
\formal{UpperTheory.V.Sd.PropMHol}
\formal{UpperTheory.Esk.HolWC}
\formal{UpperTheory.V.Z.Win.Birth.MixSeedP}
\formal{UpperTheory.V.Z.InvariantSubmoduleMonicPolynomial.E.IndCs}
\formal{UpperTheory.V.Z.Win.Birth.IndAutLP}
\end{definition}

\begin{theorem}[Rank-four reduction]\label{pf:eis:main4}
\ref{pf:eis:I1}, \ref{pf:eis:I2}, the statement of \cref{pf:eis:hfit}, \ref{pf:eis:I3}, \ref{pf:eis:I4} and \ref{pf:eis:I5} imply $(\mathrm{Sp})$.
\formallabel{UpperTheory.V.Cu.main4_alt}
\end{theorem}

\begin{proof}
Let $D$ be a good raw datum of the case of $(h,e,\theta)$ and fix $\iota$ (\lref{UpperTheory.V.Ws.ambZero}). Consider \emph{pairs} $(r,u)$ where
$r$ is a point of $D$ with values in the valuation ring of a field of characteristic $0$ or $\ell$ (multiplier $\vartheta$, reduction $\bar r$, the local
conditions at $S_b$ and $\ell$, the latter after conjugation by $u\in\GSp_4$), and say that $r$ is \emph{starred} if it is irreducible or a non-split
block extension.

\emph{Step 1: starred points are promodular.} In context II this is \lref{UpperTheory.V.Ws.proMod_of_pmEx_of_ctx2_of_kPair_of_star},
from \ref{pf:eis:I1}, \ref{pf:eis:I2}, \ref{pf:eis:I3} and \ref{pf:eis:I5} (the seed pairs of \ref{pf:eis:I3} are points on the modular loci of
the potential modularity towers of \ref{pf:eis:I5}, and the starred point is connected to them as in context III below). In context III it is
\lref{UpperTheory.V.Cu.promodularity_case_three_of_chain_reduction}, combining four results:
(a) for every admissible labelling there is a tower with $\mathrm{Ok}$ and $\mathrm{Four}$, tamely ramified above $\ell$, over which the residual
representation becomes a tensor product (\lref{UpperTheory.V.Cu.existsTensorSeedTower_holds});
(b) the chain lemma (\lref{UpperTheory.V.Cu.chain_rank_four_of_red_a_y_coin_fib_y}, from three lemmas on reduction types of residual
points, \lref{UpperTheory.V.Cu.redTanA}, \lref{UpperTheory.V.Cu.yCoin}, and \lref{UpperTheory.V.Cu.torusY}): in a potential modularity
tower over a tower with $\mathrm{Fit}$ and $\mathrm{Deep}$ and of sufficiently large degree, at a good residual point $x$, if one irreducible
component of $\Spec A_X$ through $x$ lies in $V(J)$, then all do;
(c) chains of good residual points connecting two residual points descend along enlargements of the coefficient ring
(\lref{UpperTheory.V.Cu.chainDesc4});
(d) from \ref{pf:eis:I4} and the statement of \cref{pf:eis:hfit}: for a tower $T_0$ as in (a) there is a larger tower $T_1$ such that over every deeper
tower $T$ there is a seed pair whose residual point is connected to the residual point of $D$ by a chain (\lref{UpperTheory.V.Cu.chainIV_alt}).
Given a starred point, \ref{pf:eis:I5} provides a potential modularity tower over a deep tower as in (d); the seed pair lies on $V(J)$, and by (b)
applied along the chain of (c), (d), modularity propagates from the component of the seed to all components through the residual point of
$D$, in particular to the component carrying the given point.

\emph{Step 2: consequences of promodularity.} (i) There is no starred pair in characteristic $\ell$
(\lref{UpperTheory.V.Ws.contradiction_star_promodular}: in characteristic $\ell$ a promodular point has the traces of $\bar r$, a sum of four
characters, so its traces satisfy $\tr r(xgh)=\tr r(xhg)$ and it can be neither irreducible nor a non-split block extension); (ii) promodular points of the right weights are pure
(\lref{UpperTheory.V.Ws.pureW_of_proModular_of_isZero}); (iii) irreducible promodular points of characteristic $0$ satisfy
$\mathrm{Aut}_4(\iota,r)$ (\lref{UpperTheory.V.Ws.aut_n_alt}), by the automorphy axiom of potential modularity towers and descent from
the CM field of the tower to the imaginary quadratic fields $K$.

\emph{Step 3: conclusion.} From (i)--(iii) and \ref{pf:eis:I1}, \lref{UpperTheory.V.Ws.fin0_no_end_all_irr_of_good_con_purity_aut} derives
$(\mathrm{S1})$, $(\mathrm{S2})$ and $(\mathrm{S3})_\iota$: a positive-dimensional family of points, or a block-diagonal point with spanning upper block,
would produce a starred pair in characteristic $\ell$, contradicting (i), and every point whose image spans $M_4(\Qlbar)$ is irreducible and
promodular, hence automorphic by (iii) (\lref{UpperTheory.V.Cu.modularity_rank_four_of_promodularity_case_three}).
\end{proof}

\begin{theorem}\label{pf:eis:seedII}
$\sfname{HypothesesH'}$ implies \ref{pf:eis:I3}.
\formallabel{UpperTheory.V.Z.Bc.seedIIs_alt}
\end{theorem}

\begin{proof}
\emph{Step 1: $(\mathrm{Hol})$.} By \lref{UpperTheory.V.Z.Bc.propMHolZ_alt}, $\sfname{HypothesesH'}$ together with the published existence
of Hilbert Eisenstein series of weight one (part of the literature package) implies $(\mathrm{Hol})$: this is Ribet's construction over $F$, as in
\cref{pf:eis:propEC}, with Hilbert Eisenstein series and a form congruent to $1$.

\emph{Step 2: $(\mathrm{Mix})$.} By \lref{UpperTheory.Esk.mixSeed_of_holWC_alt}, $(\mathrm{Hol})$, $(\mathrm{WC})$ (\cref{pf:eis:wc}) and
$\sfname{HypothesesH'}$ imply $(\mathrm{Mix})$ (via \lref{UpperTheory.V.Z.Win.Birth.mixOfP}): starting from a form of parallel weight
$k+1$ given by $(\mathrm{Hol})$ (with $\omega$ the Hecke character of $\xi$), $(\mathrm{WC})$ applied to the set $S$ of infinite places corresponding to
the places in $P$ produces a congruent eigenform of weight $k+1+2t$ at $S$; its Galois representation, attached by the results of the
literature package and twisted, is ordinary of weights $(k+t,-t)$ on $P$ and $(k,0)$ off $P$, residually $1\oplus\xi$ by the congruence, and
irreducible on open subgroups because the weights are not of CM type (\lref{UpperTheory.V.Z.Win.Birth.mix_birth_of_prints}).

\emph{Step 3: $(\mathrm{Ind}')$.} By \lref{UpperTheory.V.Z.InvariantSubmoduleMonicPolynomial.E.indAutLP_of_printsH_upQuad_upSol_indCs},
$(\mathrm{Ind}')$ follows from $\sfname{HypothesesH'}$, quadratic base change from totally real to CM fields
(\lref{UpperTheory.V.Z.Bc.upQuadI_of_alt}), cyclic base change in rank two (\lref{UpperTheory.V.Z.Bc.upSolI_of_alt}) and
$(\mathrm{Ind})$ (\cref{pf:eis:ind}): the representation $R$ restricted to the CM field of the tower is induced from a quadratic extension of CM
fields, so it is automorphic once the two-dimensional representation is, after base change from the totally real field.

\emph{Step 4: seed pairs.} By \lref{UpperTheory.V.Z.Bc.seedII_of_p_alt}, $(\mathrm{Mix})$ and $(\mathrm{Ind}')$ imply \ref{pf:eis:I3}: given a raw datum
in context II and a tower $T$ with $\mathrm{Ok}$, $\mathrm{Two}$, $\mathrm{Deep}$, $(\mathrm{Mix})$ over the totally real field of the tower, with $P$ the places
above $\ell$ singled out by the tower, gives $r$; the four-dimensional representation $R$ induced from a twist $r\otimes\delta$ is irreducible
(\lref{UpperTheory.induced_irreducible}), has the local properties required by the frame of $T$
(\lref{UpperTheory.GalN.V.indLoc0}), and is automorphic with the window condition by $(\mathrm{Ind}')$; its reduction lies in the space
prescribed by the frame and has the required non-zero root entries (\lref{UpperTheory.induced_final_statement_alt}), so it is a seed
pair.
\end{proof}

\begin{theorem}\label{pf:eis:seedIV}
$\sfname{HypothesesH'}$ implies \ref{pf:eis:I4}.
\formallabel{UpperTheory.V.Z.Win.Birth.fourth_initial_of_prints_h}
\end{theorem}

\begin{proof}
By \lref{UpperTheory.V.Cu.fourth_initial_of_rank_four_data}, \ref{pf:eis:I4} follows from the following rank-four data, all available
under $\sfname{HypothesesH'}$.
(a) Automorphy of tensor products: a tensor product of two two-dimensional representations over the CM field of a tower, one automorphic
over a totally real field and one induced from a character, is automorphic with the window condition
(\lref{UpperTheory.V.Cu.rank_four_automorphy_of_tensor_of_base_change_of_induction}, from $\sfname{HypothesesH'}$, quadratic and cyclic
base change, automorphic induction $(\mathrm{Ind})$ in the form \lref{UpperTheory.V.Z.InvariantSubmoduleMonicPolynomial.E.indAutCs_of_alt},
and a polarization statement \lref{UpperTheory.V.Cu.polOf4_alt}).
(b) Eisenstein-congruent two-dimensional representations, as in \cref{pf:eis:propE}
(\lref{UpperTheory.V.Z.Win.Birth.propEGP_of_printsH}).
(c) Auxiliary characters of weight four and twists adapted to the auxiliary prime $w$
(\lref{UpperTheory.V.Cu2.exists_weight_four_auxiliary_characters}, \lref{UpperTheory.V.Cu.tauEx'}, and \lref{UpperTheory.V.Cu.twEx4}).
(d) The verification that the resulting tensor product is a seed for the frame of the tower: it preserves the symplectic pairing up to the
multiplier, has the prescribed residual shape, entries and weights, is compatible with the towers in which $w$ dies, and is irreducible
(\lref{UpperTheory.V.Cu.irrOf4_of}) and automorphic (\lref{UpperTheory.V.Cu.autOf4_of}).
\end{proof}

\begin{theorem}\label{pf:eis:pmex}
$\sfname{HypothesesH'}$ implies \ref{pf:eis:I5}.
\formallabel{UpperTheory.V.HV.DescentVanishingReductions.pmEx_of_printsH}
\end{theorem}

\begin{proof}
$\sfname{HypothesesH'}$ contains the unramified local inputs (\lref{UpperTheory.V.HV.DescentVanishingReductions.adjoint_lift_of_prints})
and implies the adjoint lifting theorem (\lref{UpperTheory.V.HV.DescentVanishingReductions.adjoint_lifting_of_cohomology}). From these,
\lref{UpperTheory.Gal.projectedFormsSystemPart.pm_exists_of_unramified_adjoint_lifting} proves \ref{pf:eis:I5} in two steps.
(a) Local vanishing at modular points (\lref{UpperTheory.V.Th.local_zero_of_modular_p}): at a point of a space of points attached to an
automorphic representation, with distinct inertial characters at $\ell$, the local vanishing condition of potential modularity towers holds;
this follows from the adjoint lifting theorem by lifting eigenbases of the inertial action (\lref{UpperTheory.V.Th.eigenbasis_rank_four}).
(b) Construction of the towers (\lref{UpperTheory.Gal.projectedFormsSystemPart.promodular_existence_of_local_zero}): for a tower with
$\mathrm{Fit}$ and $\mathrm{Deep}$, after enlarging the coefficient ring, the ideal $J$ cutting out the components of a space of points that
contain modular points satisfies the axioms: the modular quotient is regular at its points
(\lref{UpperTheory.V.Ws.modular_deformation_quotient_regular_at_point_alt}), points on $V(J)$ descend to $G_\Q$
(\lref{UpperTheory.V.Z.Win.Dn.pmDown_alt}) and are pure (\lref{UpperTheory.V.Ws.pmPur_w_alt}), and the required links between
residual points exist (\lref{UpperTheory.Gal.projectedFormsSystemPart.exists_two_links_same_kappa_not_dvd_of_fit_deep}).
\end{proof}

The last two inputs are formal proofs of special cases of classical theorems. The weight-changing congruence $(\mathrm{WC})$ is a
consequence of Hida's theory of nearly ordinary Hecke algebras.

\begin{littheorem}[Hida {\cite{Hida1988,Hida1989}}]\label{pf:eis:hida}
Let $F$ be a totally real field, $p$ an odd prime, $\mathfrak n$ an ideal prime to $p$, and $\mathbf h^{\mathrm{n.o.}}$ the nearly ordinary Hecke algebra
of level $\mathfrak np^\infty$ with coefficients in the ring of integers $\OO$ of a finite extension of $\Q_p$. Then $\mathbf h^{\mathrm{n.o.}}$ is finite and
torsion-free over the Iwasawa algebra $\OO[[\mathbf T]]$ of the torsion-free part $\mathbf T$ of $(\OO_F\otimes\Z_p)^\times\times(\OO_F\otimes\Z_p)^\times$ modulo the
closure of the global units, and for every arithmetic weight $(k,w)$ (with $k_v\ge2$ for all infinite places $v$) the specialization of
$\mathbf h^{\mathrm{n.o.}}$ at the corresponding arithmetic point is the nearly ordinary Hecke algebra of weight $(k,w)$, level $\mathfrak np$ and the
corresponding character. In particular, the system of Hecke eigenvalues of a nearly ordinary eigenform of weight $(k,w)$ is congruent
modulo the maximal ideal (after the standard normalization of the $U_p$-eigenvalues) to that of a nearly ordinary eigenform of every weight
$(k',w')$ congruent to $(k,w)$ modulo $p-1$ componentwise.
\end{littheorem}

\begin{theorem}\label{pf:eis:wc}
$(\mathrm{WC})$ holds.
\formallabel{UpperTheory.Esk.holWC}
\end{theorem}

The statement $(\mathrm{WC})$ is the special case of \cref{pf:eis:hida} for $p=\ell$ split in $F$, starting from parallel weight $k_0\ge4$ and changing the
weight by $2t$ at a set $S$ of infinite places, with $(\ell-1)\mid t$; the factors $\ell^t$ in the congruence are the change of normalization of the
$U_{w_v}$-eigenvalues. The formal statement differs from \cref{pf:eis:hida} in that it only asserts the existence of congruent forms for
arbitrarily large $t$ (not for all congruent weights), requires $[F:\Q]\ge2$, Iwahori level at all places above $\ell$ and a place of norm $2$ in the level
(which makes the relevant arithmetic groups torsion-free), and makes no statement about Hecke algebras over $\OO[[\mathbf T]]$; and it is proved
directly by a comparison of traces, without $p$-adic families.

\begin{proof}
Let the data be given; let $\mathcal M_k$ denote the space of automorphic forms of weight $k$ (a function on the infinite places), level $\mathfrak c$
and central character $\omega$ to which the trace formula below applies, with the Hecke operators $T_w$ ($w\nmid\mathfrak c$) and $U_w$ ($w\mid\mathfrak c$);
$\varphi\in\mathcal M_{k_0}$. These spaces are finite-dimensional and the Hecke operators preserve them and commute
(\lref{UpperTheory.Esk.s_HFin}). Put $N=\dim\mathcal M_{k_0}$ and
$t=(\ell-1)\ell^N\max(B,1)$, so $t\ge B$, $(\ell-1)\mid t$, and $(\ell-1)\ell^N$ divides the weight shift at every infinite place.

\emph{Trace congruences.} For every Hecke operator $h$ which is a product of double cosets at a list $W$ of places containing every place
above $\ell$ at least $r\ge1$ times, the traces of $h$ on $\mathcal M_{k_0}$ and on $\mathcal M_{k_0+2t\mathbf 1_S}$, each multiplied by the normalizing factor
$\prod_{w\in W}\Nm w^{k/2-1}\prod_{v:\,w_v=w}\ell^{t_v}$, are $\ell$-integral and congruent modulo $\ell^{\min(N,r)}$ (\lref{UpperTheory.Esk.pf_TrCong}).
This is derived from the trace formula for the test functions defined by $h$ and the archimedean projectors onto the discrete series of
the two weights (\lref{UpperTheory.Esk.s_TF}): in outline, the two traces have the same geometric expansion except for the
archimedean factors, which depend on the weight through characters of the archimedean tori evaluated at the conjugacy classes; after the
normalization these factors are $\ell$-integral and congruent modulo $\ell^{N+1}$ when the weights differ by multiples of $(\ell-1)\ell^N$, and the
terms in which the $U$-operators repeated $r$ times do not act through units are divisible by $\ell^r$. The analytic inputs are bounds for
the number of contributing conjugacy classes and the convergence and integrability of the hyperbolic terms
(\lref{UpperTheory.Esk.dAx}).

\emph{From traces to eigenforms} (\lref{UpperTheory.Esk.pf_HolWC}). Taking $r\ge N$, the normalized traces of all Hecke operators
containing the $U$-operators above $\ell$ at least $r$ times agree on the two spaces modulo $\ell^N$; these operators see exactly the ordinary
parts. On $\mathcal M_{k_0}$ the eigensystem $\lambda$ of $\varphi$ is ordinary and occurs with a multiplicity $1\le d\le N$, and $v_\ell(d)<N$; comparing the
traces of suitable polynomials in the Hecke operators (which isolate the eigensystems congruent to $\lambda$ modulo the maximal ideal) shows
that the normalized eigensystem $\lambda$ also occurs modulo the maximal ideal in $\mathcal M_{k_0+2t\mathbf 1_S}$. Hence some common eigenvector
$\varphi'\in\mathcal M_{k_0+2t\mathbf 1_S}$ of all Hecke operators has eigenvalues $\lambda'$ with the asserted congruences, and it is a cusp form of the required
weight, level and character.
\end{proof}

The automorphic induction statement $(\mathrm{Ind})$ is a special case of the following theorem.

\begin{littheorem}[Arthur--Clozel {\cite[Ch.~3, Thm.~6.2]{ArthurClozel1989}}]\label{pf:eis:ac}
Let $E/F$ be a cyclic extension of number fields of prime degree $q$ and $\pi$ a cuspidal automorphic representation of $\GL_m(\A_E)$. There is an
automorphic representation $\Pi$ of $\GL_{mq}(\A_F)$, induced from cuspidal, such that at every place $v$ of $F$ at which $E$ and $\pi$ are
unramified, $\Pi_v$ is the local automorphic induction of $(\pi_w)_{w\mid v}$: its Satake parameters are the union of those of the $\pi_w$ if $v$ splits,
and the $q$-th roots of those of $\pi_w$ if $v$ is inert. $\Pi$ is cuspidal if and only if $\pi$ is not isomorphic to its conjugates under $\Gal(E/F)$.
\end{littheorem}

\begin{theorem}\label{pf:eis:ind}
$(\mathrm{Ind})$ holds.
\formallabel{UpperTheory.V.Z.Bc.indCs}
\end{theorem}

The statement $(\mathrm{Ind})$ is \cref{pf:eis:ac} for $q=2$, $m=2$ and totally imaginary $M/K$. The mixed weight condition makes $\pi$ not isomorphic
to its conjugate, so $P$ is cuspidal; the formal statement adds the description of the infinitesimal characters (the archimedean case of
local induction) and asserts the Satake parameters only at the places where both representations are unramified, under the assumption
that the ramified places of $\pi$ are split over $K$.

\begin{proof}
The formal proof compares spectral sides of trace formulas for $\GL_2$ over $M$ and $\GL_4$ over $K$ for matching test functions, in outline as
follows. (a) The matching of the geometric sides for unit elements of the spherical Hecke algebras, i.e.\ the base change fundamental lemma
for the unit element (\lref{PositivityTheory.Rg.Tf.Loc.bcflUnit}, a formal proof of Kottwitz's theorem \cite{Kottwitz1986}), gives a
global matching statement for the test functions used below (\lref{UpperTheory.ConvolutionTheory.GM.gm_of_bcfl}). (b) The contribution of
$\pi$ is isolated in the spectral side by Hecke operators at finitely many split places: Iwahori--Hecke injectivity and spectral norm
estimates for Chebyshev-type polynomials in the Hecke operators make the trace of the resulting operator on the cuspidal spectrum
non-zero on the $\pi$-isotypic part and small elsewhere (\lref{UpperTheory.V.Z.InvariantSubmoduleMonicPolynomial.cix}). (c) The
matched trace on the $\GL_4$-side is therefore non-zero, which produces a cuspidal $P$ whose unramified Satake parameters and
infinitesimal characters are those of the induced representation
(\lref{UpperTheory.V.Z.InvariantSubmoduleMonicPolynomial.E.coreInd}); finally the normalizations of the Satake parameters and
infinitesimal characters are adjusted (\lref{UpperTheory.V.Z.InvariantSubmoduleMonicPolynomial.E.indCs_of_gm}).
\end{proof}
